\documentclass[11pt]{amsart}
\usepackage[margin=1.1in]{geometry}
\usepackage{amsmath,amssymb,amsthm}
\usepackage{booktabs}
\usepackage{array}
\usepackage{microtype}
\usepackage[colorlinks=true,linkcolor=blue,citecolor=blue,urlcolor=blue]{hyperref}

\newtheorem{theorem}{Theorem}[section]
\newtheorem{lemma}[theorem]{Lemma}
\newtheorem{proposition}[theorem]{Proposition}
\newtheorem{corollary}[theorem]{Corollary}
\newtheorem{introthm}{Theorem}
\renewcommand{\theintrothm}{\Alph{introthm}}
\theoremstyle{definition}
\newtheorem{definition}[theorem]{Definition}
\theoremstyle{remark}
\newtheorem{remark}[theorem]{Remark}
\numberwithin{equation}{section}

\newcommand{\Z}{\mathbb{Z}}
\newcommand{\R}{\mathbb{R}}
\newcommand{\C}{\mathbb{C}}
\newcommand{\T}{\mathbb{T}}
\newcommand{\Disc}{\mathbb{D}}
\newcommand{\E}{\mathbb{E}}
\renewcommand{\Pr}{\mathbb{P}}
\newcommand{\tr}{\operatorname{tr}}
\newcommand{\Tr}{\operatorname{Tr}}
\newcommand{\spec}{\operatorname{spec}}
\renewcommand{\Re}{\operatorname{Re}}
\renewcommand{\Im}{\operatorname{Im}}
\newcommand{\Li}{\operatorname{Li}_2}
\newcommand{\Cl}{\operatorname{Cl}_2}
\newcommand{\FD}{F_{\mathrm{DKZ}}}
\newcommand{\IME}{I_{\mathrm{ME}}}
\newcommand{\one}{\mathbf{1}}
\newcommand{\ip}[2]{\langle #1,#2\rangle}
\newcommand{\cD}{\mathcal{D}}
\newcommand{\cG}{\mathcal{G}}
\newcommand{\cH}{\mathcal{H}}
\newcommand{\cF}{\mathcal{F}}
\newcommand{\cN}{\mathcal{N}}
\newcommand{\Sh}{\Xi^{\sharp}}
\newcommand{\Xb}{\overline{\Xi}}
\newcommand{\Phic}{\Phi_{\mathrm c}}
\newcommand{\lmin}{\lambda_{\min}}
\newcommand{\lmax}{\lambda_{\max}}
\newcommand{\pv}{\mathrm{p.v.}}
\newcommand{\dd}{\,\mathrm{d}}
\newcommand{\Q}{\mathbb{Q}}
\newcommand{\cL}{\mathcal{L}}
\newcommand{\vc}{v_{\mathrm c}}
\newcommand{\veven}{v_{\mathrm{even}}}
\newcommand{\vodd}{v_{\mathrm{odd}}}
\newcommand{\CH}{\mathrm{CH}}
\newcommand{\DKZ}{\mathrm{DKZ}}
\setlength{\emergencystretch}{3em}

\begin{document}

\title[Optimal CGLMP measurements in every dimension]{A strip inequality for projections, two-variable BMV positivity,
and optimal CGLMP measurements in every dimension}

\author{Ansh Mishra}
\author{Aryan Senthilkumar}

\date{October 2026}

\subjclass[2020]{Primary 15A42, 81P40; Secondary 47A63, 42A50, 44A10, 65G40}

\keywords{CGLMP inequality, Bell inequalities, maximally entangled states, rigidity, BMV conjecture, Stahl's theorem,
conjugate function, Stein--Weiss theorem, matrix pencils, noise resistance, statistical strength, Kullback--Leibler
divergence, computer-assisted proof, formal verification, Lean}

\begin{abstract}
We prove that for every number of outcomes $d\ge2$ the measurements of Durt, Kaszlikowski and \.Zukowski (DKZ)
maximise the CGLMP Bell expression $I_d$ over all projective measurements on maximally entangled states of any local
dimension $D$, the maximal value being $\IME(d)=\frac4{d(d-1)}\sum_{j=1}^{d-1}(d-j)\sec\frac{\pi j}{2d}$, and that
they are the only maximisers: equality forces $d\mid D$ and, up to a local unitary $u\otimes\bar u$, the strategy
is the DKZ strategy tensored with the identity. This settles, for projective measurements, the maximally entangled
clause of Part~B of Open Quantum Problem~27 of IQOQI Vienna.

The proof has an analytic part and a computer-assisted part. The analytic part is an inequality for a projection
$B$ and a Hermitian matrix $g$ of arbitrary size: the sum of an explicit harmonic function $h_\lambda$ of the strip
$0<\Re z<1$ over the eigenvalues of $B+ig$ dominates $\Tr[(P(g-\lambda)P)_+]$, $P=\one-B$, with an exact formula for
the defect. It follows from a two-variable positivity theorem of Bessis--Moussa--Villani (Stahl) type for
projections with an explicit density $\frac1{2\pi}\sum_i|\Im\xi_i(s,\tau)|$, where the $\xi_i$ are the roots of
$\det(g-s-\xi(P-\tau))$; the proof uses the Fourier-slice mechanism of Hein\"avaara's tracial joint spectral
measures. Through a noncommutative continuum version of the Stein--Weiss rearrangement theorem and a family of cell
embeddings, the strip inequality reduces the CGLMP problem to a statement about Clausen kernels: an explicit vector
lies in a convex cone generated by explicit cell vectors. We verify this cone condition for every $d$ by interval
arithmetic, with linear-programming certificates on rational cells for $d\le200$, certified Gaussian-modulated
Dirichlet constructions for $201\le d\le2000$, and a uniform interval Taylor-model argument for $d\ge2001$.

We also give a second, distribution-free proof of the two-variable positivity theorem, through a Radon identity for
matrix pencils whose proof rests on a complex inviscid Burgers equation for the pencil roots. The argument has been
formalised in Lean~4 with Mathlib: for $2\le d\le20$ the formal proof of optimality and rigidity is complete,
including the cone condition; for $d\ge21$ it is complete except for the cone condition, which is established by the
interval-arithmetic certificates; and the strip inequality and the positivity theorem are formalised in full.

We also settle the remaining clauses of Part~B as originally posed, so that, with the negative answer to Part~A by
Bancal, Gisin and Pironio, every clause of Problem~27 has a verdict. The noise-resistance clause holds for every $d$
for the CGLMP inequality, with the DKZ measurements as the unique optimum; in Gill's literal reading (uniformly random
outcomes as noise, violation of any Bell inequality) it fails for every $d\ge4$ on the maximally entangled state of two
$d$-level systems, by a projective strategy with unused outcomes whose critical visibility we compute exactly. The
Kullback--Leibler clause fails for every $d\ge4$. A strengthening of the noise clause, with white noise on the state and
every outcome used, remains open.
\end{abstract}

\maketitle

\setcounter{tocdepth}{1}
\tableofcontents

%======================================================================
\section{Introduction}\label{sec:intro}
%======================================================================

\subsection{The problem}
The CGLMP inequalities, introduced by Collins, Gisin, Linden, Massar and Popescu~\cite{CGLMP}, are Bell inequalities for
two parties with two measurement settings each and $d$ outcomes per measurement. Write $A_1,A_2$ and $B_1,B_2$ for the
outcomes (elements of $\Z_d$) of Alice's and Bob's measurements and $P(A_x=B_y+k)$ for the probability that
$A_x-B_y\equiv k\pmod d$ when the settings $x$ and $y$ are used. The CGLMP expression is
\begin{multline}\label{eq:cglmp}
I_d=\sum_{k=0}^{\lfloor d/2\rfloor-1}\Bigl(1-\frac{2k}{d-1}\Bigr)\Bigl(\bigl[P(A_1=B_1+k)+P(B_1=A_2+k+1)\\
+P(A_2=B_2+k)+P(B_2=A_1+k)\bigr]-\bigl[P(A_1=B_1-k-1)+P(B_1=A_2-k)\\
+P(A_2=B_2-k-1)+P(B_2=A_1-k-1)\bigr]\Bigr)
\end{multline}
(this is the reading of~\cite{CGLMP} under which the local bound is derived and the optimal measurements of~\cite{CGLMP}
give the value~\eqref{eq:IME} below; see~\cite[Section~2]{MS-classical}). Local hidden variable models satisfy
$I_d\le2$~\cite{CGLMP}. For $d=2$, $I_2$ is the CHSH expression; for $d=3$ see also the Clauser--Horne form
of~\cite{KKCZO}.

A \emph{projective strategy on a maximally entangled state} consists of a local dimension $D\ge1$, the state
$\Phi_D=D^{-1/2}\sum_{i=1}^D|ii\rangle\in\C^D\otimes\C^D$, and projection-valued measures $\{P^x_a\}_{a\in\Z_d}$
(Alice, $x=1,2$) and $\{Q^y_b\}_{b\in\Z_d}$ (Bob, $y=1,2$) on $\C^D$; it gives
$P(a,b|x,y)=\langle\Phi_D|P^x_a\otimes Q^y_b|\Phi_D\rangle$. The measurements found numerically by Kaszlikowski et
al.~\cite{KGZMZ} and described analytically by Durt, Kaszlikowski and \.Zukowski~\cite{DKZ} (diagonal phase shifts
followed by a discrete Fourier transform, then a measurement in the computational basis; we call this the \emph{DKZ
strategy}, $D=d$; the explicit bases are recalled in Section~\ref{subsec:dkz}) give
\begin{equation}\label{eq:IME}
\IME(d)=\frac{4}{d(d-1)}\sum_{j=1}^{d-1}(d-j)\sec\frac{\pi j}{2d},
\end{equation}
for instance $\IME(3)=2.8729\ldots$, $\IME(4)=2.8962\ldots$, $\IME(20)=2.9547\ldots$, and
$\IME(d)\to32G/\pi^2=2.9698\ldots$ as $d\to\infty$, where $G$ is Catalan's constant (Corollary~\ref{cor:limit}).
For $d\ge3$ the maximally entangled state is not the optimal state for $I_d$~\cite{ADGL}.

Open Quantum Problem~27 of IQOQI Vienna, ``The power of CGLMP inequalities''~\cite{OQP}, goes back to questions of
Gill~\cite{Gill}. Part~A asks whether every facet of the local polytope of the scenario with two settings and $d$
outcomes per party, other than the positivity facets, is of CGLMP type; this was answered negatively by Bancal,
Gisin and Pironio~\cite{BGP}. Part~B records that the observables which numerically maximise the CGLMP violation on
a maximally entangled state are measurements in the computational basis transformed by a discrete Fourier transform
and diagonal unitaries, and asks to show that this is necessarily the case; it also asks to show that these
measurements give the highest noise resistance and the best Kullback--Leibler discrimination from local
models~\cite{vDGG}. Theorems~\ref{thm:A} and~\ref{thm:B} below settle the first clause of Part~B, and
Section~\ref{sec:clauses} settles the other two.

\subsection{Main results}
Our first two results answer the first clause of Part~B for projective measurements.

\begin{introthm}[optimality; computer-assisted]\label{thm:A}
Let $d\ge2$ and $D\ge1$. For every projective strategy on the maximally entangled state $\Phi_D$,
\[
I_d\le\IME(d),
\]
and the DKZ strategy attains $\IME(d)$.
\end{introthm}

\begin{introthm}[rigidity; computer-assisted]\label{thm:B}
Let $d\ge2$, $D\ge1$, and suppose that a projective strategy on $\Phi_D$ attains $I_d=\IME(d)$. Then $d$ divides $D$,
and there is a unitary $u\colon\C^D\to\C^d\otimes\C^{D/d}$ such that
$uP^x_au^*=P^{x,\mathrm{DKZ}}_a\otimes\one$ and $\bar uQ^y_b\bar u^{\,*}=Q^{y,\mathrm{DKZ}}_b\otimes\one$ for all
$x,y,a,b$; moreover $(u\otimes\bar u)\Phi_D=\Phi_d\otimes\Phi_{D/d}$, and $u$ is unique up to $u\mapsto(\one_d\otimes u')u$
with $u'$ unitary. Consequently, if $d\nmid D$ the maximum of $I_d$ over projective strategies on $\Phi_D$ is
strictly smaller than $\IME(d)$.
\end{introthm}

The proofs of Theorems~\ref{thm:A} and~\ref{thm:B} are computer-assisted only through Theorem~\ref{thm:E} below; all
other steps are proved by hand. The argument has also been formalised in the Lean~4 theorem prover~\cite{Lean,Mathlib}
(Section~\ref{subsec:formal}): for $2\le d\le20$ the formal proof of Theorems~\ref{thm:A} and~\ref{thm:B} is complete,
including Theorem~\ref{thm:E}, and for $d\ge21$ it is complete except for the input Theorem~\ref{thm:E}. Thus
Theorems~\ref{thm:A} and~\ref{thm:B} are proved for every $d$, with a computer-assisted step (interval arithmetic) for
$d\ge21$, and are fully formalised for $d\le20$. The analytic heart of the argument is a matrix inequality, which we
now state. Let
$S=\{z\in\C:0<\Re z<1\}$ and, for $\lambda\in\R$ and $z=x+iy\in S$,
\begin{equation}\label{eq:hlambda}
h_\lambda(x+iy)=-\frac1{\pi^2}\Im\Li\bigl(e^{\pi(y-\lambda)}e^{-i\pi x}\bigr),
\end{equation}
where $\Li$ is the dilogarithm (principal branch, holomorphic on $\C\setminus[1,\infty)$). The function
$h_\lambda$ is harmonic on $S$ and extends continuously to the closed strip, with boundary values $(y-\lambda)_+$
on the line $\Re z=0$ and $0$ on the line $\Re z=1$ (Lemma~\ref{lem:h}). For $0<x<1$ let
\begin{equation}\label{eq:poisson}
K_x(u)=\frac{\sin\pi x}{2(\cosh\pi u-\cos\pi x)}\qquad(u\in\R),
\end{equation}
the Poisson kernel of the strip $S$ for its left edge. For an orthogonal projection $E$ and a Hermitian matrix $X$
we write $\Tr_E f(EXE)$ for the trace over $\operatorname{ran}E$ of $f$ applied to the compression of $X$ to
$\operatorname{ran}E$, and $X_+$ for the positive part of a Hermitian matrix $X$.

\begin{introthm}[strip inequality]\label{thm:C}
Let $M\ge1$, let $B$ be an orthogonal projection on $\C^M$, $P=\one-B$, and let $g$ be a Hermitian $M\times M$
matrix. For every $\lambda\in\R$,
\begin{equation}\label{eq:star}
\sum_{\nu\in\spec(B+ig)}h_\lambda(\nu)\;-\;\Tr\bigl[(P(g-\lambda)P)_+\bigr]
=\int_0^1\!\!\int_\R K_{1-\tau}(s-\lambda)\,\rho(s,\tau)\dd s\dd\tau\;\ge\;0,
\end{equation}
where the eigenvalues are counted with algebraic multiplicity and $\rho\ge0$ is the density of
Theorem~\ref{thm:D}. Equality holds for one $\lambda$ if and only if it holds for all $\lambda$, if and only if
$Bg=gB$.
\end{introthm}

\begin{introthm}[two-variable BMV positivity for projections]\label{thm:D}
In the situation of Theorem~\ref{thm:C} let $r=\operatorname{rank}B$, let $\lmin,\lmax$ be the extreme eigenvalues of
$g$, and for $s\in\R$ and $0<\tau<1$ let $\xi_1(s,\tau),\dots,\xi_M(s,\tau)$ be the roots, with multiplicity, of the
polynomial $\xi\mapsto\det(g-s-\xi(P-\tau))$ of degree $M$. Put
\[
\rho(s,\tau)=\frac1{2\pi}\sum_{i=1}^M\bigl|\Im\xi_i(s,\tau)\bigr|\quad(0<\tau<1),\qquad \rho(s,\tau)=0\quad(\tau\notin(0,1)).
\]
Then $\rho$ is continuous on $\R\times(0,1)$, vanishes unless $\lmin<s<\lmax$, satisfies
$\rho(s,\tau)\le\frac{M}{2\pi}\|g-s\|\,(\tau(1-\tau))^{-1/2}$, hence is integrable, and for all $(a,t)\in\C^2$
\begin{equation}\label{eq:2bmv}
\cD(a,t):=\Tr e^{ag-tP}-e^{-t}\Tr_Pe^{aPgP}-\Tr_Be^{aBgB}=a^2\int_0^1\!\!\int_\R e^{as-t\tau}\rho(s,\tau)\dd s\dd\tau .
\end{equation}
Moreover $\int_\R\rho(s,\tau)\dd s=\|BgP\|_F^2$ for every $\tau\in(0,1)$.
\end{introthm}

Taking $a=1$ in~\eqref{eq:2bmv} shows that $t\mapsto\Tr e^{g-tP}$ is the Laplace transform of the positive measure
$\Tr_B(e^{BgB})\,\delta_0+\Tr_P(e^{PgP})\,\delta_1+\bigl(\int_\R e^s\rho(s,\tau)\dd s\bigr)\dd\tau$ on $[0,1]$, which
is Stahl's theorem~\cite{Stahl} (the former Bessis--Moussa--Villani conjecture~\cite{BMV}) for a projection, with an
explicit density. Theorem~\ref{thm:D} is stronger: it asserts positivity of a measure in two variables, after the
two pinched terms are subtracted and the result is divided by $a^2$. Neither modification can be omitted:
$\cD(0,t)=0$ for all $t$, so $\cD$ itself is the Laplace transform of a positive measure only if $\cD\equiv0$, and
$a^{-2}\Tr e^{ag-tP}$ is not entire.

The last ingredient is a condition on explicit Clausen kernels. Let $N=4d$ and let
$\Cl(\theta)=\sum_{n\ge1}n^{-2}\sin n\theta$ be the Clausen function. Put
\[
\Xi(u)=-\frac{N^2}{2\pi^2}\Cl\Bigl(\frac{2\pi u}{N}\Bigr),\qquad \Sh(u)=\Xi(u)+\Xi(2d-u)\qquad(0\le u\le2d),
\]
so that $(\Sh)''(u)=2\csc\frac{\pi u}{2d}$ on $(0,2d)$. Let $\Delta_d=\{\ell\in\R^d:\ell_r\ge0,\ \sum_r\ell_r=d\}$
and $\Delta_d^+=\{\ell\in\Delta_d:\ell_r>0\text{ for all }r\}$. For $\ell\in\Delta_d$, extended $d$-periodically, let
$S_r(m)=\ell_r+\ell_{r+1}+\dots+\ell_{r+m-1}$ and define the \emph{cell vector} $u^\ell\in\R^{d-1}$ by
\begin{equation}\label{eq:uell}
u^\ell_m=\Delta^2_m\Bigl[\frac1d\sum_{r=0}^{d-1}\Sh\bigl(S_r(m)\bigr)\Bigr]\qquad(1\le m\le d-1),
\end{equation}
where $\Delta^2_m f=f(m+1)-2f(m)+f(m-1)$. Let $v\in\R^{d-1}$, $v_m=2\csc\frac{\pi m}{2d}$, and
$t_s=e_s+e_{d-s}$ ($1\le s\le d-1$), where $(e_m)$ is the standard basis.

\begin{introthm}[cone condition; computer-assisted]\label{thm:E}
For every $d\ge2$ there are a finite positive Borel measure $\mu$ on $\Delta_d$ with $\mu(\Delta_d^+)=\mu(\Delta_d)>0$
and numbers $w_s\ge0$ such that
\begin{equation}\label{eq:cone}
v=\int_{\Delta_d}u^\ell\,\mu(\dd\ell)+\sum_{s=1}^{\lfloor d/2\rfloor}w_s\,t_s .
\end{equation}
\end{introthm}

For $2\le d\le200$ the measure $\mu$ is a finite sum of point masses at rational cell vectors and $w=0$; the
certificates are verified in interval arithmetic, and for $d\le20$ also by the Lean kernel in exact rational interval
arithmetic. For $d\ge201$, $\mu$ is the law of an explicit random cell vector
(a Dirichlet vector with Gaussian-modulated parameters), and the verification is by interval arithmetic for each
$201\le d\le2000$ and by a single interval Taylor-model computation, uniform in $d$, for $d\ge2001$
(Section~\ref{sec:cone}).

\subsection{Status of Problem 27}
Table~\ref{tab:status} lists every clause of Problem~27 as originally posed~\cite{OQP} and its status after this work.
Part~A was answered negatively by Bancal, Gisin and Pironio~\cite{BGP}. Theorems~\ref{thm:A} and~\ref{thm:B} settle
the clause of Part~B on optimality and uniqueness, for projective measurements. The two remaining clauses of Part~B
are settled in Section~\ref{sec:clauses}. The verdict on the noise clause depends on how noise resistance is measured.
If it refers to the violation of the CGLMP inequality, with uniformly random outcomes as noise, the clause holds for
every $d$, and the DKZ measurements are the unique optimum (Corollary~\ref{cor:noise}). In Gill's literal reading, in
which the noisy correlations must violate some Bell inequality, it fails for every $d\ge4$ on $\Phi_d$
(Theorem~\ref{thm:literal}); it holds for $d=2$, and for $d=3$ it holds on $\Phi_3$ but fails on $\Phi_2$
(Proposition~\ref{prop:small}). A strengthening, with white noise on the state and every outcome used, remains open
(Section~\ref{subsec:white}). The Kullback--Leibler clause fails for every $d\ge4$ (Theorem~\ref{thm:kl}); for $d=3$
the DKZ measurements are a strict local maximum (Proposition~\ref{prop:kl3}). Thus every clause of Problem~27 as
originally posed now has a verdict, the verdict on the noise clause depending on its reading.

\begin{table}[ht]
\small
\begin{tabular}{@{}>{\raggedright\arraybackslash}p{0.31\textwidth}>{\raggedright\arraybackslash}p{0.31\textwidth}>{\raggedright\arraybackslash}p{0.31\textwidth}@{}}
\toprule
Clause & Status & References\\
\midrule
27A: every facet of the local polytope that is not contained in a face of the no-signalling polytope is of CGLMP
type, possibly lifted & false ($d=4$) & Bancal--Gisin--Pironio~\cite{BGP}\\
27B: the DKZ measurements are optimal for $I_d$ on maximally entangled states & proved for every $d$ (projective
measurements; computer-assisted for $d\ge21$; formally verified in Lean for $d\le20$) & Theorem~\ref{thm:A}; earlier:
$d=2$~\cite{Tsirelson}, $d\le20$~\cite{repo-cglmp}, commuting (equal-link) strategies for all $d$~\cite{MS-classical}\\
27B: they are \emph{necessarily} the optimal ones & proved for every $d$ (projective; computer-assisted for $d\ge21$;
formally verified in Lean for $d\le20$) & Theorem~\ref{thm:B}; earlier: $d\le20$~\cite{repo-cglmp}, equal-link
strategies~\cite{MS-classical}\\
27B: highest noise resistance, (a) violation of the CGLMP inequality, uniform outcome noise & true for every $d$, DKZ
the unique optimum (projective; formally verified, with the scope of Theorems~\ref{thm:A} and~\ref{thm:B}) &
Corollary~\ref{cor:noise}\\
27B: highest noise resistance, (b) Gill's reading: violation of local realism, uniform outcome noise & false for every
$d\ge4$ on $\Phi_d$ (computer-assisted for $4\le d\le9$); true for $d=2$; $d=3$: true on $\Phi_3$, false on $\Phi_2$ &
Theorem~\ref{thm:literal}, Proposition~\ref{prop:small}; mechanism~\cite{ADGL}; numerics on $\Phi_4$,
$\Phi_{16}$~\cite{BRL}\\
(c) strengthening of (b): white noise on the state, every outcome used & open for $d\ge4$; DKZ optimal in numerical
searches for $d\le8$ & Section~\ref{subsec:white}; numerics:~\cite{GLZ} ($d\le5$), \cite{repo-oqp27}\\
27B: best Kullback--Leibler discrimination & false for every $d\ge4$ (computer-assisted); $d=3$: DKZ a strict local
maximum, global optimality numerical & Theorem~\ref{thm:kl}, Proposition~\ref{prop:kl3}; $d=4$ first by
Zhang~\cite{Zhang}; $4\le d\le9$~\cite{repo-cglmp}\\
Related: maximal value of $I_d$ over all states & exact for $d=3,\dots,8$ & Ioannou--Rosset~\cite{IR} ($d=3,4$),
\cite{repo-cglmp} ($d=5,\dots,8$)\\
\bottomrule
\end{tabular}
\caption{The clauses of Open Quantum Problem~27 as originally posed~\cite{OQP}, and their status. ``Projective'' means
that POVMs are not covered by Theorems~\ref{thm:A} and~\ref{thm:B}; for $3\le d\le8$ optimality and uniqueness also hold for
arbitrary POVMs (Section~\ref{sec:open}, item~(1)). The counterexamples use projective measurements.}\label{tab:status}
\end{table}

\subsection{Strategy of the proof}\label{subsec:strategy}
The proof proceeds in six steps; Sections~\ref{sec:reduction}--\ref{sec:proofsAB} follow this order.

\emph{(1) The clock model.} By an exact reduction from~\cite{MS-classical}, recalled in Section~\ref{sec:reduction},
$I_d=4F(V)/(d(d-1))$ for a $d$-tuple $V$ of unitaries of order four on some $\C^M$, and $\IME(d)=4\FD/(d(d-1))$. The
spectral projections of the $V_k$ form a family $(Q_x)_{x\in\Z_{4d}}$ of projections, each residue class modulo $d$
being a four-outcome projection-valued measure, and $F(V)-\FD=\frac12\ip v\delta$, where $\delta\in\R^{d-1}$ collects
the deviations of $d-1$ ``net pair counts'' from their values at the DKZ configuration (Lemma~\ref{lem:linearform}).

\emph{(2) Two-variable BMV positivity and the strip inequality} (Section~\ref{sec:strip}): Theorems~\ref{thm:D}
and~\ref{thm:C}, proved by hand. Appendix~\ref{app:RI} gives a second proof of both, without distribution theory or
the Paley--Wiener--Schwartz theorem, through a Radon identity for matrix pencils.

\emph{(3) A noncommutative continuum theorem} (Section~\ref{sec:continuum}): for projection-valued step fields $B$ on
the circle, fields $0\le A\le\one-B$, and the conjugate function $g$ of $B$, the pairing $\int\tau(Ag)$ is at most
the value of the classical (commutative) problem. This is a matrix version of the circle Stein--Weiss
rearrangement theorem~\cite{SW,Laeng}; the strip inequality enters pointwise, through a harmonic function on the disc.

\emph{(4) Cell embeddings} (Section~\ref{sec:cells}): placing the projections $Q_x$ on cells of lengths
$\ell_{x\bmod d}$ on a circle of length $4d$ produces a step field to which (3) applies; averaging over rotations
gives the linear inequalities $\ip{u^\ell}\delta\le0$. Consequently Theorem~\ref{thm:E} implies $\ip v\delta\le0$,
which is Theorem~\ref{thm:A}.

\emph{(5) The cone condition} (Section~\ref{sec:cone}): Theorem~\ref{thm:E}, computer-assisted.

\emph{(6) Rigidity} (Section~\ref{sec:proofsAB}): in the equality case every inequality of the chain is an equality.
Through the equality clause of Theorem~\ref{thm:C} this forces $[B(\theta),g(\theta)]=0$ almost everywhere on a cell
embedding with all cells of positive length; a residue argument then shows that all $Q_x$ commute, and the
commutative case~\cite{MS-classical} identifies the strategy.

The continuum problem in (3) is degenerate: its maximum is attained by many configurations, including non-commuting
ones (Remark~\ref{rem:degenerate}), and the discrete problem differs from it by terms of order $d$ located at the
junctions between blocks. Step~(4) avoids any explicit estimate of these junction terms: the cell inequalities are
exact consequences of the continuum theorem, and the discrete kernel is recovered from them by the cone condition.
The mechanism behind the cone condition is that the discrete cotangent kernel is the continuum kernel averaged over
Poisson-like random cells (Remark~\ref{rem:poisson}).

\subsection{Relation to previous work}\label{subsec:previous}
For $d=2$, Theorem~\ref{thm:A} is Tsirelson's bound~\cite{Tsirelson} on maximally entangled states and
Theorem~\ref{thm:B} is the corresponding rigidity. For $d=3$, Lang, V\'ertesi and Navascu\'es~\cite{LVN} showed
with a semidefinite relaxation that the maximal value on maximally entangled states agrees with $\IME(3)$ to within
$10^{-10}$ (this is the partial result listed in~\cite{OQP}). In~\cite{repo-cglmp} we gave exact sum-of-squares
certificates proving Theorems~\ref{thm:A} and~\ref{thm:B} for $3\le d\le20$, and in~\cite{MS-classical} we proved
both theorems for every $d$ for strategies whose reduced family commutes (equivalently, whose four ``link''
operators coincide). Ioannou and Rosset~\cite{IR} obtained exact sum-of-squares bounds for the maximal value of $I_3$
and $I_4$ over all states. The SATWAP expressions of Salavrakos et al.~\cite{SATWAP} are maximally violated by
maximally entangled states with measurements of CGLMP type and underlie device-independent self-testing
results~\cite{SSKA}; they are different Bell expressions.

For the noise clause, the critical visibilities $2/\IME(d)$ of the DKZ measurements against all local models were
computed numerically by Kaszlikowski et al.~\cite{KGZMZ} ($d\le9$) and by Durt, Kaszlikowski and
\.Zukowski~\cite{DKZ} ($d\le16$), and identified with the threshold of the CGLMP inequality in~\cite[eq.~(23)]{CGLMP};
Chen et al.~\cite{Chen01} treated $d=3$ analytically, and Gruca, Laskowski and \.Zukowski~\cite{GLZ} recomputed the
thresholds numerically for $d\le5$ and found that states of Schmidt rank two are the most noise-robust. The mechanism
behind Theorem~\ref{thm:literal} goes back to Ac\'{\i}n, Durt, Gisin and Latorre~\cite{ADGL}, and Baek, Ryu and
Lee~\cite{BRL} observed numerically that coarse-grained measurements beat DKZ on $\Phi_4$ and $\Phi_{16}$ under uniform
noise (Remark~\ref{rem:literalcredit}). The Kullback--Leibler clause refers to the statistical strength of van Dam, Gill
and Gr\"unwald~\cite{vDGG}; Zhang~\cite{Zhang} found the first counterexample ($d=4$), and in~\cite{repo-cglmp} we gave
counterexamples for $4\le d\le9$.

The mechanism of our proof of Theorem~\ref{thm:D}---inverting a trace function of a pencil of Hermitian matrices
by a two-dimensional Fourier transform computed along lines, and identifying the density as
$\frac1{2\pi}\sum|\Im|$ of the roots of an auxiliary pencil by an integral identity for $\log|\det|$---is the one
introduced by Hein\"avaara~\cite{He} for tracial joint spectral measures, which yields, among other
results, a new proof of Stahl's theorem~\cite{Stahl} (see also~\cite{Eremenko,LS}). Theorem~\ref{thm:D} is a variant
adapted to a projection and to an affine chart: the subtraction of the two pinched terms and the division by $a^2$
produce a measure without singular part, which is what the strip inequality requires. We found Theorem~\ref{thm:D}
independently, but we do not claim novelty for the mechanism, and we give complete proofs. We have not found
Theorem~\ref{thm:C}, or an equivalent statement, in the literature. Its commutative case is an identity, and its
continuum consequence (Theorem~\ref{thm:continuum}) contains the circle version of the theorem of Stein and
Weiss~\cite{SW} on the distribution of the conjugate function of an indicator function, in the form of
Laeng~\cite{Laeng} (see also~\cite{Boole,CLM}).

Appendix~\ref{app:RI} gives a second proof of Theorem~\ref{thm:D} that uses no distributions, no Paley--Wiener--Schwartz
theorem and no Fourier inversion in two variables. Its key step is a Radon identity (Theorem~\ref{thm:RI}): the integral over
$\tau\in(0,1)$ of $\frac1{2\pi}\sum_i|\Im y_i(\tau)|$, where the $y_i(\tau)$ are the eigenvalues of $(P-\tau)^{-1}H$, equals
$\Tr H_+-\Tr(BHB+PHP)_+$ for every Hermitian $H$. We prove it with a complex inviscid Burgers equation satisfied by the
pencil roots. The same argument applies to an arbitrary Hermitian matrix in place of the projection, and then gives a
proof of Stahl's theorem without distribution theory or Fourier analysis (Remark~\ref{rem:RIgeneralB}); we make no
priority claim about this by-product, and refer to~\cite{Stahl,Eremenko,He} for the history of the theorem and of the
mechanism. This elementary route is the one followed by the Lean formalisation.

\subsection{Conventions}\label{subsec:conventions}
$\Tr$ is the trace on $M_M(\C)$ and $\tau=\frac1M\Tr$. Spectra are multisets, counted with algebraic multiplicity.
$\|X\|$ is the operator norm and $\|X\|_F$ the Frobenius norm. For $f\in L^1(\R)$,
$\hat f(\sigma)=\int_\R f(w)e^{-i\sigma w}\dd w$; the Hilbert transform is
$\cH f(w)=\frac1\pi\,\pv\!\int_\R\frac{f(u)}{w-u}\dd u$, which acts on $L^2(\R)$ as the Fourier multiplier
$-i\operatorname{sgn}\sigma$. On the circle $\T=\R/2\pi\Z$ we use the normalised measure $\dd m=\dd\theta/2\pi$, and
the conjugate function of $f\in L^2(\T)$ is $\pv\!\int\cot\frac{\theta-\varphi}2f(\varphi)\dd m(\varphi)$, the
multiplier $-i\operatorname{sgn}n$. Catalan's constant is $G=\Cl(\pi/2)=\sum_{n\ge0}(-1)^n(2n+1)^{-2}$, and
$\Cl(\theta)=-\int_0^\theta\log|2\sin\frac t2|\dd t$. For $\alpha,\beta\ge0$ with $\alpha+\beta\le1$ we put
\begin{equation}\label{eq:Phic}
\Phic(\alpha,\beta)=\frac1{2\pi^2}\Bigl[\Cl(2\pi\alpha)+\Cl(2\pi\beta)+\Cl\bigl(2\pi(1-\alpha-\beta)\bigr)\Bigr],
\qquad \Phic\bigl(\tfrac14,\tfrac14\bigr)=\frac{G}{\pi^2}.
\end{equation}
Results labelled \emph{computer-assisted} depend on computations in rigorous interval arithmetic with outward
rounding, or in exact rational or algebraic arithmetic (Section~\ref{sec:clauses}); everything else is proved by hand. Numerical observations that are not part of any proof are labelled as
such. The Lean formalisation (Section~\ref{subsec:formal}) is an independent check; no proof in this paper depends on
it.

%======================================================================
\section{The clock model and the Q-configuration}\label{sec:reduction}
%======================================================================

Throughout, $d\ge2$, $N=4d$, $w=e^{2\pi i/d}$, $z=e^{2\pi i/N}$ (so $z^4=w$ and $z^d=i$), and
$\kappa(u)=\cot(\pi u/N)$ for $u\in\Z_N\setminus\{0\}$, $\kappa(0)=0$.

\subsection{Strategies and the clock model}
Encode a projective strategy on $\Phi_D$ by the unitaries $U_x=\sum_aw^aP^x_a$ and $U'_y=\sum_bw^bQ^y_b$ and put
\begin{equation}\label{eq:chainR}
R_1=U_2,\qquad R_2=(U'_2)^T,\qquad R_3=U_1,\qquad R_4=(U'_1)^T,
\end{equation}
the transpose being taken in the Schmidt basis of $\Phi_D$; then $R_i^d=\one$, and conversely every quadruple of
unitaries on $\C^D$ with $R_i^d=\one$ comes from a projective strategy on $\Phi_D$, through spectral projections.
A simultaneous conjugation $R_i\mapsto uR_iu^*$ is the same as applying the local unitary $u\otimes\bar u$, which fixes
$\Phi_D$. For unitaries $V_0,\dots,V_{d-1}\in U(M)$ with $V_k^4=\one$ let
\begin{equation}\label{eq:FV}
F(V)=\sum_{0\le j<k\le d-1}\Re\Bigl[h_{k-j}\,\tau(V_jV_k^*)\Bigr],\qquad h_m=\sec\frac{\pi m}{2d}-i\csc\frac{\pi m}{2d},
\end{equation}
and $\FD=F(\one,\dots,\one)=\sum_{m=1}^{d-1}(d-m)\sec\frac{\pi m}{2d}$. Let $X$ be the cyclic shift
$X|k\rangle=|k+1\rangle$ on $\C^d$ and define the \emph{covariant strategy} $R^V$ on $\C^d\otimes\C^M$ by
\begin{equation}\label{eq:covariant}
W=\sum_{k=0}^{d-1}|k\rangle\langle k|\otimes z^kV_k,\qquad R^V_1=X\otimes\one_M,\qquad R^V_{i+1}=WR^V_iW^*\quad(i=1,2,3).
\end{equation}

\begin{theorem}[reduction to the clock model {\cite[Theorem~2.3]{MS-classical}}]\label{thm:reduction}
Let $d\ge2$.
\begin{enumerate}
\item For every covariant strategy, $I_d(R^V)=4F(V)/(d(d-1))$.
\item For every projective strategy $R$ on $\Phi_D$ there are $M$ (one can take $M=4D$) and unitaries
$V_0,\dots,V_{d-1}\in U(M)$ with $V_k^4=\one$, unique up to a simultaneous unitary conjugation, such that $R$ is
unitarily equivalent (by a simultaneous conjugation of $R_1,\dots,R_4$) to a direct summand of $R^V$ and
$I_d(R)=I_d(R^V)=4F(V)/(d(d-1))$.
\end{enumerate}
In particular $I_d\le\IME(d)$ holds for all projective strategies on maximally entangled states if and only if
$F(V)\le\FD$ for all $M$ and all $V_0,\dots,V_{d-1}\in U(M)$ with $V_k^4=\one$.
\end{theorem}

The family $V$ in (2) is the \emph{reduced family} of $R$; it is obtained from a twirl over a cyclic symmetry of order
$4d$ followed by the Stone--von Neumann theorem for $\Z_d$.

\subsection{The DKZ strategy and the commutative sector}\label{subsec:dkz}
Let $W_0=\mathrm{diag}(z^k)_{k=0}^{d-1}$ and $R^{(0)}=R^V$ for $M=1$, $V\equiv1$. The DKZ measurements are the
measurements in the orthonormal bases $|k\rangle_{A,x}=d^{-1/2}\sum_je^{-2\pi ij(k+\alpha_x)/d}|j\rangle$ and
$|l\rangle_{B,y}=d^{-1/2}\sum_je^{2\pi ij(l-\beta_y)/d}|j\rangle$ with $\alpha=(0,\frac12)$ and
$\beta=(\frac14,-\frac14)$ (the measurements of~\cite[eqs.~(12)--(15)]{CGLMP}), and their chain~\eqref{eq:chainR} is
$GR^{(0)}G^*$ with $G=W_0^{-2}$~\cite[Section~2.4]{MS-classical}. Hence $I_d(\mathrm{DKZ})=4\FD/(d(d-1))=\IME(d)$.

\begin{theorem}[commutative sector {\cite[Theorem~B, Corollary~C, Proposition~2.4]{MS-classical}}]\label{thm:commutative}
Let $d\ge2$.
\begin{enumerate}
\item If $V_0,\dots,V_{d-1}\in U(M)$ commute pairwise and $V_k^4=\one$, then $F(V)\le\FD$, with equality if and only if
every joint eigenvector $v$ of the $V_k$ satisfies $V_kv=i^{-a_k}v$ with a \emph{one-step} configuration
$a_k=c+[k\ge r]$ ($c\in\Z_4$, $0\le r\le d-1$; $c$ and $r$ may depend on $v$).
\item A projective strategy on $\Phi_D$ has commuting reduced family if and only if its four link operators
coincide, $R_1^nR_2^{-n}=R_2^nR_3^{-n}=R_3^nR_4^{-n}=w^{-n}R_4^nR_1^{-n}$ for all $n$.
\item If a projective strategy on $\Phi_D$ with commuting reduced family attains $I_d=\IME(d)$, then $d\mid D$ and
there is a unitary $u\colon\C^D\to\C^d\otimes\C^{D/d}$ with $uR_iu^*=R_i^{\mathrm{DKZ}}\otimes\one$ for $i=1,\dots,4$.
\end{enumerate}
\end{theorem}

\subsection{Q-configurations}
\begin{definition}\label{def:Qconf}
A \emph{Q-configuration} on $\C^M$ is a family $(Q_x)_{x\in\Z_N}$ of orthogonal projections on $\C^M$ such that for every
$k\in\{0,\dots,d-1\}$ the four projections $Q_{k-da}$, $a\in\Z_4$, are pairwise orthogonal with sum $\one$. We put
$A_x=Q_x$ and $B_x=Q_{x+d}$, and for $j\in\Z_N$ we let $(\mathrm{rot}_jQ)_x=Q_{x-j}$.
\end{definition}

Every element of $\Z_N$ is uniquely of the form $k-da$ with $0\le k\le d-1$ and $a\in\Z_4$. Given a family $V$ as
above, let $Q_{k-da}$ be the spectral projection of $E_k:=V_k^*$ for the eigenvalue $i^a$; this is a
Q-configuration, $E_k=\sum_{a\in\Z_4}i^aQ_{k-da}$, and conversely every Q-configuration arises from a unique family
$V$ in this way. The following properties are immediate: $A_xB_x=0$ (two outcomes of one measurement); for every
$d$-periodic function $\ell$ on $\Z_N$,
\begin{equation}\label{eq:opconstraint}
\sum_{x\in\Z_N}\ell_x\,B_x=\sum_{x\in\Z_N}\ell_x\,A_x=\Bigl(\sum_{r=0}^{d-1}\ell_r\Bigr)\one ;
\end{equation}
and every rotation $\mathrm{rot}_jQ$ is again a Q-configuration. The Q-configuration of $V\equiv\one$ ($M=1$) is the
\emph{window} $Q_x=\one_{[0,d)}(x)$; its rotations are the windows $\one_{[c,c+d)}$.

\begin{lemma}[projection form]\label{lem:projform}
For every family $V$ and its Q-configuration,
\[
F(V)=\frac12\ip AB_\tau-\frac d2,\qquad \ip AB_\tau:=\sum_{x,y\in\Z_N}\kappa(x-y)\,\tau(A_xB_y).
\]
\end{lemma}

\begin{proof}
This is~\cite[eq.~(13)]{MS-classical}; we include the short proof. For $j<k$, $m=k-j$ and $a,b\in\Z_4$ put
$x=j-da$, $y=k-db$. Then $\tau(V_jV_k^*)=\tau(E_j^*E_k)=\sum_{a,b}i^{b-a}\tau(Q_xQ_y)$ and
$\Re[h_mi^{s}]=\sec(\frac{\pi m}{2d}-\frac{\pi s}2)$ for $s\in\Z_4$ (check $s=0,1,2,3$), which equals
$\sec\frac{\pi(y-x)}{2d}$ because $y-x=m-d(b-a)$ and $u\mapsto\sec\frac{\pi u}{2d}$ is even with period $4d$. Since
$\tau(Q_xQ_y)=\tau(Q_yQ_x)$,
\[
F(V)=\frac12\sum_{x\not\equiv y\ (d)}\tau(Q_xQ_y)\sec\frac{\pi(x-y)}{2d}.
\]
The identity $\sec t=\frac12[\cot(\frac t2+\frac\pi4)-\cot(\frac t2-\frac\pi4)]$ with $t=2\pi(x-y)/N$ gives
$\sec\frac{\pi(x-y)}{2d}=\frac12[\kappa(x-y+d)-\kappa(x-y-d)]$ for $x\not\equiv y\pmod d$, and the two halves give equal
sums (exchange $x$ and $y$ and use that $\kappa$ is odd). Hence
$F(V)=\frac12\sum_{x\not\equiv y}\tau(Q_xQ_y)\kappa(x-y+d)=\frac12\sum_{x\not\equiv y'}\tau(A_xB_{y'})\kappa(x-y')$ with
$y'=y-d$. In the full sum $\ip AB_\tau$ the pairs with $x\equiv y'\pmod d$ contribute only for $x=y'+d$ (otherwise
$Q_xQ_{y'+d}=0$, or $\kappa(0)=0$ when $x=y'$), and these contribute $\kappa(d)\sum_x\tau(Q_x)=d$. So
$\ip AB_\tau=2F(V)+d$.
\end{proof}

\subsection{Net pair counts}
For a Q-configuration put
\[
T(m)=\sum_{y\in\Z_N}\tau(A_{y+m}B_y),\qquad \cN(m)=T(m)-T(-m),\qquad \delta_m=\cN(m)-m\quad(1\le m\le d-1),
\]
and $\delta=(\delta_m)_{m=1}^{d-1}$. For sites $0\le j<k\le d-1$ let
$\pi_{jk}(t)=\sum_{b\in\Z_4}\tau(Q_{j-db}Q_{k-d(b+t)})$, $t\in\Z_4$; this is a probability vector on $\Z_4$
(nonnegative, since $\tau(QQ')\ge0$ for projections, with sum $\sum_b\tau(Q_{j-db})=1$).

\begin{lemma}[linear form]\label{lem:linearform}
For every Q-configuration:
\begin{enumerate}
\item $\cN(2d-m)=\cN(m)$ for $1\le m\le 2d-1$, and $\cN(d)=d$;
\item $F(V)=\sum_{m=1}^{d-1}\csc\frac{\pi m}{2d}\,\cN(m)$, the window has $\cN(m)=m$, and
\[
F(V)-\FD=\frac12\ip v\delta,\qquad v_m=2\csc\frac{\pi m}{2d};
\]
\item for $1\le m\le d-1$,
$\cN(m)=\sum_{k-j=m}[\pi_{jk}(1)-\pi_{jk}(3)]+\sum_{k-j=d-m}[\pi_{jk}(0)-\pi_{jk}(2)]$; consequently
$\ip{t_s}\delta=\cN(s)+\cN(d-s)-d\le0$ for $1\le s\le d-1$.
\end{enumerate}
\end{lemma}

\begin{proof}
(1) Put $C(n)=\sum_u\tau(Q_uQ_{u+n})$; then $C(-n)=C(n)$ by cyclicity of the trace, and
$T(m)=\sum_y\tau(Q_{y+m}Q_{y+d})=C(d-m)$, so $\cN(m)=C(m-d)-C(m+d)$. Hence
$\cN(2d-m)=C(d-m)-C(3d-m)=C(m-d)-C(m+d)=\cN(m)$ (indices modulo $N$), and
$\cN(d)=C(0)-C(2d)=\sum_u\tau(Q_u)-0=d$, because $Q_u$ and $Q_{u+2d}$ are distinct outcomes of one measurement.

(2) Since $\kappa$ is odd and $\kappa(2d)=0$, $\ip AB_\tau=\sum_{m\in\Z_N}\kappa(m)T(m)=\sum_{m=1}^{2d-1}\kappa(m)\cN(m)$,
and by (1) this is $\kappa(d)d+\sum_{m=1}^{d-1}(\kappa(m)+\kappa(2d-m))\cN(m)=d+\sum_mv_m\cN(m)$, because
$\cot\frac{\pi m}{4d}+\tan\frac{\pi m}{4d}=2\csc\frac{\pi m}{2d}$. Lemma~\ref{lem:projform} gives the formula for $F(V)$.
For the window, $T(m)=\#\{y\in[-d,0):y+m\in[0,d)\}=m$ and $T(-m)=0$ for $1\le m\le d-1$, so $\cN(m)=m$; thus
$\FD=\sum_mm\csc\frac{\pi m}{2d}$ (which agrees with $\sum_m(d-m)\sec\frac{\pi m}{2d}$ after $m\mapsto d-m$), and the
formula for $F(V)-\FD$ follows.

(3) Write $y+d=j-db$ with $0\le j\le d-1$. If $j+m<d$, then $y+m=(j+m)-d(b+1)$ is the outcome $b+1$ of site $k=j+m$; if
$j+m\ge d$, then $y+m=(j+m-d)-db$ is the outcome $b$ of site $k'=j+m-d<j$. Summing over $b$ gives
$T(m)=\sum_{k-j=m}\pi_{jk}(1)+\sum_{k-j=d-m}\pi_{jk}(0)$; in the same way
$T(-m)=\sum_{k-j=m}\pi_{jk}(3)+\sum_{k-j=d-m}\pi_{jk}(2)$. Hence
$\cN(s)+\cN(d-s)=\sum[\pi_{jk}(0)+\pi_{jk}(1)-\pi_{jk}(2)-\pi_{jk}(3)]$, the sum running over the $d$ pairs at distance
$s$ or $d-s$ (counted twice if $s=d/2$), and each term is at most $1$.
\end{proof}

By Theorem~\ref{thm:reduction} and Lemma~\ref{lem:linearform}, Theorem~\ref{thm:A} is equivalent to
\begin{equation}\label{eq:target}
\ip v\delta\le0\qquad\text{for every Q-configuration, on every }\C^M .
\end{equation}
The inequalities $\ip{t_s}\delta\le0$ of Lemma~\ref{lem:linearform}(3) hold pairwise and carry no information
about $\ip v\delta$ by themselves; they are the ``trivial directions'' in Theorem~\ref{thm:E}.

%======================================================================
\section{Two-variable BMV positivity and the strip inequality}\label{sec:strip}
%======================================================================

Throughout this section $M\ge1$, $B$ is an orthogonal projection on $\C^M$ of rank $r$, $P=\one-B$, $g$ is Hermitian,
$g_d=BgB+PgP$ is the pinching of $g$, $\lmin\le\lmax$ are the extreme eigenvalues of $g$, and
$K=[\lmin,\lmax]\times[0,1]\subseteq\R^2$.

\subsection{The strip, its Poisson kernel and the functions \texorpdfstring{$h_\lambda$}{h}}

\begin{lemma}[strip Poisson kernel]\label{lem:poisson}
Let $0<x<1$ and let $K_x$ be as in~\eqref{eq:poisson}.
\begin{enumerate}
\item $K_x>0$ and $\int_\R K_x=1-x$. For $|u|\ge1$, $K_x(u)\le2\sin(\pi x)e^{-\pi|u|}$; if $x\ge\frac12$, then
$K_x(u)\le\sin(\pi x)e^{-\pi|u|}$ for all $u$.
\item For $u\ne0$, $K_x(u)=\sum_{n\ge1}\sin(n\pi x)e^{-n\pi|u|}$.
\item For $a\in\C$ with $|\Re a|<\pi$,
$\displaystyle\int_\R e^{-au}K_x(u)\dd u=\frac{\sin(a(1-x))}{\sin a}$ (interpreted as $1-x$ at $a=0$).
\end{enumerate}
\end{lemma}

\begin{proof}
(1) $\cosh\pi u\ge1>\cos\pi x$. For $|u|\ge1$, $\cosh\pi u-\cos\pi x\ge\cosh\pi u-1\ge\frac12e^{\pi|u|}-1\ge\frac14e^{\pi|u|}$;
for $x\ge\frac12$, $\cos\pi x\le0$ and $\cosh\pi u-\cos\pi x\ge\frac12e^{\pi|u|}$. The mass is (3) at $a=0$.

(2) With $q=e^{-\pi|u|}\in(0,1)$ and $\theta=\pi x$,
$\sum_nq^n\sin n\theta=\Im\frac{qe^{i\theta}}{1-qe^{i\theta}}=\frac{q\sin\theta}{1-2q\cos\theta+q^2}=\frac{\sin\theta}{2\cosh\pi u-2\cos\theta}$.

(3) First let $a=i\kappa$ with $\kappa\in\R\setminus\{0\}$, and $f(u)=e^{-i\kappa u}/(\cosh\pi u-\cos\pi x)$. Integrate $f$ over the
boundary of the rectangle $[-L,L]\times[0,2]$; the vertical sides tend to $0$ as $L\to\infty$, and
$f(u+2i)=e^{2\kappa}f(u)$, so the contour integral tends to $(1-e^{2\kappa})\int_\R f$. Inside the rectangle the
denominator vanishes at $u=ix$ and $u=i(2-x)$, simply, with residues of $f$ equal to $e^{\kappa x}/(i\pi\sin\pi x)$ and
$-e^{\kappa(2-x)}/(i\pi\sin\pi x)$. Hence
$\int_\R f=\frac{2(e^{\kappa x}-e^{\kappa(2-x)})}{(1-e^{2\kappa})\sin\pi x}=\frac{2\sinh(\kappa(1-x))}{\sinh(\kappa)\sin\pi x}$, that is,
$\int e^{-i\kappa u}K_x(u)\dd u=\frac{\sinh(\kappa(1-x))}{\sinh\kappa}=\frac{\sin(a(1-x))}{\sin a}$. Both sides of (3) are
holomorphic on $|\Re a|<\pi$ (the left side by the exponential decay in (1); on the right the only zero of $\sin a$ in
this strip, $a=0$, is removable), and they agree on $i\R\setminus\{0\}$, hence everywhere.
\end{proof}

For $\lambda\in\R$ let
\begin{equation}\label{eq:philambda}
\phi_\lambda(z)=\frac{i}{\pi^2}\Li\bigl(e^{-i\pi z-\pi\lambda}\bigr)\qquad(z\in S).
\end{equation}
For $z=x+iy\in S$ the argument $e^{-i\pi z-\pi\lambda}=e^{\pi(y-\lambda)}e^{-i\pi x}$ has argument $-\pi x\in(-\pi,0)$, so it
avoids the cut $[1,\infty)$; thus $\phi_\lambda$ is holomorphic on $S$, and $\Re\phi_\lambda=h_\lambda$ by~\eqref{eq:hlambda}.
Moreover $h_\lambda(z)=h_0(z-i\lambda)$.

\begin{lemma}[the strip functions]\label{lem:h}
Let $\lambda\in\R$.
\begin{enumerate}
\item For $z=x+iy\in S$, $h_\lambda(z)=\int_\R(w-\lambda)_+K_x(y-w)\dd w$.
\item $h_\lambda$ extends continuously to the closed strip $\bar S$, with $h_\lambda(iy)=(y-\lambda)_+$ and
$h_\lambda(1+iy)=0$. (For $y>\lambda$ the point $iy$ corresponds to an argument of $\Li$ on the cut $[1,\infty)$; the
boundary value is the limit from inside $S$, $\Im\Li(\varrho-i0)=-\pi\log\varrho$ for $\varrho>1$.)
\item $0\le h_\lambda(x+iy)\le(1-x)(y-\lambda)_++2$ on $\bar S$.
\item For fixed $x\in(0,1)$, $\partial_y^2h_\lambda(x+iy)=K_x(y-\lambda)>0$. For fixed $y$, $x\mapsto h_\lambda(x+iy)$ is
concave on $[0,1]$.
\item For $\nu=x+iy\in\bar S$ let $\omega_\nu$ be the measure $K_x(y-w)\dd w$ if $0<x<1$, $\delta_y$ if $x=0$, and $0$ if
$x=1$. Then $h_\lambda(\nu)=\int(w-\lambda)_+\,\omega_\nu(\dd w)$, and for $a\in\C$ with $|\Re a|<\pi$,
$\int e^{aw}\omega_\nu(\dd w)=e^{ay}\sin(a(1-x))/\sin a$.
\end{enumerate}
\end{lemma}

\begin{proof}
(1) Let $\Phi(z)$ denote the right side. Since $\Re\cot(A+iB)=\sin2A/(\cosh2B-\cos2A)$,
$K_x(y-w)=\frac12\Re\cot\frac{\pi(z-iw)}2$, which is harmonic in $z\in S$; differentiation under the integral is
justified by Lemma~\ref{lem:poisson}(1), so $\Phi$ is harmonic on $S$. So is $h_\lambda=\Re\phi_\lambda$. For $y<\lambda$,
the integrand vanishes unless $w>\lambda>y$, where by Lemma~\ref{lem:poisson}(2)
$K_x(y-w)=\sum_n\sin(n\pi x)e^{-n\pi(w-y)}$, uniformly convergent for $w\ge\lambda$; integrating term by term,
\[
\Phi(z)=\sum_{n\ge1}\sin(n\pi x)\frac{e^{n\pi(y-\lambda)}}{(n\pi)^2}
=-\frac1{\pi^2}\Im\sum_{n\ge1}\frac{(e^{\pi(y-\lambda)}e^{-i\pi x})^n}{n^2}=h_\lambda(z).
\]
Two harmonic functions on the connected set $S$ that agree on the open set $\{y<\lambda\}$ agree everywhere.

(2) As $x\to0$, $K_x(y-\cdot)$ is an approximate identity at $y$: it is positive, has mass $1-x$, and
$\int|u|K_x(u)\dd u\to0$. Indeed $\cosh\pi u-\cos\pi x=2\sinh^2\frac{\pi u}2+2\sin^2\frac{\pi x}2\ge2(u^2+x^2)$ (as
$\sinh t\ge t$ and $\sin\frac{\pi x}2\ge x$ on $[0,1]$), so $K_x(u)\le\frac{\pi x}{4(u^2+x^2)}$ and
$\int_{|u|\le1}|u|K_x(u)\dd u\le\frac{\pi x}4\log(1+x^{-2})\to0$, while for $|u|\ge1$ we use Lemma~\ref{lem:poisson}(1). Since $|(w-\lambda)_+-(y-\lambda)_+|\le|w-y|$, $h_\lambda(x+iy')\to(y-\lambda)_+$ as
$(x,y')\to(0,y)$. As $x\to1$, Lemma~\ref{lem:poisson}(1) gives $h_\lambda(x+iy)\le\sin(\pi x)\int(w-\lambda)_+e^{-\pi|y-w|}\dd w\to0$,
locally uniformly in $y$. The statement about $\Li$ on the cut is the classical
$\Im\Li(\varrho\pm i0)=\pm\pi\log\varrho$ for $\varrho>1$.

(3) $(w-\lambda)_+\le(y-\lambda)_++|w-y|$, $\int K_x=1-x$, and
$\int|u|K_x(u)\dd u\le\int_{|u|\le1}K_x+4\int_1^\infty ue^{-\pi u}\dd u\le1+1$.

(4) For fixed $x$, $y\mapsto h_\lambda(x+iy)$ is the convolution of the convex function $(\cdot-\lambda)_+$ with the
integrable, smooth, exponentially decaying kernel $K_x$; its second derivative is the convolution of
$\delta_\lambda$ with $K_x$, that is, $K_x(y-\lambda)$. As $h_\lambda$ is harmonic, $\partial_x^2h_\lambda=-K_x(y-\lambda)<0$ on
$(0,1)$, and $x\mapsto h_\lambda(x+iy)$ is continuous on $[0,1]$ by (2).

(5) The first formula is (1) and (2); the second follows from Lemma~\ref{lem:poisson}(3) after the substitution
$u=y-w$.
\end{proof}

\subsection{The pencil}
For $s\in\R$ and $0<\tau<1$ put
\[
p_{s,\tau}(\xi)=\det\bigl(g-s-\xi(P-\tau)\bigr),\qquad q_{s,\tau}(\xi)=\det\bigl(g_d-s-\xi(P-\tau)\bigr)\qquad(\xi\in\C).
\]

\begin{lemma}[the pencil]\label{lem:pencil}
Let $s\in\R$ and $0<\tau<1$.
\begin{enumerate}
\item $p_{s,\tau}$ and $q_{s,\tau}$ have degree $M$, the same leading coefficient $c_\tau=\det(\tau-P)=(\tau-1)^{M-r}\tau^r\ne0$, and
the same sum of roots.
\item Every root $\xi$ of $p_{s,\tau}$ or of $q_{s,\tau}$ satisfies $|\xi|\le\|g-s\|/\min(\tau,1-\tau)$.
\item Every non-real root $\xi$ of $p_{s,\tau}$ satisfies $|\xi|\le\|g-s\|/\sqrt{\tau(1-\tau)}$.
\item If $s\le\lmin$ or $s\ge\lmax$, all roots of $p_{s,\tau}$ are real.
\item All roots of $q_{s,\tau}$ are real.
\item The roots of $p_{s,\tau}$, as an unordered $M$-tuple, depend continuously on $(s,\tau)\in\R\times(0,1)$.
\end{enumerate}
Consequently $\rho$ is continuous on $\R\times(0,1)$, vanishes unless $\lmin<s<\lmax$, satisfies
$\rho(s,\tau)\le\frac M{2\pi}\|g-s\|(\tau(1-\tau))^{-1/2}$, has support in $K$, and is integrable on $\R^2$.
\end{lemma}

\begin{proof}
(1) $g-s-\xi(P-\tau)=(\tau-P)\bigl[\xi-(P-\tau)^{-1}(g-s)\bigr]$, so $p_{s,\tau}(\xi)=c_\tau\det(\xi-(P-\tau)^{-1}(g-s))$ and the
roots are the eigenvalues of $(P-\tau)^{-1}(g-s)$, with sum $\Tr[(P-\tau)^{-1}(g-s)]$; likewise for $q_{s,\tau}$ with $g_d$.
Since $(P-\tau)^{-1}=\frac P{1-\tau}-\frac B\tau$ commutes with $B$, $\Tr[(P-\tau)^{-1}(BgP+PgB)]=0$, and
$g-g_d=BgP+PgB$.

(2) A root $\xi$ comes with a unit vector $v$ such that $(g-s)v=\xi(P-\tau)v$. As
$|(P-\tau)v|^2=(1-\tau)^2|Pv|^2+\tau^2|Bv|^2\ge\min(\tau,1-\tau)^2$, $|\xi|\min(\tau,1-\tau)\le\|g-s\|$. For $q_{s,\tau}$ use
$\|g_d-s\|\le\|g-s\|$.

(3) With $v$ as in (2), $\ip v{(g-s)v}=\xi\ip v{(P-\tau)v}$, and both inner products are real. If $\Im\xi\ne0$, then
$\ip v{(P-\tau)v}=0$ and $\ip v{(g-s)v}=0$; the first gives $|Pv|^2=\tau$, $|Bv|^2=1-\tau$, hence
$|(P-\tau)v|^2=\tau(1-\tau)$, and the bound follows as in (2).

(4) If $s\le\lmin$ then $g-s\ge0$, and a non-real root would give $\ip v{(g-s)v}=0$, hence $(g-s)v=0$, hence
$\xi(P-\tau)v=0$, which is impossible. The case $s\ge\lmax$ is the same.

(5) Relative to $\C^M=\operatorname{ran}B\oplus\operatorname{ran}P$, $g_d-s-\xi(P-\tau)=(BgB-s+\xi\tau)\oplus(PgP-s-\xi(1-\tau))$; it is
singular if and only if $\xi=(s-\beta)/\tau$ for an eigenvalue $\beta$ of $BgB|_{\operatorname{ran}B}$ or
$\xi=(\mu-s)/(1-\tau)$ for an eigenvalue $\mu$ of $PgP|_{\operatorname{ran}P}$.

(6) The coefficients of $p_{s,\tau}$ are polynomials in $(s,\tau)$, and the leading coefficient $c_\tau$ does not vanish on
$(0,1)$; the roots of a polynomial of fixed degree with non-vanishing leading coefficient depend continuously on the
coefficients. The final assertions follow from (3), (4), (6), and $\int_0^1(\tau(1-\tau))^{-1/2}\dd\tau=\pi$.
\end{proof}

\subsection{A Jensen-type identity}

\begin{lemma}\label{lem:jensen}
Let $p(\xi)=c\prod_{i=1}^M(\xi-\xi_i)$ and $q(\xi)=c\prod_{i=1}^M(\xi-y_i)$ with $c\ne0$, all $y_i$ real, and
$\sum_i\xi_i=\sum_iy_i$. Then $\log|p/q|\in L^1(\R)$ and
\[
\int_\R\log\Bigl|\frac{p(\xi)}{q(\xi)}\Bigr|\dd\xi=\pi\sum_{i=1}^M|\Im\xi_i| .
\]
\end{lemma}

\begin{proof}
Write $\xi_i=\alpha_i+i\beta_i$. Then $\log|p/q|=A+C$ with
$A(\xi)=\frac12\sum_i\log\bigl(1+\beta_i^2/(\xi-\alpha_i)^2\bigr)\ge0$ and $C(\xi)=\sum_i[\log|\xi-\alpha_i|-\log|\xi-y_i|]$. Each
summand of $A$ is integrable, with $\int_\R\frac12\log(1+b^2/u^2)\dd u=|b|\int_0^\infty\log(1+v^{-2})\dd v=\pi|b|$, since
$v\log(1+v^{-2})+2\arctan v$ is an antiderivative of $\log(1+v^{-2})$ with limits $0$ at $0$ and $\pi$ at $\infty$. The
function $C$ has only logarithmic singularities, and $C(\xi)=\sum_i[\log|1-\alpha_i/\xi|-\log|1-y_i/\xi|]=O(\xi^{-2})$
as $|\xi|\to\infty$ because $\sum\alpha_i=\Re\sum\xi_i=\sum y_i$; so $C\in L^1(\R)$. For $L>\max_i(|\alpha_i|,|y_i|)$ and real
$a$ with $|a|<L$,
$m_L(a):=\int_{-L}^L\log|\xi-a|\dd\xi=(L-a)\log(L-a)+(L+a)\log(L+a)-2L$ is even in $a$ with
$m_L''(a)=2L/(L^2-a^2)$, so $m_L(a)=m_L(0)+a^2/L+O(a^4/L^3)$. Hence
$\int_{-L}^LC=\sum_i(\alpha_i^2-y_i^2)/L+O(L^{-3})\to0$, and $\int_\R C=0$.
\end{proof}

\subsection{Real slices and limits}
For $\xi\in\R$ let $\lambda_1(\xi),\dots,\lambda_M(\xi)$ be the eigenvalues of $g-\xi P$ and
$\lambda^0_1(\xi),\dots,\lambda^0_M(\xi)$ those of $g_d-\xi P$, that is, the eigenvalues of $BgB|_{\operatorname{ran}B}$
together with those of $PgP|_{\operatorname{ran}P}$ shifted by $-\xi$. Put
\[
U_\xi(w)=\sum_{j=1}^M\bigl[(w-\lambda_j(\xi))_+-(w-\lambda^0_j(\xi))_+\bigr],\qquad
n_\xi(w)=\sum_{j=1}^M\bigl[\one_{w>\lambda_j(\xi)}-\one_{w>\lambda^0_j(\xi)}\bigr].
\]

\begin{lemma}[real slices]\label{lem:slices}
Let $\xi\in\R$.
\begin{enumerate}
\item $U_\xi$ is continuous, piecewise linear and compactly supported, and $U_\xi'=n_\xi$ almost everywhere; $n_\xi$ is
bounded with compact support.
\item For every $a\in\C\setminus\{0\}$, $\int_\R e^{aw}U_\xi(w)\dd w=a^{-2}\cD(a,a\xi)$.
\item For $w\notin\spec(g-\xi P)\cup\spec(g_d-\xi P)$,
$(\cH n_\xi)(w)=\frac1\pi\log\bigl|\det(g-\xi P-w)/\det(g_d-\xi P-w)\bigr|$.
\end{enumerate}
\end{lemma}

\begin{proof}
(1) $U_\xi$ vanishes below all eigenvalues, and above all of them it equals
$\sum_j(\lambda^0_j-\lambda_j)=\Tr(g_d-\xi P)-\Tr(g-\xi P)=0$. (2) As distributions,
$U_\xi''=\sum_j(\delta_{\lambda_j(\xi)}-\delta_{\lambda^0_j(\xi)})$, and $U_\xi$ has compact support, so two integrations by
parts give $\int e^{aw}U_\xi=a^{-2}\sum_j(e^{a\lambda_j(\xi)}-e^{a\lambda^0_j(\xi)})
=a^{-2}[\Tr e^{a(g-\xi P)}-\Tr_Be^{aBgB}-e^{-a\xi}\Tr_Pe^{aPgP}]=a^{-2}\cD(a,a\xi)$. (3) Pair $\lambda_j$ with $\lambda^0_j$: the
function $\one_{w>\lambda_j}-\one_{w>\lambda^0_j}$ is $\pm$ the indicator of the interval between them, and
$\frac1\pi\int_\alpha^\beta\frac{\dd u}{w-u}=\frac1\pi\log\bigl|\frac{w-\alpha}{w-\beta}\bigr|$; summing over $j$ gives
$\frac1\pi\log|\prod_j(w-\lambda_j)/\prod_j(w-\lambda^0_j)|$.
\end{proof}

\begin{lemma}[limits]\label{lem:limits}
For every Hermitian matrix $A$,
\[
\lim_{t\to+\infty}\Tr e^{A-tP}=\Tr_Be^{BAB}\qquad\text{and}\qquad\lim_{t\to-\infty}e^t\Tr e^{A-tP}=\Tr_Pe^{PAP}.
\]
\end{lemma}

\begin{proof}
Both sides of the first statement are multiplied by $e^c$ if $A$ is replaced by $A+c$, so we may assume $A\ge\one$. Let
$\mu_1\ge\dots\ge\mu_M$ be the eigenvalues of $H_t=A-tP$ and $\beta_1\ge\dots\ge\beta_r$ those of $BAB|_{\operatorname{ran}B}$.
The min-max principle with subspaces of $\operatorname{ran}B$, on which $H_t=A$, gives $\mu_k\ge\beta_k$ for $k\le r$. For
$v=b+p$ ($b\in\operatorname{ran}B$, $p\in\operatorname{ran}P$) and $t>\|A\|$,
$\ip v{H_tv}\le\ip b{Ab}+2\|A\||b||p|-(t-\|A\|)|p|^2\le\ip b{(A+\varepsilon_t)b}$ with
$\varepsilon_t=\|A\|^2/(t-\|A\|)$. The form $v\mapsto\ip{Bv}{(A+\varepsilon_t)Bv}$ has the eigenvalues
$\beta_k+\varepsilon_t>0$ and $0$, so $\mu_k\le\beta_k+\varepsilon_t$ for $k\le r$. Finally $\mu_{r+1}\le\|A\|-t$ by the min-max
principle with the subspace $\operatorname{ran}P$ of codimension $r$. Hence $\Tr e^{H_t}\to\sum_ke^{\beta_k}$. The second
statement is the first one with $B$ and $P$ exchanged, since $e^t\Tr e^{A-tP}=\Tr e^{A-|t|B}$ for $t<0$.
\end{proof}

\subsection{Proof of Theorem~\ref{thm:D}}\label{subsec:proofD}
The properties of $\rho$ are Lemma~\ref{lem:pencil}. We prove~\eqref{eq:2bmv} in four steps.

\emph{Step 0: $\cD/a^2$ is entire.} $\cD$ is entire on $\C^2$. Since $e^{-tP}=B+e^{-t}P$,
$\cD(0,t)=\Tr e^{-tP}-(M-r)e^{-t}-r=0$ and
$\partial_a\cD(0,t)=\Tr(ge^{-tP})-e^{-t}\Tr(PgP)-\Tr(BgB)=0$. Hence
$\cG(a,t):=\cD(a,t)/a^2=\int_0^1(1-u)\,\partial_a^2\cD(ua,t)\dd u$ is entire on $\C^2$.

\emph{Step 1: a distribution with support in $K$.} For $\zeta=(\zeta_1,\zeta_2)\in\C^2$ put
$\hat\cG(\zeta)=\cG(-i\zeta_1,i\zeta_2)$, so that $e^{-i(\zeta_1s+\zeta_2\tau)}=e^{as-t\tau}$ when $(a,t)=(-i\zeta_1,i\zeta_2)$. Let
$\eta=\Im\zeta$ and $H_K(\eta)=\max_{(s,\tau)\in K}(s\eta_1+\tau\eta_2)$, the supporting function of $K$. For $X=ag-tP$ the
Hermitian part of $X$ is $\eta_1g+\eta_2P$, because $\Re a=\eta_1$ and $\Re t=-\eta_2$. If $u(\theta)=e^{\theta X}v$ then
$\frac{\dd}{\dd\theta}|u|^2=2\ip u{(\eta_1g+\eta_2P)u}\le2\lambda_{\max}(\eta_1g+\eta_2P)|u|^2$, so
$|\Tr e^X|\le Me^{\lambda_{\max}(\eta_1g+\eta_2P)}\le Me^{H_K(\eta)}$, the last step because
$\lambda_{\max}(\eta_1g+\eta_2P)=\max_{|v|=1}(\eta_1\ip v{gv}+\eta_2\ip v{Pv})$ and $(\ip v{gv},\ip v{Pv})\in K$. Similarly
$|e^{-t}\Tr_Pe^{aPgP}|\le(M-r)e^{H_K(\eta)}$ and $|\Tr_Be^{aBgB}|\le re^{H_K(\eta)}$, since the eigenvalues of the
compressions lie in $[\lmin,\lmax]$ and $(\mu,1),(\beta,0)\in K$. Hence $|\cD|\le2Me^{H_K(\eta)}$, and
$|\hat\cG(\zeta)|\le2Me^{H_K(\eta)}$ for $|\zeta_1|\ge1$. For $|\zeta_1|<1$ the maximum principle for
$z\mapsto\hat\cG(z,\zeta_2)$ on the disc $|z|\le2$ gives
$|\hat\cG(\zeta)|\le\frac M2\max_{|z|=2}e^{H_K(\Im z,\eta_2)}\le\frac M2e^{H_K(\eta)+3c_K}$, where
$c_K=\max(|\lmin|,|\lmax|)$ is a Lipschitz constant of $H_K$ in $\eta_1$ and $|\Im z-\eta_1|\le3$. So
$|\hat\cG(\zeta)|\le Ce^{H_K(\Im\zeta)}$ on $\C^2$. By the Paley--Wiener--Schwartz theorem~\cite[Theorem~7.3.1]{Hormander}
there is a distribution $\Theta$ with support in the convex compact set $K$ whose Fourier--Laplace transform
$\Theta(e^{-i\ip\cdot\zeta})$ is $\hat\cG(\zeta)$; equivalently,
\begin{equation}\label{eq:Theta}
\Theta\bigl(e^{as-t\tau}\bigr)=\cG(a,t)=\frac{\cD(a,t)}{a^2}\qquad\text{for all }(a,t)\in\C^2 .
\end{equation}
On $\R^2$, $|\hat\cG|\le C$, since $H_K(0)=0$.

\emph{Step 2: $\Theta=\rho$ on the open strip $\R\times(0,1)$.} Let $\varphi\in C_c^\infty(\R^2)$ and
$\hat\varphi(k)=\int\varphi(z)e^{-ik\cdot z}\dd z$. The Fourier transform of $\Theta$ is the bounded function $\hat\cG$
on $\R^2$~\cite[Theorem~7.1.14]{Hormander}, so
\begin{equation}\label{eq:Tphi}
\Theta(\varphi)=(2\pi)^{-2}\int_{\R^2}\hat\cG(k)\hat\varphi(-k)\dd k,
\end{equation}
with an integrable integrand. The map $(\sigma,\xi)\mapsto k=(\sigma,-\sigma\xi)$ is a diffeomorphism of
$(\R\setminus\{0\})\times\R$ onto the complement of the null set $\{k_1=0\}$, with Jacobian $|\sigma|$; by Fubini,
\[
\Theta(\varphi)=(2\pi)^{-2}\int_\R\dd\xi\int_\R\dd\sigma\,|\sigma|\,\hat\cG(\sigma,-\sigma\xi)\,\hat\varphi(-\sigma,\sigma\xi).
\]
Fix $\xi$. (i) By Lemma~\ref{lem:slices}(2) with $a=-i\sigma$,
$\hat\cG(\sigma,-\sigma\xi)=\cD(a,a\xi)/a^2=\hat U_\xi(\sigma)$. (ii)
$\hat\varphi(-\sigma,\sigma\xi)=\iint\varphi(s,\tau)e^{i\sigma(s-\xi\tau)}\dd s\dd\tau=\hat\Psi_\xi(-\sigma)$, where
$\Psi_\xi(w)=\int_\R\varphi(w+\xi\tau,\tau)\dd\tau$ belongs to $C_c^\infty(\R)$. (iii) By Lemma~\ref{lem:slices}(1),
$i\sigma\hat U_\xi=\hat n_\xi$, so $|\sigma|\hat U_\xi(\sigma)=-i\operatorname{sgn}(\sigma)\hat n_\xi(\sigma)=(\cH n_\xi)^\wedge(\sigma)$,
with $\cH n_\xi\in L^2(\R)$. By Parseval's identity
$\int fh=(2\pi)^{-1}\int\hat f(\sigma)\hat h(-\sigma)\dd\sigma$ ($f\in L^2$, $h$ Schwartz),
\[
\frac1{2\pi}\int_\R|\sigma|\hat U_\xi(\sigma)\hat\Psi_\xi(-\sigma)\dd\sigma=\int_\R(\cH n_\xi)(w)\Psi_\xi(w)\dd w
=\iint(\cH n_\xi)(s-\xi\tau)\varphi(s,\tau)\dd s\dd\tau .
\]
(iv) Since $g-\xi P-(s-\xi\tau)=g-s-\xi(P-\tau)$, Lemma~\ref{lem:slices}(3) gives
$(\cH n_\xi)(s-\xi\tau)=\frac1\pi\log|p_{s,\tau}(\xi)/q_{s,\tau}(\xi)|$ for $\tau\notin\{0,1\}$ and almost every $s$. Therefore
\begin{equation}\label{eq:backproj}
\Theta(\varphi)=\frac1{2\pi^2}\int_\R\dd\xi\iint\varphi(s,\tau)\log\Bigl|\frac{p_{s,\tau}(\xi)}{q_{s,\tau}(\xi)}\Bigr|\dd s\dd\tau .
\end{equation}
Now assume $\operatorname{supp}\varphi\subseteq[-S_0,S_0]\times[\epsilon,1-\epsilon]$ with $0<\epsilon<\frac12$. For $(s,\tau)$ in this set,
Lemma~\ref{lem:pencil}(2) puts all roots of $p_{s,\tau}$ and $q_{s,\tau}$ in the disc $|\xi|\le R:=(\|g\|+S_0)/\epsilon$, and by
Lemma~\ref{lem:pencil}(1) the two polynomials have the same leading coefficient and the same sum of roots. For
$|\xi|\le2R$, $|\log|p/q||\le\sum_i(|\log|\xi-\xi_i||+|\log|\xi-y_i||)$, whose integral over $[-2R,2R]$ is bounded
uniformly. For $|\xi|>2R$, using $|\log|1-z|+\Re z|\le|z|^2$ for $|z|\le\frac12$ and $\Re\sum_i(\xi_i-y_i)=0$,
$|\log|p/q|(\xi)|=\bigl|\sum_i(\log|1-\xi_i/\xi|-\log|1-y_i/\xi|)\bigr|\le2MR^2/\xi^2$. So the integrand
of~\eqref{eq:backproj} is absolutely integrable on $\R\times\operatorname{supp}\varphi$, and Fubini and
Lemma~\ref{lem:jensen} (applicable by Lemma~\ref{lem:pencil}(1),(5)) give
\[
\Theta(\varphi)=\iint\varphi(s,\tau)\Bigl[\frac1{2\pi^2}\int_\R\log\Bigl|\frac{p_{s,\tau}(\xi)}{q_{s,\tau}(\xi)}\Bigr|\dd\xi\Bigr]\dd s\dd\tau
=\iint\varphi(s,\tau)\rho(s,\tau)\dd s\dd\tau .
\]

\emph{Step 3: no singular part on the edges.} By Lemma~\ref{lem:pencil}, $\rho$ defines a distribution with support in
$K$, and $\Theta-\rho$ vanishes on $\R\times(0,1)$ and outside $K$, so its support lies in the union of the disjoint
segments $L_0=[\lmin,\lmax]\times\{0\}$ and $L_1=[\lmin,\lmax]\times\{1\}$. Splitting $\Theta-\rho$ with smooth cutoffs
and applying the structure theorem for distributions supported in a hyperplane~\cite[Theorem~2.3.5]{Hormander} to
each piece, there are $N_0\ge0$ and compactly supported distributions $\alpha_k,\beta_k$ on $\R$ with
\[
\Theta-\rho=\sum_{k=0}^{N_0}\bigl[\alpha_k\otimes\delta^{(k)}(\tau)+\beta_k\otimes\delta^{(k)}(\tau-1)\bigr].
\]
Evaluating on $e^{as-t\tau}$, and using $\ip{\delta^{(k)}}f=(-1)^kf^{(k)}(0)$,
\[
\cG(a,t)-\iint e^{as-t\tau}\rho(s,\tau)\dd s\dd\tau=\sum_{k=0}^{N_0}t^k\bigl[A_k(a)+e^{-t}B_k(a)\bigr],
\]
where $A_k(a)=\alpha_k(e^{as})$ and $B_k(a)=\beta_k(e^{as})$. Fix a real $a\ne0$. As $t\to+\infty$: $\cG(a,t)\to0$ by Lemma~\ref{lem:limits} (with $A=ag$), the integral tends to $0$ by
dominated convergence ($\rho\in L^1$ with compact support, $e^{-t\tau}\le1$, and $e^{-t\tau}\to0$ for $\tau>0$), and
$t^ke^{-t}B_k(a)\to0$; so the polynomial $\sum_kA_k(a)t^k$ tends to $0$, and every $A_k(a)=0$. As $t\to-\infty$:
$e^t\cG(a,t)=a^{-2}[e^t\Tr e^{ag-tP}-\Tr_Pe^{aPgP}-e^t\Tr_Be^{aBgB}]\to0$ by Lemma~\ref{lem:limits}, and
$e^t\iint e^{as-t\tau}\rho=\iint e^{as+t(1-\tau)}\rho\to0$; so $\sum_kB_k(a)t^k\to0$ and every $B_k(a)=0$. The functions
$A_k,B_k$ are entire (Fourier--Laplace transforms of compactly supported distributions) and vanish on
$\R\setminus\{0\}$, hence identically, so $\alpha_k=\beta_k=0$. Thus $\Theta=\rho$, and~\eqref{eq:Theta}
is~\eqref{eq:2bmv}.

\emph{Step 4: the $\tau$-marginals.} By Duhamel's formula,
\[
\partial_a^2\Tr e^{ag-tP}\big|_{a=0}=\int_0^1\Tr\bigl(ge^{-\theta tP}ge^{-(1-\theta)tP}\bigr)\dd\theta ,
\]
and with $e^{-\theta tP}=B+e^{-\theta t}P$ this equals $\Tr(BgB)^2+e^{-t}\Tr(PgP)^2+2\|BgP\|_F^2(1-e^{-t})/t$. The first two
terms are removed by the subtractions in $\cD$, so
\[
\cG(0,t)=\tfrac12\partial_a^2\cD(0,t)=\|BgP\|_F^2\,\frac{1-e^{-t}}t=\|BgP\|_F^2\int_0^1e^{-t\tau}\dd\tau .
\]
By~\eqref{eq:Theta} at $a=0$, where now $\Theta=\rho$, and uniqueness of Laplace transforms,
$\bigl(\int\rho(s,\tau)\dd s\bigr)\dd\tau=\|BgP\|_F^2\dd\tau$ on $[0,1]$; the
marginal is continuous in $\tau\in(0,1)$ (dominated convergence, Lemma~\ref{lem:pencil}), so the identity holds for every
$\tau$. This completes the proof of Theorem~\ref{thm:D}.\qed

\subsection{Proof of Theorem~\ref{thm:C}}\label{subsec:proofC}
Let $\nu_k=x_k+iy_k$ ($k=1,\dots,M$) be the eigenvalues of $X=B+ig$. Since $\Re\ip v{Xv}=\ip v{Bv}\in[0,1]$ for unit $v$,
$x_k\in[0,1]$. Let $\omega=\sum_k\omega_{\nu_k}$ (Lemma~\ref{lem:h}(5)) and $\sigma=\sum_\mu\delta_\mu$, the sum over the
eigenvalues $\mu$ of $PgP|_{\operatorname{ran}P}$. By Lemma~\ref{lem:h}(5), the left side of~\eqref{eq:star} is
$\int(w-\lambda)_+\dd(\omega-\sigma)(w)$; both measures are finite, and $\omega$ is the sum of finitely many atoms and
a measure with a density that is $O(e^{-\pi|w|})$ as $|w|\to\infty$ (Lemma~\ref{lem:poisson}(1)).

\emph{(a) Laplace transforms.} For real $a$ with $0<|a|<\pi$, Lemma~\ref{lem:h}(5) gives
\[
\hat\omega(a):=\int e^{aw}\omega(\dd w)=\sum_ke^{ay_k}\frac{\sin(a(1-x_k))}{\sin a}=\frac{\Tr e^{ag+iaP}-\Tr e^{ag-iaP}}{2i\sin a},
\]
because $e^{-ia(X-\one)}=e^{ag+iaP}$ has the eigenvalues $e^{ay_k}e^{ia(1-x_k)}$ and $e^{ia(X^*-\one)}=e^{ag-iaP}$ has the
eigenvalues $e^{ay_k}e^{-ia(1-x_k)}$. Also $\hat\sigma(a)=\Tr_Pe^{aPgP}$. In $\cD(a,-ia)-\cD(a,ia)$ the $B$-terms cancel,
and
\[
\cD(a,-ia)-\cD(a,ia)=\Tr e^{ag+iaP}-\Tr e^{ag-iaP}-(e^{ia}-e^{-ia})\Tr_Pe^{aPgP}=2i\sin a\,\bigl[\hat\omega(a)-\hat\sigma(a)\bigr].
\]
By Theorem~\ref{thm:D} at $t=\mp ia$, the left side is $2ia^2\iint e^{as}\sin(a\tau)\rho(s,\tau)\dd s\dd\tau$. With
Lemma~\ref{lem:poisson}(3) for $x=1-\tau$, $\sin(a\tau)/\sin a=\int e^{-au}K_{1-\tau}(u)\dd u$, and Tonelli's theorem (all
integrands are nonnegative for real $a$),
\begin{equation}\label{eq:laplaceV}
\hat\omega(a)-\hat\sigma(a)=a^2\int_\R e^{aw}V(w)\dd w,\qquad V(w):=\int_0^1\!\!\int_\R K_{1-\tau}(s-w)\rho(s,\tau)\dd s\dd\tau .
\end{equation}
$V$ is finite, continuous and $O(e^{-\pi|w|})$: $\int K_{1-\tau}=\tau\le1$, $\sup_s\rho(s,\tau)\le C(\tau(1-\tau))^{-1/2}$,
$\rho$ has compact support, and $K_{1-\tau}(u)\le2\sin(\pi\tau)e^{-\pi|u|}$ for $|u|\ge1$.

\emph{(b) Moments.} By~\eqref{eq:laplaceV}, $\hat\omega-\hat\sigma=O(a^2)$ as $a\to0$ (both sides extend holomorphically to
$|\Re a|<\pi$), so $\omega$ and $\sigma$ have the same mass and the same first moment. Put
$W(w)=\int(w-u)_+\dd(\omega-\sigma)(u)$. Then $W$ is continuous; for $w$ below all atoms of $\omega$ and $\sigma$,
$W(w)=\int_{u<w}(w-u)\dd\omega(u)=O(e^{-\pi|w|})$, and for $w$ above all atoms, by the moment identities,
$W(w)=\int_{u>w}(u-w)\dd\omega(u)=O(e^{-\pi w})$; the same holds for $W'$. As $W''=\omega-\sigma$ in the sense of distributions, two integrations by
parts give $\int e^{aw}W(w)\dd w=a^{-2}(\hat\omega(a)-\hat\sigma(a))=\int e^{aw}V(w)\dd w$ for real $0<|a|<\pi$. Both sides are
holomorphic in $|\Re a|<\pi$, so they agree on $i\R$: the integrable continuous functions $W$ and $V$ have the same
Fourier transform, and $W=V$.

\emph{(c) Conclusion.} Since $(u-\lambda)_+=(\lambda-u)_++(u-\lambda)$ and $\omega-\sigma$ has zero mass and zero first moment,
$\int(u-\lambda)_+\dd(\omega-\sigma)=W(\lambda)=V(\lambda)$, which is~\eqref{eq:star}.

\emph{(d) Equality.} $K_{1-\tau}(s-\lambda)>0$ for $0<\tau<1$, and $\rho\ge0$ is continuous on $\R\times(0,1)$; so $V(\lambda)=0$ for
one $\lambda$ forces $\rho=0$, hence $\cD\equiv0$ by Theorem~\ref{thm:D}, hence $\|BgP\|_F^2=0$ by Step~4 of
Section~\ref{subsec:proofD}, that is, $BgP=0=PgB$ and $[B,g]=BgP-PgB=0$. Conversely, if $[B,g]=0$ then $g=g_d$,
$p_{s,\tau}=q_{s,\tau}$, $\rho=0$ and $V\equiv0$.\qed

\subsection{Remarks}

\begin{remark}[$M=2$]\label{rem:M2}
Let $B=\mathrm{diag}(1,0)$ and $g=\begin{pmatrix}\alpha&c\\\bar c&\beta\end{pmatrix}$. Then $p_{s,\tau}(\xi)$ has discriminant
$(m(\tau)-s)^2-4\tau(1-\tau)|c|^2$ with $m(\tau)=(1-\tau)\alpha+\tau\beta$, so
\[
\rho(s,\tau)=\frac{\bigl(4\tau(1-\tau)|c|^2-(s-m(\tau))^2\bigr)_+^{1/2}}{2\pi\tau(1-\tau)} :
\]
each slice $\tau=\mathrm{const}$ is a semicircle law of radius $2|c|\sqrt{\tau(1-\tau)}$ and mass $|c|^2=\|BgP\|_F^2$,
centred on the segment from $\alpha$ (the eigenvalue of $BgB$, at $\tau=0$) to $\beta$ (that of $PgP$, at $\tau=1$).
\end{remark}

\begin{remark}[convex order]\label{rem:convex}
In the proof of Theorem~\ref{thm:C}, $\omega$ (the sweep of the eigenvalues of $B+ig$ onto the left edge of $\bar S$ by
harmonic measure) and $\sigma$ (the eigenvalue measure of $PgP|_{\operatorname{ran}P}$) have equal masses $M-r$ and equal
means, and $\int(\lambda-u)_+\dd(\omega-\sigma)\ge0$ for all $\lambda$. Hence Theorem~\ref{thm:C} says exactly that
$\sigma$ is dominated by $\omega$ in the convex order: $\Tr_Pf(PgP)\le\int f\dd\omega$ for every convex $f$ of at most
linear growth, with equality for all such $f$ if and only if $[B,g]=0$. Applying this to $(P,g)$, whose matrix
$P+ig=\one-(B+ig)^*$ has the eigenvalues $1-\bar\nu$, gives the mirror statement: the eigenvalue measure of
$BgB|_{\operatorname{ran}B}$ is dominated in convex order by the sweep of $\spec(B+ig)$ onto the right edge. When $B$
and $g$ commute both statements are identities. In this sense Theorem~\ref{thm:C} is a Schur--Horn type majorisation
for the non-normal matrix $B+ig$; first-order (distribution-function) versions are false whenever $[B,g]\ne0$,
because $\omega$ has a density as soon as one eigenvalue lies in the open strip.
\end{remark}

\begin{remark}[relation to Hein\"avaara's tracial joint spectral measures]\label{rem:heinavaara}
For Hermitian $A,B$, Hein\"avaara~\cite{He} constructs a positive measure $\mu_{A,B}$ on $\R^2$ such that
$\tr\,\mathcal{T}(f)(xA+yB)=\int f(ax+by)\dd\mu_{A,B}(a,b)$ for a large class of $f$, where
$\mathcal T(f)(x)=\int_0^1\frac{1-t}tf(xt)\dd t$ (so $\mathcal T$ maps $x^k$ to $x^k/(k(k+1))$), computes the absolutely
continuous part of $\mu_{A,B}$ explicitly as $\frac1{2\pi}\sum_i|\Im\lambda_i(C_{a,b})|$ for an auxiliary pencil
$C_{a,b}$ (\cite[Theorem~3.1]{He}), and deduces Stahl's theorem. That proof computes a Fourier transform in polar coordinates,
one line through the origin at a time, and uses an integral identity for $\log|\det|$ of a pencil. Our proof of
Theorem~\ref{thm:D} follows the same route in an affine chart adapted to the projection $P$: the lines are the real
slices of Lemma~\ref{lem:slices}, the division by $a^2$ plays the role of $\mathcal T$, the pinched subtractions remove
the contributions of the edges $\tau\in\{0,1\}$, and Lemma~\ref{lem:jensen} is the corresponding $\log|\det|$ identity. We
use only the statement of Theorem~\ref{thm:D} in the form given above, and we have given a complete proof of it.
Appendix~\ref{app:RI} gives a second proof, in which the Fourier inversion along lines is replaced by a Radon identity
for the pencil (Theorem~\ref{thm:RI}), proved with a complex Burgers equation for its roots.
\end{remark}

\begin{remark}[an independent proof for rank one, and an automorphism of order three]\label{rem:rankone}
For $\operatorname{rank}B\le1$ there is a second proof of Theorem~\ref{thm:C} (inequality part), by a different method;
we sketch it, the details being in the accompanying notes~\cite{repo-oqp27}. Let $B=vv^*$, and let
$\nu_k=x_k+iy_k$ be the eigenvalues of $B+ig$, so that $x_k\in[0,1]$ and $\sum_kx_k=1$. If all $x_k>0$, the matrix
determinant lemma gives $\det(g-iB-t)/\det(g+iB-t)=\exp(-2\pi iF(t))$ for real $t$, where
$F(t)=\sum_k[\frac12+\frac1\pi\arctan\frac{t-y_k}{x_k}]$ is the distribution function of the sum of the Cauchy laws with
centres $y_k$ and scales $x_k$, and the left side equals $(1-iG(t))/(1+iG(t))$ with $G(t)=\ip v{(g-t)^{-1}v}$; the zeros of
$G$ are the eigenvalues of $PgP|_{\operatorname{ran}P}$. Hence these eigenvalues are the integer quantiles
$\gamma_j=\min\{t:F(t)\ge j\}$, $1\le j\le M-1$, and conversely every configuration with $x_k>0$ and $\sum x_k=1$ arises.
Theorem~\ref{thm:C} for rank one thus becomes the inequality $\sum_kh_0(\nu_k)\ge\sum_j(\gamma_j)_+$ for $M$ points in the
strip. It is proved by induction on $M$, for the more general constraint $\sum_kx_k=s\le1$ with a correction term: a
penalised minimiser is either reducible (a point on or near $\partial S$, handled by the induction hypothesis) or an
interior critical point; at an interior critical point all $\nu_k$ are zeros, counted with multiplicity, of a holomorphic
function on $S$ of the form $\phi_0'+\frac i\pi\sum_j\frac{c_j}{z-i\gamma_j}-m_0+iaz+\varepsilon\pi\cos\pi z$ ($c_j,\varepsilon>0$, $a\ge0$),
and an argument-principle count shows that such a function has at most $J+1\le M-1$ zeros in $S$, where $J$ is the
number of positive quantiles. Combined with a duality argument it also gives Theorem~\ref{thm:C} whenever
$\operatorname{rank}P\le3$, hence for all $M\le5$.

The derivative $\phi_0'(z)=-\frac1\pi\log(1-e^{-i\pi z})$ has a curious structure. Let $\chi=i\phi_0'$. In the coordinate
$w=e^{-i\pi z}$, which maps $S$ onto the lower half-plane, $\chi$ acts as the M\"obius map $w\mapsto1/(1-w)$; hence
$\chi$ is an automorphism of $S$ of order three, with fixed point $\frac13$, and $z+\chi(z)+\chi(\chi(z))=1$ on $S$ (the
three values of $w$ have product $-1$ and arguments in $(-\pi,0)$). By the holomorphic functional calculus,
$X+\chi(X)+\chi(\chi(X))=\one$ for every matrix $X$ with spectrum in $S$. If $B+ig$ has no eigenvalue on $\partial S$, the
derivative of $g\mapsto\sum_{\nu\in\spec(B+ig)}h_0(\nu)$ in a Hermitian direction $\delta g$ is $\Tr[\Re\chi(B+ig)\,\delta g]$,
where $\Re Y=\frac12(Y+Y^*)$; so the critical points of the defect in~\eqref{eq:star} (at $\lambda=0$, where
$PgP|_{\operatorname{ran}P}$ is invertible) are the solutions of $\Re\chi(B+ig)=Q$, with $Q$ the spectral projection of
$PgP$ for $(0,\infty)$.
\end{remark}

%======================================================================
\section{A noncommutative continuum theorem}\label{sec:continuum}
%======================================================================

Let $\T=\R/2\pi\Z$ with $\dd m=\dd\theta/2\pi$, and let $\Disc$ be the open unit disc. A \emph{projection step field} on
$\C^M$ is given by points $t_0<t_1<\dots<t_{n-1}<t_0+2\pi$ (indices modulo $n$, $t_n=t_0+2\pi$), the open arcs
$J_j=(t_j,t_{j+1})$, and orthogonal projections $B_j$ on $\C^M$; it is the function $B(\theta)=B_j$ for $\theta\in J_j$
(values at the points $t_j$ play no role). Its \emph{conjugate function} is
\begin{equation}\label{eq:conj}
g(\theta)=\pv\!\int_\T\cot\frac{\theta-\varphi}2\,B(\varphi)\dd m(\varphi)
=\frac1\pi\sum_{j=0}^{n-1}(B_j-B_{j-1})\log\Bigl|\sin\frac{\theta-t_j}2\Bigr|\qquad(\theta\notin\{t_j\}),
\end{equation}
where the second expression follows from $\frac{\dd}{\dd\varphi}\log|\sin\frac{\theta-\varphi}2|=-\frac12\cot\frac{\theta-\varphi}2$
and summation by parts. Thus $g$ is Hermitian, real-analytic on each arc $J_j$, and in $L^p(\T)$ for every $p<\infty$.
We write $P(\theta)=\one-B(\theta)$ and $\bar B=\int_\T B\dd m$.

\begin{lemma}[Herglotz function of a step field]\label{lem:herglotz}
For $z\in\Disc$ let $\cF(z)=\int_\T\frac{e^{i\varphi}+z}{e^{i\varphi}-z}B(\varphi)\dd m(\varphi)$. Then:
\begin{enumerate}
\item $\cF(z)=\bar B+\frac i\pi\sum_{j}(B_j-B_{j-1})\operatorname{Log}(1-ze^{-it_j})$ (principal logarithm);
\item $\cF$ is holomorphic on $\Disc$ and $\cF(0)=\bar B$; its Hermitian part
$\Re\cF(z):=\frac12(\cF(z)+\cF(z)^*)$ is $\int_\T P_z(\varphi)B(\varphi)\dd m(\varphi)$, where $P_z>0$ is the Poisson kernel;
\item $\cF$ extends continuously to $\bar\Disc\setminus\{e^{it_j}\}$, with $\cF(e^{i\theta})=B(\theta)+ig(\theta)$ for
$\theta\notin\{t_j\}$;
\item for $\frac12\le r<1$, $\|\cF(re^{i\theta})\|\le C_0+\frac1\pi\sum_j|\log|\sin\frac{\theta-t_j}2||$, with $C_0=1+n(\log2+\frac\pi2)/\pi$.
\end{enumerate}
\end{lemma}

\begin{proof}
(1) Since $\frac{e^{i\varphi}+z}{e^{i\varphi}-z}=1+\frac2i\frac{\dd}{\dd\varphi}\operatorname{Log}(1-ze^{-i\varphi})$ and $\Re(1-ze^{-i\varphi})>0$,
$\frac1{2\pi}\int_a^b\frac{e^{i\varphi}+z}{e^{i\varphi}-z}\dd\varphi=\frac{b-a}{2\pi}-\frac i\pi[\operatorname{Log}(1-ze^{-ib})-\operatorname{Log}(1-ze^{-ia})]$.
Multiply by $B_j$ for the arc $J_j$, sum over $j$, and regroup the logarithms by endpoints. (2) is clear. (3) The
right side of (1) is continuous on $\bar\Disc\setminus\{e^{it_j}\}$. For $u\in(0,2\pi)$,
$1-e^{iu}=2\sin\frac u2\,e^{i(u-\pi)/2}$, so $\operatorname{Log}(1-e^{iu})=\log|2\sin\frac u2|+\frac i2(u-\pi)$. The
imaginary part of $\cF(e^{i\theta})$ is therefore $\frac1\pi\sum_j(B_j-B_{j-1})\log|2\sin\frac{\theta-t_j}2|=g(\theta)$ (the
$\log2$ terms cancel since $\sum_j(B_j-B_{j-1})=0$). For the real part let $\theta\in J_x$ and $u_j=(\theta-t_j)\bmod2\pi\in(0,2\pi)$,
so $u_j=\theta-t_j$ for $j\le x$ and $u_j=\theta-t_j+2\pi$ for $j>x$ (with $0\le j\le n-1$). Then
$\Re\cF(e^{i\theta})=\bar B-\frac1{2\pi}\sum_j(B_j-B_{j-1})u_j$, and summation by parts gives
$\sum_j(B_j-B_{j-1})u_j=-\sum_j(B_j-B_{j-1})t_j+2\pi\sum_{j>x}(B_j-B_{j-1})=2\pi(\bar B-B_{n-1})+2\pi(B_{n-1}-B_x)$, so
$\Re\cF(e^{i\theta})=B_x$. (4) $\|B_j-B_{j-1}\|\le1$, and for $\frac12\le r<1$,
$\sqrt2|\sin\frac u2|\le|1-re^{iu}|\le2$, hence $|\operatorname{Log}(1-re^{iu})|\le\log2+\frac\pi2+|\log|\sin\frac u2||$.
\end{proof}

\begin{theorem}[continuum theorem]\label{thm:continuum}
Let $B$ be a projection step field on $\C^M$ with conjugate function $g$, and let $A$ be a measurable field of
Hermitian matrices with $0\le A(\theta)\le P(\theta)$. Put $\alpha=\int\tau(A)\dd m$, $\beta=\int\tau(B)\dd m$, and let
$\beta_1,\dots,\beta_M\in[0,1]$ be the eigenvalues of $\bar B$. For $\lambda\in\R$ and $\theta\notin\{t_j\}$ put
\begin{align*}
D_\lambda(\theta)&=\sum_{\nu\in\spec(B(\theta)+ig(\theta))}h_\lambda(\nu)-\Tr\bigl[(P(\theta)(g(\theta)-\lambda)P(\theta))_+\bigr],\\
S_\lambda(\theta)&=\tau\bigl[(P(\theta)(g(\theta)-\lambda)P(\theta))_+\bigr]-\tau\bigl[A(\theta)(g(\theta)-\lambda)\bigr].
\end{align*}
Then $D_\lambda\ge0$, $S_\lambda\ge0$, both are integrable, and
\begin{equation}\label{eq:deficit}
\lambda\alpha+h_\lambda(\beta)-\int_\T\tau(Ag)\dd m=\Bigl[h_\lambda(\beta)-\frac1M\sum_{i=1}^Mh_\lambda(\beta_i)\Bigr]
+\frac1M\int_\T D_\lambda\dd m+\int_\T S_\lambda\dd m ,
\end{equation}
where the bracket is also nonnegative. In particular $\int\tau(Ag)\dd m\le\lambda\alpha+h_\lambda(\beta)$ for every $\lambda$, and if
$0<\alpha,\beta$ and $\alpha+\beta<1$, then $\int\tau(Ag)\dd m\le\Phic(\alpha,\beta)$.
\end{theorem}

\begin{proof}
\emph{Signs.} $D_\lambda\ge0$ is Theorem~\ref{thm:C} for the pair $(B(\theta),g(\theta))$ (note that
$\Tr[(P(g-\lambda)P)_+]$ is the same whether $P(g-\lambda)P$ is regarded on $\C^M$ or on $\operatorname{ran}P$). For
$S_\lambda$: $0\le A\le P$ gives $BAB=0$, hence $AB=0$ and $A=PAP$; with $Y=P(g-\lambda)P$,
$\tau(A(g-\lambda))=\tau(AY)\le\tau(AY_+)\le\tau(PY_+)=\tau(Y_+)$. The bracket is nonnegative because
$x\mapsto h_\lambda(x)$ is concave on $[0,1]$ (Lemma~\ref{lem:h}(4)) and $\frac1M\sum_i\beta_i=\tau(\bar B)=\beta$.

\emph{A harmonic function.} Let $K_0=\{v:B_jv=0\ \forall j\}$ and $K_1=\{v:B_jv=v\ \forall j\}$; these subspaces and
$K'=(K_0\oplus K_1)^\perp$ reduce every $B_j$, and $\cF(z)=0\oplus\one\oplus\cF'(z)$ accordingly. For a unit vector
$v\in K'$ and $z\in\Disc$, $\ip v{\Re\cF(z)v}=\sum_j\omega_j(z)\ip v{B_jv}$ with harmonic measures $\omega_j(z)>0$,
$\sum_j\omega_j(z)=1$; this is $0$ only if $v\in K_0$ and $1$ only if $v\in K_1$, which is impossible. By compactness the
numerical range of $\Re\cF'(z)$ is a compact subset of $(0,1)$, so $\spec\cF'(z)\subseteq S$. By the holomorphic
functional calculus (a Riesz--Dunford integral over a contour in $S$, locally constant in $z$),
$z\mapsto\Re\Tr\phi_\lambda(\cF'(z))$ is harmonic on $\Disc$. Hence
\[
u_\lambda(z):=\sum_{\nu\in\spec\cF(z)}h_\lambda(\nu)=\dim K_0\,h_\lambda(0)+\dim K_1\,h_\lambda(1)+\Re\Tr\phi_\lambda(\cF'(z))
\]
is harmonic on $\Disc$, and $u_\lambda(0)=\sum_ih_\lambda(\beta_i)$ because $\cF(0)=\bar B$.

\emph{Boundary passage.} Every eigenvalue $\nu$ of $\cF(z)$, $z\in\bar\Disc\setminus\{e^{it_j}\}$, lies in $\bar S$ and
satisfies $|\Im\nu|\le\|\cF(z)\|$; by Lemma~\ref{lem:h}(3), $0\le u_\lambda(re^{i\theta})\le M(\|\cF(re^{i\theta})\|+|\lambda|+2)$,
which by Lemma~\ref{lem:herglotz}(4) is bounded, for $\frac12\le r<1$, by an integrable function of $\theta$. For
$\theta\notin\{t_j\}$, $\cF(re^{i\theta})\to B(\theta)+ig(\theta)$ as $r\to1$ (Lemma~\ref{lem:herglotz}(3)), the eigenvalues
converge, and $h_\lambda$ is continuous on $\bar S$; so $u_\lambda(re^{i\theta})\to u^*_\lambda(\theta):=\sum_{\nu\in\spec(B(\theta)+ig(\theta))}h_\lambda(\nu)$.
The mean value property on $|z|=r$ and dominated convergence give $\int_\T u^*_\lambda\dd m=\sum_ih_\lambda(\beta_i)$.

\emph{Conclusion.} By definition, $\tau(A(g-\lambda))=\frac1Mu^*_\lambda-\frac1MD_\lambda-S_\lambda$ off $\{t_j\}$, and every term is
integrable ($g\in L^p$, $A$ bounded, $u^*_\lambda$ dominated as above). Integrating,
$\int\tau(Ag)\dd m-\lambda\alpha=\frac1M\sum_ih_\lambda(\beta_i)-\frac1M\int D_\lambda-\int S_\lambda$, which
is~\eqref{eq:deficit}. The last assertion follows from~\eqref{eq:deficit} and Lemma~\ref{lem:classical} below, whose
proof uses only the identity~\eqref{eq:deficit} for $M=1$.
\end{proof}

\begin{lemma}[two adjacent arcs]\label{lem:twoarcs}
Let $\alpha,\beta\ge0$ with $\alpha+\beta\le1$, and let $g_\beta$ be the conjugate function of the scalar step field
$\one_{(-2\pi\beta,0)}$. Then $\int_0^{2\pi\alpha}g_\beta\dd m=\Phic(\alpha,\beta)$. Equivalently, for $a,b\ge0$ with $a+b\le N$,
$\int_0^a\int_{-b}^0\cot\frac{\pi(\xi-\eta)}N\dd\eta\dd\xi=N^2\Phic(a/N,b/N)$.
\end{lemma}

\begin{proof}
The function $\Xi(u)=-\frac{N^2}{2\pi^2}\Cl(2\pi u/N)$ is continuous, odd and $N$-periodic, with $\Xi''(u)=\cot\frac{\pi u}N$
for $u\notin N\Z$ (since $\Cl''(\theta)=-\frac12\cot\frac\theta2$) and $\Xi'$ locally integrable. Integrating twice,
$\int_0^a\int_{-b}^0\Xi''(\xi-\eta)\dd\eta\dd\xi=\Xi(a+b)-\Xi(a)-\Xi(b)+\Xi(0)$, which equals
$\frac{N^2}{2\pi^2}[\Cl(2\pi\frac aN)+\Cl(2\pi\frac bN)-\Cl(2\pi\frac{a+b}N)]=N^2\Phic(\frac aN,\frac bN)$, because
$-\Cl(\theta)=\Cl(2\pi-\theta)$. The first form follows by the substitution $\theta=2\pi\xi/N$, $\varphi=2\pi\eta/N$.
\end{proof}

\begin{lemma}[the classical value]\label{lem:classical}
Let $\alpha,\beta>0$ with $\alpha+\beta<1$. The function $f(\lambda)=\lambda\alpha+h_\lambda(\beta)$ is strictly convex on $\R$, and
\[
\min_{\lambda\in\R}f(\lambda)=\Phic(\alpha,\beta),\qquad\text{attained only at}\quad
\lambda_0=\frac1\pi\log\frac{\sin\pi(\alpha+\beta)}{\sin\pi\alpha}.
\]
In particular $\min_\lambda\bigl[\frac\lambda4+h_\lambda(\frac14)\bigr]=\Phic(\frac14,\frac14)=G/\pi^2$, attained only at $\lambda^*=\frac{\log2}{2\pi}$.
\end{lemma}

\begin{proof}
Since $h_\lambda(\beta)=h_0(\beta-i\lambda)$, Lemma~\ref{lem:h}(4) gives $f''(\lambda)=K_\beta(\lambda)>0$. Apply
Theorem~\ref{thm:continuum} with $M=1$ to $B=\one_{(-2\pi\beta,0)}$ and $A=\one_{(0,2\pi\alpha)}$. The conjugate function
$g_\beta(\theta)=\frac1\pi\log|\sin\frac{\theta+2\pi\beta}2/\sin\frac\theta2|$ has derivative
$\frac1{2\pi}[\cot\frac{\theta+2\pi\beta}2-\cot\frac\theta2]<0$ on $(0,2\pi(1-\beta))$, where it decreases from $+\infty$ to $-\infty$.
With $\lambda_0=g_\beta(2\pi\alpha)$, which is the stated value, $S_{\lambda_0}=0$ off finitely many points: on $(0,2\pi\alpha)$ both
terms equal $g_\beta-\lambda_0>0$, on $(2\pi\alpha,2\pi(1-\beta))$ both vanish, and on the arc of $B$ both vanish. For $M=1$,
$D_\lambda\equiv0$ (if $B(\theta)=1$ both terms of $D_\lambda$ vanish, if $B(\theta)=0$ both equal $(g-\lambda)_+$) and the bracket
in~\eqref{eq:deficit} vanishes. Hence $f(\lambda)-\int Ag_\beta\dd m=\int S_\lambda\ge0$ for all $\lambda$, with equality at $\lambda_0$, and
$\int Ag_\beta\dd m=\Phic(\alpha,\beta)$ by Lemma~\ref{lem:twoarcs}. For $\alpha=\beta=\frac14$,
$\lambda_0=\frac1\pi\log\frac1{\sin(\pi/4)}=\frac{\log2}{2\pi}$ and $\Phic=\frac1{2\pi^2}[2\Cl(\frac\pi2)+\Cl(\pi)]=G/\pi^2$.
\end{proof}

The identity $\min_\lambda f=\Phic$ is a Legendre-type identity, $\inf_y[h_0(x+iy)-py]=\Phic(x,p)$; it can also be derived
from the Bloch--Wigner evaluation of $\Im\Li$ at the vertex of a triangle with angles
$\pi(\alpha,\beta,1-\alpha-\beta)$~\cite{Zagier}. At $\alpha=\beta=\frac14$ it is the classical value
$\Im\Li(\frac{1-i}2)=\frac\pi8\log2-G$~\cite{Lewin}.

\begin{corollary}[equality]\label{cor:equality}
Let $B$ be a projection step field on $\C^M$ with $\int B\dd m=\frac14\one$, and let $0\le A\le\one-B$ be measurable with
$\int\tau(A)\dd m=\frac14$. Then
\[
\frac{G}{\pi^2}-\int_\T\tau(Ag)\dd m=\frac1M\int_\T D_{\lambda^*}\dd m+\int_\T S_{\lambda^*}\dd m\ \ge0 .
\]
If equality holds, then $[B(\theta),g(\theta)]=0$ for almost every $\theta$.
\end{corollary}

\begin{proof}
All $\beta_i$ equal $\frac14$, so the bracket in~\eqref{eq:deficit} vanishes; take $\lambda=\lambda^*$ and use
Lemma~\ref{lem:classical}. In the equality case $D_{\lambda^*}=0$ almost everywhere, and the equality clause of
Theorem~\ref{thm:C} gives $[B(\theta),g(\theta)]=0$.
\end{proof}

\begin{remark}[scope]\label{rem:measurable}
For $M=1$, Theorem~\ref{thm:continuum} contains the circle version of the bathtub bound obtained from the Stein--Weiss
theorem~\cite{SW,Laeng} (cf.~\cite[Proposition~4.6]{MS-classical}); the proof above uses only the mean value property of
$u_\lambda$. The theorem extends to measurable projection fields, by the standard boundary theory of Herglotz functions
and the exponential integrability of conjugate functions of bounded functions~\cite{Zygmund}; we only need step fields.
\end{remark}

\begin{remark}[non-commuting maximisers in the continuum]\label{rem:degenerate}
The bound $\Phic(\alpha,\beta)$ is attained by many fields, including non-commuting ones. For instance, if $\Theta$ is a
rational $M\times M$ inner function with $\Theta(0)=0$ (a finite Blaschke--Potapov product), the spectral projections
$B=\one_{(-2\pi\beta,0)}(\Theta(e^{i\theta}))$ and $A=\one_{(0,2\pi\alpha)}(\Theta(e^{i\theta}))$ satisfy $\int B\dd m=\beta\one$ and
$\int A\dd m=\alpha\one$ (because $\int\Theta(e^{i\theta})^k\dd m=\Theta(0)^k=0$ for $k\ge1$), the conjugate function of $B$ is
$g_\beta(\Theta(e^{i\theta}))$, and in the measurable version of Theorem~\ref{thm:continuum} (Remark~\ref{rem:measurable})
all inequalities of~\eqref{eq:deficit} are equalities at $\lambda_0$, since $A$, $B$ and $g$ commute pointwise; the field $B$ does not
commute with itself at different points unless the factors of $\Theta$ commute. Thus the continuum problem alone
cannot detect the structure of the optimal discrete strategies. In the proof of Theorem~\ref{thm:B} the needed
rigidity comes instead from the cell structure (Proposition~\ref{prop:residue}). We do not use this remark.
\end{remark}

%======================================================================
\section{Cell embeddings}\label{sec:cells}
%======================================================================

\subsection{Cells and kernels}
Let $\ell\in\Delta_d$, extended to $\Z$ by $\ell_x=\ell_{x\bmod d}$. Put $L_0=0$ and $L_{x+1}=L_x+\ell_x$, so that
$L_{x+d}=L_x+d$ and $L_{x+N}=L_x+N$, and define the cells $I_x=[L_x,L_{x+1})\subseteq\T_N:=\R/N\Z$, $x\in\Z_N$; their lengths
are $\ell_x$, and $S_r(m)=L_{r+m}-L_r$. We write $\kappa(u)=\cot(\pi u/N)$ also for real $u\notin N\Z$. For $x\ne y$ in $\Z_N$
let
\[
K^\ell(x,y)=\int_{I_x}\int_{I_y}\kappa(\xi-\eta)\dd\eta\dd\xi ,\qquad K^\ell(x,x)=0,
\]
and for a Q-configuration $Q$ put $Q^\ell(Q)=\sum_{x,y\in\Z_N}K^\ell(x,y)\,\tau(A_xB_y)$.

\begin{lemma}\label{lem:kernel}
For $x\ne y$ the integral defining $K^\ell(x,y)$ converges absolutely, and
\[
K^\ell(x,y)=\Xi(L_{x+1}-L_y)-\Xi(L_x-L_y)-\Xi(L_{x+1}-L_{y+1})+\Xi(L_x-L_{y+1}).
\]
Moreover $K^\ell(x+d,y+d)=K^\ell(x,y)=-K^\ell(y,x)$, and $K^\ell(x,y)=0$ if $\ell_x=0$ or $\ell_y=0$.
\end{lemma}

\begin{proof}
The cells are disjoint; $\kappa(\xi-\eta)$ is singular on $\bar I_x\times\bar I_y$ only at a common endpoint, where
$|\kappa(\xi-\eta)|\le C/|\xi-\eta|$ with $\xi,\eta$ on opposite sides of it, and
$\int_0^1\int_0^1\frac{\dd s\dd t}{s+t}<\infty$. The formula follows by integrating $\Xi''=\kappa$ twice, as in
Lemma~\ref{lem:twoarcs}. The symmetries follow from $I_{x+d}=I_x+d$ and the oddness of $\kappa$.
\end{proof}

\begin{proposition}[cell inequalities]\label{prop:cellineq}
For every $\ell\in\Delta_d$ and every Q-configuration $Q$ on $\C^M$,
\[
Q^\ell(Q)\le\frac{N^2G}{\pi^2},
\]
with equality for every window. More precisely, let $B$ and $A$ be the step fields on $\T$ equal to $B_x$ and $A_x$ on the
arc $J_x=(2\pi L_x/N,2\pi L_{x+1}/N)$, for every $x$ with $\ell_x>0$, and let $g$ be the conjugate function of $B$; then
$\int B\dd m=\int A\dd m=\frac14\one$, $Q^\ell(Q)=N^2\int_\T\tau(Ag)\dd m$, and Corollary~\ref{cor:equality} applies.
\end{proposition}

\begin{proof}
The arcs $J_x$ with $\ell_x>0$ are disjoint, nonempty and cover $\T$ up to finitely many points, so $B$ is a projection
step field; $0\le A\le\one-B$ because $A_xB_x=0$; and by~\eqref{eq:opconstraint},
$\int B\dd m=\frac1N\sum_x\ell_xB_x=\frac dN\one=\frac14\one$, and likewise $\int A\dd m=\frac14\one$, so
$\int\tau(A)\dd m=\frac14$. For $\theta\in J_x$, $\tau(A(\theta)g(\theta))=\sum_{y\ne x}\tau(A_xB_y)\int_{J_y}\cot\frac{\theta-\varphi}2\dd m(\varphi)$,
because $\tau(A_xB_x)=0$ (no principal value is needed). Integrating over $J_x$, summing over $x$, and substituting
$\theta=2\pi\xi/N$, $\varphi=2\pi\eta/N$ gives $\int\tau(Ag)\dd m=N^{-2}\sum_{x\ne y}K^\ell(x,y)\tau(A_xB_y)=N^{-2}Q^\ell(Q)$, all
integrals being absolutely convergent by Lemma~\ref{lem:kernel}. Corollary~\ref{cor:equality} gives the inequality. For
the window $\one_{[c,c+d)}$ ($M=1$), $A=1$ exactly on $[L_c,L_c+d)$ and $B=1$ exactly on $[L_c-d,L_c)$ (each window contains
every residue once), two adjacent arcs of length $d$; Lemma~\ref{lem:twoarcs} gives $Q^\ell=N^2\Phic(\frac14,\frac14)=N^2G/\pi^2$.
\end{proof}

\subsection{Rotation averages}
\begin{proposition}[averaged cell inequalities]\label{prop:average}
For every $\ell\in\Delta_d$ and every Q-configuration,
\begin{equation}\label{eq:average}
\frac1d\sum_{j=0}^{d-1}\Bigl[Q^\ell(\mathrm{rot}_jQ)-\frac{N^2G}{\pi^2}\Bigr]=\ip{u^\ell}\delta ,
\end{equation}
with $u^\ell$ as in~\eqref{eq:uell}. Equivalently $u^\ell_m=\tilde K(m)+\tilde K(2d-m)$ with
$\tilde K(m)=\frac1d\sum_{r=0}^{d-1}K^\ell(r+m,r)$. Consequently $\ip{u^\ell}\delta\le0$.
\end{proposition}

\begin{proof}
By definition $Q^\ell(\mathrm{rot}_jQ)=\sum_{x,y}K^\ell(x+j,y+j)\tau(A_xB_y)$. Since $K^\ell$ is invariant under
$(x,y)\mapsto(x+d,y+d)$, the average over $0\le j\le d-1$ equals the average over $j\in\Z_N$, which depends only on $x-y$
and equals $\tilde K(x-y)$. Hence the left side of~\eqref{eq:average}, plus $N^2G/\pi^2$, is
$\sum_{m\in\Z_N}\tilde K(m)T(m)$. The function $\tilde K$ is odd on $\Z_N$ (Lemma~\ref{lem:kernel}), so $\tilde K(0)=\tilde K(2d)=0$
and, by Lemma~\ref{lem:linearform}(1),
\[
\sum_m\tilde K(m)T(m)=\sum_{m=1}^{2d-1}\tilde K(m)\cN(m)=d\,\tilde K(d)+\sum_{m=1}^{d-1}\bigl[\tilde K(m)+\tilde K(2d-m)\bigr]\cN(m).
\]
For the window, $\cN(m)=m$ and every rotation is a window, so Proposition~\ref{prop:cellineq} gives
$N^2G/\pi^2=d\tilde K(d)+\sum_m[\tilde K(m)+\tilde K(2d-m)]m$. Subtracting gives~\eqref{eq:average} with
$u^\ell_m=\tilde K(m)+\tilde K(2d-m)$. It remains to identify this with~\eqref{eq:uell}. By Lemma~\ref{lem:kernel},
$K^\ell(r+m,r)=\Xi(S_r(m+1))-\Xi(S_r(m))-\Xi(S_{r+1}(m))+\Xi(S_{r+1}(m-1))$, and $S_{r+d}=S_r$, so
$\tilde K(m)=\Delta^2_m\Pi(m)$ with $\Pi(k)=\frac1d\sum_r\Xi(S_r(k))$. For $0\le j\le d$,
$S_r(2d-j)=d+S_r(d-j)=2d-S_{r+d-j}(j)$, so $\Pi(2d-j)=\frac1d\sum_r\Xi(2d-S_r(j))$, and
$\tilde K(2d-m)=\Delta^2_m\bigl[\frac1d\sum_r\Xi(2d-S_r(m))\bigr]$ for $1\le m\le d-1$. Adding gives~\eqref{eq:uell}. The
last assertion follows from Proposition~\ref{prop:cellineq}.
\end{proof}

\begin{proposition}[the cone condition implies optimality]\label{prop:conetoopt}
Let $d\ge2$, and suppose that $v=\int u^\ell\mu(\dd\ell)+\sum_sw_st_s$ for a finite positive Borel measure $\mu$ on $\Delta_d$
and $w_s\ge0$. Then $\ip v\delta\le0$ for every Q-configuration on every $\C^M$; hence $F(V)\le\FD$ for every family $V$ of
order-four unitaries, and $I_d\le\IME(d)$ for every projective strategy on a maximally entangled state.
\end{proposition}

\begin{proof}
The map $\ell\mapsto u^\ell$ is continuous on the compact set $\Delta_d$ (as $\Sh$ is continuous on $[0,2d]$), so the integral
is defined, and $\ip v\delta=\int\ip{u^\ell}\delta\mu(\dd\ell)+\sum_sw_s\ip{t_s}\delta\le0$ by
Proposition~\ref{prop:average} and Lemma~\ref{lem:linearform}(3). The rest is Lemma~\ref{lem:linearform}(2) and
Theorem~\ref{thm:reduction}.
\end{proof}

\begin{remark}[uniform cells and junction terms]\label{rem:uniform}
For uniform cells $\ell\equiv1$, $K^1(x,y)=\bar\kappa(x-y)$ with $\bar\kappa(m)=\int_{-1}^1(1-|s|)\kappa(m+s)\dd s$, and
$u^1=v+j$, where $j_m=|e_{\rm J}(m)|+|e_{\rm J}(2d-m)|\ge0$ and $e_{\rm J}(m)=\kappa(m)-\bar\kappa(m)$ is the junction function
of~\cite[Lemma~4.9]{MS-classical}. So the uniform embedding gives $\ip v\delta\le-\ip j\delta$, a bound with an error made of
junction terms; this is how the commutative case is treated in~\cite{MS-classical}, where those terms are estimated
using the cyclic order of optimal colourings, an argument with no known operator analogue. The cone condition replaces
this estimate: since $u^1$ is not proportional to $v$, non-uniform cells are needed.
\end{remark}

\begin{remark}[Poisson cells]\label{rem:poisson}
Let $\ell=d\cdot(Y_0,\dots,Y_{d-1})$ with $(Y_r)$ uniformly distributed on the probability simplex (the Dirichlet law
with all parameters $1$; equivalently, the cell boundaries are $d$ independent uniform points on a circle of length
$d$). Then $S_r(m)/d$ has the Beta$(m,d-m)$ law, and the Bernstein derivative formula gives
$\E u^\ell_m=\int_0^d(\Sh)''(s)\,b_{m,d}(s/d)\dd s$ with $b_{m,d}(x)=\binom dmx^m(1-x)^{d-m}$. As
$(\Sh)''(s)=\frac{4d}{\pi s}+O(1)$ and $\int_0^1b_{m,d}(x)\frac{\dd x}x=\frac1m$, the averaged cell vector reproduces
\emph{exactly} the singular part $\frac{4d}{\pi m}$ of $v_m$: the discrete kernel is the continuum kernel averaged over
Poisson-like cells. The residual $\E u^\ell-v$ is smooth and of relative size $O(1/d)$; the constructions of
Section~\ref{sec:cone} correct it by modulating the Dirichlet parameters. (This remark explains the construction and
is not used in the proofs; Lemma~\ref{lem:singular} below is the precise statement that is used.)
\end{remark}

%======================================================================
\section{The cone condition for every \texorpdfstring{$d$}{d} (computer-assisted)}\label{sec:cone}
%======================================================================

This section proves Theorem~\ref{thm:E}. The proof is computer-assisted: it depends on computations in rigorous
interval arithmetic (the interval type of the Python library \texttt{mpmath}~\cite{mpmath}, with outward rounding, at
$160$ bits for $d\le2000$ and at $120$--$240$ bits for $d\ge2001$). Below we state the mathematical reductions, with
proofs, and then exactly what is verified, by which program, and with which margins. Programs, certificates and logs
are listed in Section~\ref{sec:verification}.
There are three regimes: $2\le d\le200$ (Section~\ref{subsec:regimeI}), $201\le d\le2000$
(Section~\ref{subsec:regimeII}) and $d\ge2001$ (Section~\ref{subsec:regimeIII}).

\subsection{Potentials}
For a probability measure $\mu$ on $\Delta_d$ and $0\le m\le d$ let
$P_\mu(m)=\int\frac1d\sum_{r=0}^{d-1}\Sh(S_r(m))\,\mu(\dd\ell)$. Since $S_r(0)=0$ and $S_r(d)=d$, $P_\mu(0)=\Sh(0)$ and
$P_\mu(d)=\Sh(d)$ for every $\mu$, and $\int u^\ell\mu(\dd\ell)=\Delta^2P_\mu$ on $\{1,\dots,d-1\}$. Let $P_v$ be the solution of
$\Delta^2P_v=v$ on $\{1,\dots,d-1\}$ with $P_v(0)=\Sh(0)$ and $P_v(d)=\Sh(d)$. For $X\colon\{0,\dots,d\}\to\R$ write
$X_A(m)=\frac12(X(m)-X(d-m))$ and $X_S(m)=\frac12(X(m)+X(d-m))$.

\begin{lemma}[reduction to potentials]\label{lem:R1}
Suppose that $\mu$ satisfies (i) $(P_\mu)_A=(P_v)_A$, and (ii) $W:=(P_v-P_\mu)_S$ has $\Delta^2W(m)\ge0$ for $1\le m\le d-1$.
Then~\eqref{eq:cone} holds with this $\mu$ and suitable $w_s\ge0$.
\end{lemma}

\begin{proof}
$W$ is symmetric with $W(0)=W(d)=0$, and by (i) $P_\mu+W=P_v$; applying $\Delta^2$ gives $\int u^\ell\dd\mu+\Delta^2W=v$. The
vector $\Delta^2W$ is symmetric under $m\mapsto d-m$ and nonnegative, hence a nonnegative combination of the $t_s$.
\end{proof}

\subsection{Regime I: \texorpdfstring{$2\le d\le200$}{d <= 200}}\label{subsec:regimeI}

\begin{proposition}[computer-assisted]\label{prop:regimeI}
For every $2\le d\le200$ there are $K_d\le6.3\,d$ cell vectors $\ell^{(k)}=n^{(k)}/16$ with integer entries
$n^{(k)}_r\ge1$ and $\sum_rn^{(k)}_r=16d$, and weights $c_k>0$ with $d\cdot\min_kc_k\ge0.0816$, such that
$\sum_kc_ku^{\ell^{(k)}}=v$. In particular Theorem~\ref{thm:E} holds with $\mu=\sum_kc_k\delta_{\ell^{(k)}}$, which is supported in
$\Delta_d^+$, and $w=0$.
\end{proposition}

\emph{What is verified.} Since all window sums of $\ell=n/16$ are multiples of $1/16$, only the values $\Sh(j/16)$,
$0\le j\le16d$, are needed, and
$\Sh(j/q)=-\frac{N^2}{2\pi^2}[\Cl(\pi j/K_q)+\Cl(\pi(K_q-j)/K_q)]$ with $q=16$, $K_q=2dq$. Each value $\Cl(t)$, $0<t\le\pi$, is
enclosed using
\[
\Cl(t)=t-t\log t+\sum_{n=1}^{70}\frac{|B_{2n}|\,t^{2n+1}}{2n\,(2n+1)!}+R_{70}(t),\qquad
0\le R_{70}(t)\le\frac{t\,\pi^2/6}{\frac34\cdot71\cdot143\cdot4^{71}},
\]
with exact rational Bernoulli numbers; the tail bound follows from $|B_{2n}|=2(2n)!\zeta(2n)/(2\pi)^{2n}$,
$\zeta(2n)\le\zeta(2)$ and $(t/2\pi)^{2n}\le4^{-n}$. The vectors $u^{\ell^{(k)}}$ are then enclosed through the window-sum
form~\eqref{eq:uell}. The cell vectors and approximate weights $\tilde c$ were found by floating-point linear
programming (HiGHS~\cite{HiGHS}); this search is not part of the verification. The exact weights are
$c=\tilde c+U^{\mathsf T}y$, where $U$ is the matrix with columns $u^{\ell^{(k)}}$ and $(UU^{\mathsf T})y=v-U\tilde c$; then
$Uc=v$ exactly. With an approximate inverse $R$ of $UU^{\mathsf T}$, the bound
$\|y\|_\infty\le\|R(v-U\tilde c)\|_\infty/(1-\|\one-RUU^{\mathsf T}\|_\infty)$ holds as soon as $\|\one-RUU^{\mathsf T}\|_\infty<1$; all
quantities are computed in interval arithmetic (160 bits), so the bound holds for the exact $U$. The program
\texttt{verify\_cone.py} (driver \texttt{run\_verify.py}) certifies, for every $2\le d\le200$:
$\|\one-RUU^{\mathsf T}\|_\infty\le9.4\cdot10^{-34}$, $\|y\|_\infty\le2.1\cdot10^{-11}$, and $\min_kc_k>0$. Every $d$ was certified
at the first attempt; $3.5d\le K_d\le6.3d$; the smallest certified weight is $4.25\cdot10^{-4}$ (at $d=198$), and
$0.0816\le d\cdot\min_kc_k\le0.2572$. The certificates contain the integer vectors $n^{(k)}$ and the weights. A complete
re-verification of all 199 stored certificates (program \texttt{lp\_reverify.py}, which reruns \texttt{verify\_cone.py})
passed and reproduced every certified bound. For $2\le d\le20$ the same certificates are also re-evaluated by the Lean
kernel in exact rational interval arithmetic (Section~\ref{subsec:formal}).

\emph{Independent check.} The program \texttt{crosscheck\_certs.py} recomputes $\sum_kc_ku^{\ell^{(k)}}-v$ without interval
arithmetic, with \texttt{mpmath.clsin} at 50 digits and, for $d\le10$, with the cell-pair kernel of
Lemma~\ref{lem:kernel} instead of the window-sum form: the maximal deviation is between $3\cdot10^{-15}$ and
$2.5\cdot10^{-13}$ for $d\in\{3,5,8,10,20,37\}$, and at most $1.5\cdot10^{-10}$ for $d\in\{150,200\}$ (weights stored in
double precision).

\subsection{Dirichlet cells}
For $0<b<d$ let $X_b$ have the Beta$(b,d-b)$ law, and put $\Xb(b)=\E\,\Sh(dX_b)$, $\Xb(0)=\Sh(0)$, $\Xb(d)=\Sh(d)$. If
$\beta_0,\dots,\beta_{d-1}>0$ with $\sum_r\beta_r=d$ and $\ell/d$ has the Dirichlet law with parameters $(\beta_r)$, then by the
aggregation property of the Dirichlet law $S_r(m)/d$ has the Beta$(B_r(m),d-B_r(m))$ law, $B_r(m)=\beta_r+\dots+\beta_{r+m-1}$,
so that
\begin{equation}\label{eq:dirichlet}
\E\Bigl[\frac1d\sum_r\Sh(S_r(m))\Bigr]=\frac1d\sum_r\Xb\bigl(B_r(m)\bigr).
\end{equation}

\begin{lemma}[the singular part is reproduced exactly]\label{lem:singular}
Write $\Sh=\Xi_s+\Xi_r$ with $\Xi_s(u)=\frac{4d}\pi u\log\frac ud$ and $\Xi_r(u)=d^2g_r(u/d)$, where
\[
g_r(x)=-\frac8{\pi^2}\Bigl[\Cl\bigl(\tfrac{\pi x}2\bigr)+\Cl\bigl(\pi-\tfrac{\pi x}2\bigr)\Bigr]-\frac4\pi x\log x
\]
extends to an odd analytic function on the disc $|x|<2$. Let $\Xb_s,\Xb_r$ be the corresponding parts of $\Xb$, and let
$P_{v,s},P_{v,r}$ be the solutions of $\Delta^2P=\Xi_s''$, respectively $\Xi_r''$, on $\{1,\dots,d-1\}$ with the boundary values
of $\Xi_s$, respectively $\Xi_r$, at $0$ and $d$ (so $P_v=P_{v,s}+P_{v,r}$). Then for every integer $0\le m\le d$,
\[
\Xb_s(m)=P_{v,s}(m)=\frac{4d}\pi\Bigl[m\psi(m)+1-\frac md-m\psi(d)\Bigr],
\]
where $\psi$ is the digamma function and $m\psi(m)$ is read as $-1$ at $m=0$. Consequently $P_v-\Xb=P_{v,r}-\Xb_r$
involves only the analytic function $g_r$.
\end{lemma}

\begin{proof}
Near $x=0$, $\Cl(\theta)=\theta-\theta\log\theta+(\text{odd power series})$ shows that the logarithms cancel; $g_r$ is analytic for
$|x|<2$ because $\Cl$ is analytic on $(0,2\pi)$. For $X\sim$ Beta$(b,d-b)$, $\E[X\log X]=\frac bd[\psi(b+1)-\psi(d+1)]$, so
$\Xb_s(b)=\frac{4d}\pi b[\psi(b+1)-\psi(d+1)]$; at integers use $m\psi(m+1)=m\psi(m)+1$ and $\psi(d+1)=\psi(d)+\frac1d$. The right
side vanishes at $m=0$ and $m=d$, as do $\Xi_s(0)$ and $\Xi_s(d)$, and $\Delta^2[m\psi(m)]=\frac1m$ for $m\ge1$, while
$\Xi_s''(m)=\frac{4d}{\pi m}$; so both sides solve the same discrete Dirichlet problem.
\end{proof}

\subsection{Gaussian modulation}
Call $\Upsilon\colon\{0,\dots,d\}\to[0,\infty)$ a \emph{variogram} if $\Upsilon(m)=\Upsilon(d-m)$, $\Upsilon(0)=\Upsilon(d)=0$, and, extending
$\Upsilon$ $d$-periodically,
\[
\hat a_k(\Upsilon):=-\frac2d\sum_{m=0}^{d-1}\Upsilon(m)\cos\frac{2\pi km}d\ \ (1\le k<\tfrac d2),\qquad
\hat a_{d/2}(\Upsilon):=-\frac1d\sum_{m=0}^{d-1}(-1)^m\Upsilon(m)\ \ (d\text{ even})
\]
are all nonnegative. If all $\hat a_k(\Upsilon)>0$, the function $c(j)=\frac12(\Upsilon(j+1)-2\Upsilon(j)+\Upsilon(j-1))$ on $\Z_d$ has
discrete Fourier transform $\hat c(k)=d\sin^2(\pi k/d)\,\hat a_k(\Upsilon)$ for $1\le k<d/2$ (and $2d\,\hat a_{d/2}$ at $k=d/2$),
and $\hat c(0)=0$; so there is a centred Gaussian vector $(\eta_r)_{r\in\Z_d}$, invariant in law under rotations, with
covariance $c$, with $\sum_r\eta_r=0$, with $\mathrm{Var}(\eta_r+\dots+\eta_{r+m-1})=\Upsilon(m)$, and with a positive
density on the hyperplane $\{\sum_r\eta_r=0\}$. Let $E=\{\max_r|\eta_r|<\frac34\}$ and let $\mu_\Upsilon$ be the law of the random
cell vector $\ell$ which, given $\eta$, is $d$ times a Dirichlet vector with parameters $(1+\eta_r)$ on $E$ and with all
parameters $1$ on the complement $E^c$.

For $1\le m\le d-1$ let $L_m=\frac34\min(m,d-m)$, and for $t\ge0$ and $\zeta\sim N(0,1)$ let
\[
H(m,t)=\E\bigl[\Xb(m+\sqrt t\zeta);|\sqrt t\zeta|<L_m\bigr]+\Xb(m)\,\Pr(|\sqrt t\zeta|\ge L_m).
\]
For $1\le m<d/2$ put $\mathrm{Hc}_m(t)=[H(m,t)-\Xb(m)]-[H(d-m,t)-\Xb(d-m)]$ and $\tau_m=\mathrm{Dl}(m)-\mathrm{Dl}(d-m)$ with
$\mathrm{Dl}=P_v-\Xb$. The \emph{decoupled equations} are
\begin{equation}\label{eq:decoupled}
\mathrm{Hc}_m(\Upsilon(m))=\tau_m\qquad(1\le m<d/2).
\end{equation}

\begin{lemma}[decoupling]\label{lem:decoupling}
Let $\Upsilon$ be a variogram with all $\hat a_k(\Upsilon)>0$, and put
\begin{gather*}
\varepsilon_{\rm t}=\frac{3d}8\,\sup_{0<b<d}|\Xb'(b)|\cdot2d\exp\Bigl(-\frac9{32\,\Upsilon(1)}\Bigr),\\
e_m(\Upsilon)=[P_{\mu_\Upsilon}(m)-\Xb(m)]-[P_{\mu_\Upsilon}(d-m)-\Xb(d-m)]-\mathrm{Hc}_m(\Upsilon(m)).
\end{gather*}
Then $|P_{\mu_\Upsilon}(m)-H(m,\Upsilon(m))|\le\varepsilon_{\rm t}$ for $1\le m\le d-1$, and $|e_m(\Upsilon)|\le2\varepsilon_{\rm t}$ for $1\le m<d/2$.
\end{lemma}

\begin{proof}
By~\eqref{eq:dirichlet}, rotation invariance of $(\eta,E)$, and $B_r(m)=m+Z_r(m)$ with $Z_r(m)=\eta_r+\dots+\eta_{r+m-1}$,
$P_{\mu_\Upsilon}(m)=\E[\Xb(m+Z);E]+\Pr(E^c)\Xb(m)$ with $Z=Z_0(m)\sim N(0,\Upsilon(m))$. On $E$, $|Z|<\frac34m$, and since
$\sum_r\eta_r=0$ also $|Z|<\frac34(d-m)$; so $\{|Z|\ge L_m\}\subseteq E^c$, and
$P_{\mu_\Upsilon}(m)-H(m,\Upsilon(m))=-\E[(\Xb(m+Z)-\Xb(m))\one_{|Z|<L_m}\one_{E^c}]$, which is bounded by $L_m\sup|\Xb'|\Pr(E^c)$.
Finally $\Pr(E^c)\le\sum_r\Pr(|\eta_r|\ge\frac34)\le2d\exp(-\frac9{32\Upsilon(1)})$ since $\mathrm{Var}\,\eta_r=\Upsilon(1)$, and
$L_m\le\frac{3d}8$. The bound on $e_m$ follows since $\Upsilon(d-m)=\Upsilon(m)$.
\end{proof}

\begin{lemma}[from the decoupled to the exact equations]\label{lem:B}
Let $\mathrm{lo}_m<\mathrm{hi}_m$ for $1\le m<d/2$, and let $\mathcal Q$ be the set of variograms with
$\Upsilon(m)\in[\mathrm{lo}_m,\mathrm{hi}_m]$ for $1\le m<d/2$ (and a fixed value $\Upsilon(d/2)$ if $d$ is even). Suppose that
\begin{enumerate}
\item[(B1)] $\hat a_k(\Upsilon)>0$ for all $k$ and all $\Upsilon\in\mathcal Q$;
\item[(B2)] $\varepsilon_{\rm t}\le\bar\varepsilon$ for all $\Upsilon\in\mathcal Q$;
\item[(B3)] $\mathrm{Hc}_m(\mathrm{lo}_m)<\tau_m-2\bar\varepsilon$ and $\mathrm{Hc}_m(\mathrm{hi}_m)>\tau_m+2\bar\varepsilon$ for all $m$.
\end{enumerate}
Then some $\Upsilon^\#\in\mathcal Q$ satisfies $(P_{\mu_{\Upsilon^\#}})_A=(P_v)_A$, that is, condition (i) of Lemma~\ref{lem:R1}.
\end{lemma}

\begin{proof}
Condition (i) for $\mu_\Upsilon$ is the system $F_m(\Upsilon):=\mathrm{Hc}_m(\Upsilon(m))+e_m(\Upsilon)-\tau_m=0$, $1\le m<d/2$. By (B1)
the Gaussian law depends continuously on $\Upsilon\in\mathcal Q$ and has a positive density on the hyperplane
$\{\sum\eta_r=0\}$, so the boundary of $E$ is a null set and $\Upsilon\mapsto P_{\mu_\Upsilon}$ is continuous; hence $F$ is continuous on
the box $\mathcal Q$. By Lemma~\ref{lem:decoupling}, (B2) and (B3), $F_m<0$ on the face $\Upsilon(m)=\mathrm{lo}_m$ and $F_m>0$ on the
face $\Upsilon(m)=\mathrm{hi}_m$. The Poincar\'e--Miranda theorem~\cite{Kulpa} gives a zero of $F$ in $\mathcal Q$.
\end{proof}

Condition (B3) is obtained from a sign change and a derivative bound: if $\mathrm{Hc}_m(a_m)<\tau_m<\mathrm{Hc}_m(b_m)$,
$\mathrm{lo}_m\le a_m-\varrho$, $\mathrm{hi}_m\ge b_m+\varrho$, and $\mathrm{Hc}_m'\ge c_*>0$ on $[\mathrm{lo}_m,\mathrm{hi}_m]$ with
$c_*\varrho\ge2\bar\varepsilon$, then $\mathrm{Hc}_m(\mathrm{lo}_m)\le\mathrm{Hc}_m(a_m)-c_*\varrho<\tau_m-2\bar\varepsilon$, and
similarly at $\mathrm{hi}_m$, so (B3) holds. For condition (ii) of Lemma~\ref{lem:R1} at $\Upsilon^\#$, note that
$W=(P_v-P_{\mu_{\Upsilon^\#}})_S$ and that $|P_{\mu_\Upsilon}(m)-H(m,\Upsilon(m))|\le\bar\varepsilon$ by
Lemma~\ref{lem:decoupling}, while both functions take the same values at $m=0$ and $m=d$. So it suffices that
$\Delta^2(P_v-H(\cdot,\Upsilon))_S>4\bar\varepsilon$ on $\{1,\dots,d-1\}$ for every $\Upsilon\in\mathcal Q$. Since moreover
$(P_v-P_{\mu_{\Upsilon^\#}})_A=0$, one can instead use the one-sided slack $X_\Upsilon(m):=P_v(m)-H(m,\Upsilon(m))$
($1\le m\le d/2$), extended by $X_\Upsilon(d-m)=X_\Upsilon(m)$ and $X_\Upsilon(0)=0$: then $W=X_{\Upsilon^\#}+e$ with
$|e|\le\bar\varepsilon$, and it suffices that $\Delta^2X_\Upsilon>4\bar\varepsilon$ for every $\Upsilon\in\mathcal Q$.

\subsection{Regime II: \texorpdfstring{$201\le d\le2000$}{201 <= d <= 2000}}\label{subsec:regimeII}

\begin{proposition}[computer-assisted]\label{prop:regimeII}
For every $201\le d\le2000$ the hypotheses of Lemma~\ref{lem:B} and condition (ii) of Lemma~\ref{lem:R1} hold for explicit
intervals $[\mathrm{lo}_m,\mathrm{hi}_m]$. Hence Theorem~\ref{thm:E} holds with $\mu=\mu_{\Upsilon^\#}$, a probability measure with
$\mu(\Delta_d^+)=1$.
\end{proposition}

\emph{What is verified.} For each $201\le d\le2000$, one run of the program \texttt{verify\_cert.py} certifies all
hypotheses of Lemma~\ref{lem:B} and condition (ii) of Lemma~\ref{lem:R1}. Every quantity is computed in one process, at
one precision ($160$-bit interval arithmetic throughout) and by one version of the code, whose SHA-256 hash is stored in
every certificate (Section~\ref{sec:verification}). The run computes:
\begin{enumerate}
\item enclosures of $\Xb^{(2j)}(m)$ for $j\le7$ and all integers $1\le m\le d-1$ (and bounds for $\Xb^{(16)}$ on discs, used
in (3)): the singular part through $\Xb_s(b)=\frac{4d}\pi b[\psi(b+1)-\psi(d+1)]$ and the polygamma values
$\psi^{(n)}(m+1)=(-1)^{n+1}n![\zeta(n+1)-H^{(n+1)}_m]$, where $\zeta(s)$, $2\le s\le17$, is enclosed by the
Euler--Maclaurin formula in exact rational arithmetic with its remainder bound; the regular part through the series
$\Xb_r(b)=d^2\sum_{i\text{ odd}}g_i\frac{(b)_i}{(d)_i}$ ($g_i$ the Taylor coefficients of $g_r$, $(b)_i$ the rising
factorial, $\E X_b^i=(b)_i/(d)_i$), whose Taylor coefficients at the integers are computed by exact truncated products for
$i\le141$, with a tail bound from $0\le\frac{\dd^k}{\dd b^k}\frac{(b)_i}{(d)_i}\le k!\binom ik/(d)_k$ on $[0,d]$;
\item the targets $\tau_m$, through Lemma~\ref{lem:singular} and the discrete potential $P_{v,r}$ (prefix sums), with
$g_r(1)$ enclosed including the tail of its series;
\item root boxes $[a_m,b_m]$ for~\eqref{eq:decoupled}: the Gaussian average is expanded exactly to order $14$,
$H(m,t)-\Xb(m)=\sum_{j=1}^7t^j\Xb^{(2j)}(m)/(2^jj!)+R(m,t)$, where $R$ contains the order-$16$ Lagrange remainder (bounded
through $\sup|\Xb^{(16)}|$ on $|z|\le\frac12\min(m,d-m)$ and a Gaussian tail estimate on the rest of the window, using
$|\psi^{(n)}(z)|\le(n-1)!z^{-n}+n!z^{-n-1}$) and the truncation $|\sqrt t\zeta|<L_m$ (truncated Gaussian moments); an
untrusted search proposes $[a_m,b_m]$, and a certified sign change of $\mathrm{Hc}_m-\tau_m$ at its ends is checked;
\item the tail quantities, in logarithms: $\sup|\Xb'|\le\frac{4d}\pi(H_d+1)+S_1d$ ($H_d$ the harmonic number,
$S_1=\sum_{i\text{ odd}}i|g_i|\le1.8884$), $\Pr(E^c)$ and $\bar\varepsilon$ from Lemma~\ref{lem:decoupling}, and a
radius $\varrho=2^{e}$, $e\in\Z$, with $\varrho\ge8\bar\varepsilon/c_{1,\min}$, where
$c_{1,\min}=\min_m\frac12(\Xb''(m)-\Xb''(d-m))$;
\item the box $\mathcal Q=\prod_m[\mathrm{lo}_m,\mathrm{hi}_m]$ with $\mathrm{lo}_m\le a_m-\varrho$ and
$\mathrm{hi}_m\ge b_m+\varrho$ (outward rounding), and, on the whole of $\mathcal Q$: condition (B1); condition (B2),
using $\Upsilon(1)\le\mathrm{hi}_1$; the bound $c_D:=\min_m\inf\{\mathrm{Hc}_m'(t):t\in[\mathrm{lo}_m,\mathrm{hi}_m]\}>0$
with $c_D\varrho>2\bar\varepsilon$, which gives (B3) by the remark after Lemma~\ref{lem:B} (the derivative of the
order-$14$ expansion and of the remainder is bounded with the heat equation $\partial_t\phi_t=\frac12\phi_t''$ for the
Gaussian density and two integrations by parts on $[-L_m,L_m]$); condition (B3) once more, directly, by interval
evaluation of $\mathrm{Hc}_m-\tau_m$ at $\mathrm{lo}_m$ and $\mathrm{hi}_m$; and
$\Delta^2(P_v-H(\cdot,\Upsilon))_S>4\bar\varepsilon$ on $\{1,\dots,d-1\}$ for all $\Upsilon\in\mathcal Q$, which gives
condition (ii) by the same remark.
\end{enumerate}
All factorials, binomial coefficients, Bernoulli numbers and rational constants enter as exact integers or rationals
with outward rounding, the results of the elementary interval functions of \texttt{mpmath} are widened by a further
$16$ units in the last place, every decision is a comparison of interval endpoints, and every certified bound is stored
as the exact decimal expansion of the corresponding endpoint. The program \texttt{audit\_cert\_g.py} re-checks every
condition from these stored values, together with the code hash: all $1800$ certificates pass. The smallest certified
values over $201\le d\le2000$ are: $d^2\min_k\hat a_k\ge0.06319$ for even $d$ (at $d=2000$; the minimum over $k$ is at
$k=\lfloor d/2\rfloor$) and $d^2\min_k\hat a_k\ge0.12638$ for odd $d$; $d\cdot\min_m\Delta^2(P_v-H)_S\ge0.660$;
$c_D/(c_{1,\min}/4)\ge4.0009$ (a ratio above $1$ is needed) and $c_D\varrho/(2\bar\varepsilon)\ge10^{0.61}$; the face
values of (B3) exceed $2\bar\varepsilon$ by a factor of at least $10^{16.7}$; $\Pr(E^c)\le4.2\cdot10^{-37}$,
$\bar\varepsilon\le10^{-31.16}$ and $\varrho\le2^{-94}$. (Earlier runs, which used separate programs for the Gaussian and
the fixed-point conditions and $100$-bit arithmetic for $d\le600$, gave lower apparent margins, $0.0285$ and $0.0798$;
they are superseded by the certificates described here.)

\subsection{Regime III: \texorpdfstring{$d\ge2001$}{d >= 2001}}\label{subsec:regimeIII}
For large $d$ the variogram coefficients of the decoupled solution are of size $d^{-2}$ at high frequencies and are
controlled by a local criterion. Let $\varepsilon=1/d$, $A_0=\frac\pi2-1$,
\[
\mathrm{B}(b)=b\,[2\psi'(b+1)+b\,\psi''(b+1)],\qquad f(b)=b^2\Bigl(\frac1{\mathrm{B}(b)}-1\Bigr)\ (b>0),\qquad f(0)=0,
\]
$a(m)=\frac16-f(m)$ for $m\ge0$, and $F_A(m)=A_0\varepsilon[f(m)+f(d-m)-f(d)]$ for $0\le m\le d$. For a symmetric function $X$ on $\Z_d$ let $\Delta^4X$ be its
centred periodic fourth difference.

\begin{lemma}[positivity criterion]\label{lem:P}
\begin{enumerate}
\item \emph{(computer-assisted)} The sequence $(a(m))_{m\ge0}$ is positive and convex, and $m^2a(m)\le\frac{13}{180}$ for $m\ge1$.
Consequently, for all $\theta$,
\[
\mathcal A(\theta):=a(0)+2\sum_{m\ge1}a(m)\cos m\theta\ \ge\ \Delta^2a(1)=2f(1)-f(2)\ \ge\ 0.105487,
\]
where $f(1)=(\pi^2/3-2\zeta(3))^{-1}-1$ and $f(2)=4\,(2\pi^2/3+4-8\zeta(3))^{-1}-4$.
\item Let $d\ge5$, let $\Upsilon_e$ be symmetric on $\Z_d$ with $\Upsilon_e(0)=0$, and let $D_B=\Delta^4(\Upsilon_e-F_A)$. If $D_B(1)\ge-\gamma\varepsilon^2$
and $D_B(m)\ge\beta'\varepsilon^3$ for $2\le m\le d/2$, with $\gamma\ge0$ and $\beta'>0$, then for $1\le k\le d/2$
\[
\hat a_k(\Upsilon_e)\ge w_k\frac{\varepsilon^2}8\Bigl[P_d-\frac{4\gamma^2}{\beta'd(1-4/d)}\Bigr],\qquad P_d=16A_0\Bigl(\inf_\theta\mathcal A(\theta)-\frac{0.1446}{d-1}\Bigr),
\]
where $w_k=1$ for $k<d/2$ and $w_{d/2}=\frac12$.
\item If $|\Upsilon(m)-\Upsilon_e(m)|\le\delta_{\max}$ for all $m$, then $\hat a_k(\Upsilon)\ge\hat a_k(\Upsilon_e)-2w_k\delta_{\max}$.
\end{enumerate}
\end{lemma}

\begin{proof}
(1) The positivity, convexity and the bound $m^2a(m)\le\frac{13}{180}$ are verified in interval arithmetic for $m\le3000$,
and for $m>3000$ from the asymptotic expansion $f(b)=\frac16-\frac{13}{180}b^{-2}+\frac{683}{7560}b^{-4}+E(b)$,
$|E(b)|\le0.409\,b^{-6}$, obtained from the Binet expansion of the polygamma functions with its explicit remainder
(program \texttt{t35\_polya.py}). Since $a$ is convex with $a(m)\to0$, P\'olya's argument writes
$\mathcal A=\sum_{L\ge1}\Delta^2a(L)\,L\,\mathcal K_L$ with the nonnegative Fej\'er kernels
$\mathcal K_L(\theta)=L^{-2}\bigl(\sin\frac{L\theta}2/\sin\frac\theta2\bigr)^2$, and $L\mathcal K_L\equiv1$ for $L=1$. The values of
$f(1),f(2)$ follow from $\psi'(2)=\frac{\pi^2}6-1$, $\psi''(2)=2-2\zeta(3)$, $\psi'(3)=\frac{\pi^2}6-\frac54$, $\psi''(3)=\frac94-2\zeta(3)$.

(2) For symmetric $X$ on $\Z_d$ with $X(0)=0$, $\theta=2\pi k/d$ and $s=\sin\frac\theta2$, summation by parts gives
$8ds^4\hat a_k(X)=w_k\sum_{m=1}^{d-1}\Delta^4X(m)(1-\cos m\theta)$, because the symbol of $\Delta^4$ is $16s^4$ and
$\sum_m\Delta^4X(m)=0$. For $X=F_A$, using $F_A(m)=A_0\varepsilon[\frac16+a(d)-a(m)-a(d-m)]$ (valid also at $m=0$),
$\sum_m\Delta^4F_A(m)(1-\cos m\theta)=16A_0\varepsilon s^4[a(0)+2\sum_{m=1}^{d-1}a(m)\cos m\theta+a(d)]
\ge16A_0\varepsilon s^4[\mathcal A(\theta)-2\sum_{m\ge d}a(m)]\ge P_d\,\varepsilon s^4$, since $2\sum_{m\ge d}a(m)\le\frac{26}{180(d-1)}<\frac{0.1446}{d-1}$.
For the remainder, $D_B(d-m)=D_B(m)$, $1-\cos\theta=2s^2$, and $\sum_{m=2}^{d-2}(1-\cos m\theta)=d-2+2\cos\theta\ge d-4$, so
$\sum_mD_B(m)(1-\cos m\theta)\ge-4\gamma\varepsilon^2s^2+\beta'\varepsilon^2(1-4\varepsilon)$. Dividing by $8ds^4$ and putting $y=s^{-2}\ge1$,
$\hat a_k(\Upsilon_e)\ge w_k\frac{\varepsilon^2}8[P_d-4\gamma\varepsilon y+\beta'\varepsilon(1-4\varepsilon)y^2]$, and the minimum over $y\in\R$ of the
bracket is the stated bound.

(3) $|\hat a_k(\delta)|\le\frac2d\sum_m|\delta(m)|\le2\delta_{\max}$, and half of this for $k=d/2$.
\end{proof}

For example, $\gamma=3$, $\beta'=\frac14$ and $\delta_{\max}=\varepsilon^2/20$ in Lemma~\ref{lem:P}(2),(3) give
$\hat a_k(\Upsilon)\ge w_k\varepsilon^2\bigl[\frac18\bigl(0.962733-\frac{36}{0.25\cdot2001\cdot(1-4/2001)}\bigr)-\frac1{10}\bigr]
\ge0.011\,w_k\varepsilon^2$ for every $d\ge2001$; the bounds certified below are much better.

\begin{proposition}[computer-assisted, uniform in $d$]\label{prop:regimeIII}
Let $\Upsilon_e$ be the explicit model of Appendix~\ref{app:model}. For every $d\ge2001$, with $\varepsilon=1/d$, let
$\bar\varepsilon_d$ be the quantity $\varepsilon_{\rm t}$ of Lemma~\ref{lem:decoupling} evaluated with
$\Upsilon(1)\le0.6451\,\varepsilon$ and $\sup|\Xb'|\le\frac{4d}\pi(H_d+1)+1.8884\,d$, let $c_d=0.1935\,\varepsilon$ and
$\varrho_d=2\bar\varepsilon_d/c_d$. Then:
\begin{enumerate}
\item for each $1\le m<d/2$ there are $a_m<b_m$ within $6.6468\cdot10^{-4}\varepsilon^2$ of $\Upsilon_e(m)$ with
$\mathrm{Hc}_m(a_m)<\tau_m<\mathrm{Hc}_m(b_m)$, and $\mathrm{Hc}_m'\ge c_d$ on $[a_m-\varrho_d,b_m+\varrho_d]$;
\item $D_B[\Upsilon_e](1)\ge-2.765192\,\varepsilon^2$ and $D_B[\Upsilon_e](m)\ge6.022431\,\varepsilon^3$ for $2\le m\le d/2$;
\item for every symmetric $\Upsilon$ with $|\Upsilon(m)-\Upsilon_e(m)|\le6.6468\cdot10^{-4}\varepsilon^2+\varrho_d$ for
$1\le m\le d/2$, the one-sided slack $X_\Upsilon$ of the remark after Lemma~\ref{lem:B} satisfies
$\Delta^2X_\Upsilon(m)\ge0.620154\,\varepsilon-4\varrho_dd^2$ for $1\le m\le d-1$;
\item every such $\Upsilon$ has $\Upsilon(1)\le0.6451\,\varepsilon$; consequently, for every $d\ge2001$ and in
natural logarithms, $\log\Pr(E^c)\le-864$, $\log\bar\varepsilon_d\le-847$ and $\log\varrho_d\le-837$ (each bound
decreases in $d$, so $d=2001$ is the worst case); in particular $2\varrho_d$, $4(\bar\varepsilon_d+\varrho_dd^2)$ and
$\varrho_d$ are smaller than $0.1186944\,\varepsilon^2$, $0.620154\,\varepsilon$ and $10^{-6}\varepsilon$, respectively,
by factors larger than $e^{800}$.
\end{enumerate}
\end{proposition}

\emph{What is verified.} The model $\Upsilon_e$ is real-analytic in $(x,\varepsilon)=(m/d,1/d)$, and every statement is
proved by interval Taylor models over boxes in $(x,\varepsilon)$ or $(w,\varepsilon)$, $w=1/m$, with
$\varepsilon\in[0,1/2001]$, so that each box covers all $d\ge2001$ at once (Appendix~\ref{app:model}); the box endpoints lie
on exact decimal grids and the $\varepsilon$-ranges are stored exactly. Item (1) is the \emph{pointwise lemma} (program
\texttt{a10\_pointwise.py}, $159$ boxes; it also certifies that the model root is unique) together with the derivative
bound of \texttt{a12\_fp\_alld.py}: a certified lower bound $0.15204$ for the derivative of the scaled decoupled equation
on all boxes, which gives $\mathrm{Hc}_m'\ge\frac4\pi\,0.15204\,\varepsilon\ge c_d$. Item (2) is assembled from
\texttt{a06\_zoneO.py} (outer zone $x\in[0.011,0.501]$: a fourth-derivative density $\ge7.419$), \texttt{a07\_zoneI.py}
(inner zone $x\le0.013$, $m\ge3$: $\ge6.022431$), \texttt{a08\_A2.py} (the leading inner term $A_2(m)=\Delta^4\varphi_2(m)$
is $\ge0$ for all $m\ge2$, and $A_2(1)=-2.7559456\ldots$) and \texttt{a09\_smallm.py} ($m=1,2$:
$D_B(1)\ge-2.765192\,\varepsilon^2$, $D_B(2)\ge206.884\,\varepsilon^3$); the contribution of the term $A_0\varepsilon f(d-m)$
of $F_A$ is at most $2.03\cdot10^{-11}\varepsilon^3$ (\texttt{a13\_constants.py}). Item (3) is \texttt{a11\_slack.py}
(slack densities of the model $\ge0.651054$ in the outer zone, $\ge0.8260$ in the inner zone, $\ge0.9100$ at $m=1$)
together with a perturbation bound of $0.030900\,\varepsilon$: with $\mathsf G(b,t)=H(b,t)-\Xb(b)$ and $\mathsf G_4$ its expansion to order
$t^4$, the program \texttt{a13\_constants.py} certifies $|\mathsf G-\mathsf G_4|\le138.103\,\varepsilon^3$ (orders $5$--$7$ and the
remainder) and $|\partial_t\mathsf G_4|\le11.5703\,d$, and the Euler--Maclaurin defect of the model target is bounded
analytically. Item (4) follows from the scan of \texttt{a10\_pointwise.py} ($u_e\le0.645042$ for $m=1$),
Lemma~\ref{lem:decoupling} and item (1); every margin is monotone in $d$, so $d=2001$ is the worst case. The program
\texttt{audit\_alld.py} checks the coverage of all boxes exactly, with rational arithmetic, and evaluates every margin in
interval arithmetic from the logs. A final rigour pass (documented with the certificates) completed the written proofs of
the model identities used in Appendix~\ref{app:model}, among them the first-order identity of the inner zone, which was
also checked symbolically and numerically to $6\cdot10^{-25}$, and corrected several code-level enclosure defects; one
of them had left a sliver of size about $10^{-20}$ of the parameter range uncovered at $d=2001$. No conclusion changed.

\begin{proof}[Proof of Theorem~\ref{thm:E} for $d\ge2001$]
Let $\mathrm{lo}_m=a_m-\varrho_d$ and $\mathrm{hi}_m=b_m+\varrho_d$. Since $c_d\varrho_d=2\bar\varepsilon_d$, the remark after
Lemma~\ref{lem:B} and Proposition~\ref{prop:regimeIII}(1),(4) give (B2) and (B3) with $\bar\varepsilon=\bar\varepsilon_d$. For
$\Upsilon$ in the resulting box, $|\Upsilon(m)-\Upsilon_e(m)|\le6.6468\cdot10^{-4}\varepsilon^2+\varrho_d$, so
Lemma~\ref{lem:P}, with $\gamma=2.765192$, $\beta'=6.022431$, $P_d\ge0.962733$ (from $\inf\mathcal A\ge0.1054878997577$)
and $d\ge2001$, gives
\begin{align*}
\hat a_k(\Upsilon)&\ge w_k\Bigl[\frac{\varepsilon^2}8\Bigl(0.962733-\frac{4\gamma^2}{\beta'\cdot2001\cdot(1-4/2001)}\Bigr)
-2\bigl(6.6468\cdot10^{-4}\varepsilon^2+\varrho_d\bigr)\Bigr]\\
&\ge w_k\bigl(0.1186944\,\varepsilon^2-2\varrho_d\bigr)>0,
\end{align*}
which is (B1). Condition (ii) of Lemma~\ref{lem:R1} at the point $\Upsilon^\#$ provided by Lemma~\ref{lem:B} follows from
the remark after Lemma~\ref{lem:B} in its one-sided form, since by Proposition~\ref{prop:regimeIII}(3),(4)
$\Delta^2X_\Upsilon\ge0.620154\,\varepsilon-4\varrho_dd^2>4\bar\varepsilon_d$ on the box. Lemmas~\ref{lem:B}
and~\ref{lem:R1} give Theorem~\ref{thm:E} with $\mu=\mu_{\Upsilon^\#}$, and $\mu(\Delta_d^+)=1$ because, given $\eta$,
$\ell/d$ is a Dirichlet vector with all parameters $\ge\frac14$.
\end{proof}

\subsection{Summary}\label{subsec:conesummary}
Propositions~\ref{prop:regimeI}, \ref{prop:regimeII} and~\ref{prop:regimeIII} prove Theorem~\ref{thm:E} for every $d\ge2$. In
all three regimes $\mu(\Delta_d\setminus\Delta_d^+)=0$: in Regime~I every cell vector has $n^{(k)}_r\ge1$, and in
Regimes~II and~III the cell vector is almost surely in the open simplex. This property is used in the proof of
Theorem~\ref{thm:B}.

All values are certified, and a single program checks all three regimes. It verifies the $199$ certificates of
Regime~I (their exact structure, and a matching verification record of \texttt{verify\_cone.py}), runs
\texttt{audit\_cert\_g.py} on the $1800$ certificates of Regime~II ($1800$ of $1800$ pass, all with the same code hash),
checks every box and every margin of Regime~III, and prints ALL-D CERTIFIED (every $d\ge2$). This program,
\texttt{audit\_alld.py}, and its output are included with the certificates.

%======================================================================
\section{Proofs of Theorems~\ref{thm:A} and~\ref{thm:B}}\label{sec:proofsAB}
%======================================================================

\subsection{Proof of Theorem~\ref{thm:A}}
By Theorem~\ref{thm:E} and Proposition~\ref{prop:conetoopt}, $I_d\le\IME(d)$ for every projective strategy on every
maximally entangled state, and the DKZ strategy attains $\IME(d)$ (Section~\ref{subsec:dkz}).\qed

\begin{remark}[a deficit identity]\label{rem:deficit}
The proof gives more than an inequality. For a family $V$ with Q-configuration $Q$, Lemma~\ref{lem:linearform},
Theorem~\ref{thm:E}, Proposition~\ref{prop:average} and Corollary~\ref{cor:equality} give
\begin{multline*}
\FD-F(V)=\frac{N^2}{2}\int_{\Delta_d}\frac1d\sum_{j=0}^{d-1}\Bigl[\frac1M\int_\T D^{(\ell,j)}_{\lambda^*}\dd m
+\int_\T S^{(\ell,j)}_{\lambda^*}\dd m\Bigr]\mu(\dd\ell)\\
+\frac12\sum_sw_s\bigl[d-\cN(s)-\cN(d-s)\bigr],
\end{multline*}
where $D^{(\ell,j)}_{\lambda^*}$ and $S^{(\ell,j)}_{\lambda^*}$ are the (nonnegative) strip and bathtub defects of
Theorem~\ref{thm:continuum} for the cell embedding of $\mathrm{rot}_jQ$ with cells $\ell$. Thus the CGLMP deficit of any
projective strategy on a maximally entangled state is an explicit positive combination of integrated defects of
Theorem~\ref{thm:C}, of bathtub slacks, and of pair slacks.
\end{remark}

\subsection{Proof of Theorem~\ref{thm:B}}
We need one lemma about cotangent sums.

\begin{lemma}[residue lemma]\label{lem:residue}
Let $t_0,\dots,t_{n-1}\in\R$ be pairwise distinct modulo $2\pi$ and $J_0,\dots,J_{n-1}\in M_M(\C)$. If
$\sum_jJ_j\cot\frac{\theta-t_j}2=0$ for all $\theta$ in a nonempty open interval $I\subseteq\R\setminus\bigcup_j(t_j+2\pi\Z)$, then
$J_0=\dots=J_{n-1}=0$.
\end{lemma}

\begin{proof}
Entrywise, $R(z)=\sum_jJ_j\cot\frac{z-t_j}2$ is holomorphic on the connected open set $\Omega=\C\setminus\bigcup_j(t_j+2\pi\Z)$ and
vanishes on $I$, so $R=0$ on $\Omega$. Near $z=t_j$ the terms with index $j'\ne j$ are holomorphic, and
$\cot\frac{z-t_j}2=\frac2{z-t_j}+O(z-t_j)$; so the residue of $R$ at $t_j$ is $2J_j$, which must vanish.
\end{proof}

\begin{proposition}[pointwise commutation forces commutation]\label{prop:residue}
Let $\ell\in\Delta_d^+$, let $Q$ be a Q-configuration, and let $B$ and $g$ be the step field and its conjugate function
from Proposition~\ref{prop:cellineq}. If $[B(\theta),g(\theta)]=0$ for almost every $\theta$, then $Q_xQ_y=Q_yQ_x$ for all
$x,y\in\Z_N$.
\end{proposition}

\begin{proof}
Since all cells have positive length, the points $t_x=2\pi L_x/N$ ($x\in\Z_N$) are pairwise distinct modulo $2\pi$ and the arcs
$J_x=(t_x,t_{x+1})$ are nonempty. Fix $x$. On $J_x$, $B(\theta)=B_x$ and $g$ is real-analytic by~\eqref{eq:conj}, so
$\theta\mapsto[B_x,g(\theta)]$ is continuous on $J_x$ and vanishes almost everywhere there, hence identically; so does its
derivative,
\[
0=\frac{\dd}{\dd\theta}[B_x,g(\theta)]=\frac1{2\pi}\sum_{j\in\Z_N}[B_x,B_j-B_{j-1}]\cot\frac{\theta-t_j}2\qquad(\theta\in J_x).
\]
By Lemma~\ref{lem:residue}, $[B_x,B_j]=[B_x,B_{j-1}]$ for every $j$; so $j\mapsto[B_x,B_j]$ is constant on $\Z_N$, with value
$[B_x,B_x]=0$. As $x$ runs over $\Z_N$, $B_x=Q_{x+d}$ runs over all the $Q$'s.
\end{proof}

\begin{proof}[Proof of Theorem~\ref{thm:B}]
Let a projective strategy $R$ on $\Phi_D$ attain $\IME(d)$, let $V$ be its reduced family (Theorem~\ref{thm:reduction}) and
$Q$ the Q-configuration of $V$. Then $F(V)=\FD$, so $\ip v\delta=0$ by Lemma~\ref{lem:linearform}.

\emph{Step 1: a tight cell inequality with all cells positive.} By Theorem~\ref{thm:E},
$0=\ip v\delta=\int\ip{u^\ell}\delta\mu(\dd\ell)+\sum_sw_s\ip{t_s}\delta$, and every term is nonpositive
(Proposition~\ref{prop:average} and Lemma~\ref{lem:linearform}(3)). Hence the continuous nonpositive function
$\ell\mapsto\ip{u^\ell}\delta$ vanishes $\mu$-almost everywhere. As $\mu(\Delta_d^+)=\mu(\Delta_d)>0$, there is a cell vector
$\ell^*\in\Delta_d^+$ with $\ip{u^{\ell^*}}\delta=0$. By Proposition~\ref{prop:average}, the $d$ numbers
$Q^{\ell^*}(\mathrm{rot}_jQ)-N^2G/\pi^2$ are nonpositive with mean $0$; so $Q^{\ell^*}(Q)=N^2G/\pi^2$.

\emph{Step 2: commutation almost everywhere.} For the step fields of Proposition~\ref{prop:cellineq} built from $\ell^*$,
$\int\tau(Ag)\dd m=N^{-2}Q^{\ell^*}(Q)=G/\pi^2$; by Corollary~\ref{cor:equality}, $[B(\theta),g(\theta)]=0$ for almost every
$\theta$.

\emph{Step 3: commutation.} By Proposition~\ref{prop:residue} all $Q_x$ commute. Hence the unitaries
$E_k=\sum_ai^aQ_{k-da}$ commute, and so do the $V_k=E_k^*$.

\emph{Step 4: the commutative case.} By Theorem~\ref{thm:commutative}(2),(3), $d\mid D$ and there is a unitary
$u\colon\C^D\to\C^d\otimes\C^{D/d}$ with $uR_iu^*=R_i^{\mathrm{DKZ}}\otimes\one$ for $i=1,\dots,4$.

\emph{Step 5: measurements, state and uniqueness.} From $uU_xu^*=U_x^{\mathrm{DKZ}}\otimes\one$ we get
$uP^x_au^*=P^{x,\mathrm{DKZ}}_a\otimes\one$ (spectral projections). From $u(U'_y)^{\mathsf T}u^*=(U'^{\mathrm{DKZ}}_y)^{\mathsf T}\otimes\one$, by
transposition, $\bar uU'_y\bar u^{\,*}=U'^{\mathrm{DKZ}}_y\otimes\one$, so $\bar uQ^y_b\bar u^{\,*}=Q^{y,\mathrm{DKZ}}_b\otimes\one$. For any
unitary $u$, $(u\otimes\bar u)\sum_i|ii\rangle=\sum_{j,k}(uu^*)_{jk}|jk\rangle=\sum_j|jj\rangle$, so $(u\otimes\bar u)\Phi_D=\Phi_d\otimes\Phi_{D/d}$
in the product basis. If $u_1,u_2$ both have these properties, then $u_2u_1^*$ commutes with all $R^{\mathrm{DKZ}}_i\otimes\one$;
the $R_i^{\mathrm{DKZ}}=GR_i^{(0)}G^*$ generate $M_d(\C)$ (the algebra generated by the $R^{(0)}_i$ contains $X$ and
$R^{(0)}_1(R^{(0)}_2)^{-1}=z^{-1}(\one+(i-1)|0\rangle\langle0|)$, hence all matrix units~\cite[proof of
Corollary~C]{MS-classical}), so $u_2u_1^*\in\one_d\otimes U(D/d)$. Finally, for fixed $D$ the projective strategies on $\Phi_D$ form a
compact set on which $I_d$ is continuous, so the maximum is attained; if it were $\IME(d)$, then $d\mid D$.
\end{proof}

\begin{remark}[why all cells must be positive]\label{rem:positivecells}
If $\ell_r=0$ for some residue $r$, then $K^\ell(x,y)=0$ whenever $x\equiv r$ or $y\equiv r\pmod d$, so $Q^\ell$ does not depend on the
measurement of site $r$, and a tight cell inequality cannot force that measurement to commute with the others. The
uniform cell vector lies in $\Delta_d^+$, but $u^1$ is not proportional to $v$ (Remark~\ref{rem:uniform}), so $\ip v\delta=0$ does
not make the uniform cell inequality tight. Step~1 needs a cone representation charging $\Delta_d^+$, which
Theorem~\ref{thm:E} provides.
\end{remark}

%======================================================================
\section{The continuum limit}\label{sec:consequences}
%======================================================================

\begin{corollary}\label{cor:limit}
As $d\to\infty$,
\[
\IME(d)=\frac{32G}{\pi^2}-\frac{c_1}d+O(d^{-2}),\qquad c_1=2+\frac4\pi-\frac{32G}{\pi^2}=0.30342\ldots,
\]
so the maximum of $I_d$ over projective strategies on maximally entangled states tends to $32G/\pi^2=2.96981\ldots$.
\end{corollary}

\begin{proof}
By Lemma~\ref{lem:linearform}(2), $\FD=\sum_{m=1}^{d-1}m\csc x_m$ with $x_m=\frac{\pi m}{2d}$. Write
$m\csc x_m=\frac{2d}\pi(1+f(x_m))$ with $f(x)=x\csc x-1$, which is smooth on $[0,\frac\pi2]$ with $f(0)=0$, $f(\frac\pi2)=\frac\pi2-1$ and
$\int_0^{\pi/2}f=2G-\frac\pi2$ (because $\int_0^{\pi/2}x\csc x\dd x=2G$). The trapezoidal rule gives
$\sum_{m=1}^{d-1}f(x_m)=\frac{2d}\pi(2G-\frac\pi2)-\frac12(\frac\pi2-1)+O(\frac1d)$, hence
$\FD=\frac{8G}{\pi^2}d^2-(\frac12+\frac1\pi)d+O(1)$ and $\IME(d)=\frac{4\FD}{d(d-1)}$ has the stated expansion.
\end{proof}

In the language of Sections~\ref{sec:continuum} and~\ref{sec:cells}, $N^{-2}(2\FD+d)\to\Phic(\frac14,\frac14)=G/\pi^2$: the
discrete value converges to the value of the continuum problem, which by Theorem~\ref{thm:continuum} is the same for
matrix-valued and for scalar configurations. The continuum theorem alone does not give the exact discrete bound: the
discrete and continuum problems differ by terms of order $d$ (Remark~\ref{rem:uniform}), and the exact bound requires
the cell embeddings and Theorem~\ref{thm:E}.

%======================================================================
\section{The remaining clauses of Problem~27}\label{sec:clauses}
%======================================================================

Besides the clause settled by Theorems~\ref{thm:A} and~\ref{thm:B}, Part~B of Problem~27~\cite{OQP} asks to show that
the DKZ measurements realise the highest resistance of the violation to noise, and the best discrimination against
local realism in the sense of the Kullback--Leibler divergence~\cite{vDGG}. In this section we settle both clauses: the
noise clause in two readings (a third, stronger reading remains open), and the Kullback--Leibler clause for every
$d\ge4$. Throughout, the settings are $x,y\in\{1,2\}$ and the outcomes $a,b\in\Z_d$, as in~\eqref{eq:cglmp}, and
$\DKZ_d$ denotes the DKZ strategy on $\Phi_d$. The statements and proofs, the certificates, the verification programs
and the reports of the independent verifications of this section are in the directory \texttt{complete-resolution/}
of~\cite{repo-oqp27}; its file \texttt{LEDGER.md} records the wording of Problem~27 on the archived page~\cite{OQP} and
the verdict on each clause.

\subsection{Critical visibilities}\label{subsec:noisesetting}
A \emph{behaviour} is a family $p=(p(\cdot,\cdot|x,y))_{x,y}$ of probability distributions on $\Z_d\times\Z_d$. For
$\lambda=(a_1,a_2,b_1,b_2)\in\Z_d^4$ let $\delta_\lambda(a,b|x,y)=[a=a_x][b=b_y]$. The local polytope is
$\cL_d=\operatorname{conv}\{\delta_\lambda:\lambda\in\Z_d^4\}$; a behaviour is \emph{local} if it lies in $\cL_d$, that
is, if it is the behaviour of a random deterministic strategy (a \emph{local model}). Relabellings of the outcomes of
each measurement, of the settings and of the parties, and local post-processing of the outcomes, map $\cL_d$ into
itself. Let $p_0(a,b|x,y)=d^{-2}$ be the behaviour of completely random outcomes; it is local. For a behaviour $p$ and
$0\le v\le1$ put $p_v=vp+(1-v)p_0$ and
\[
\vc(p)=\max\{v\in[0,1]:p_v\in\cL_d\}.
\]
As $\cL_d$ is closed and convex and contains $p_0$, $p_v$ is local for $v\le\vc(p)$ and nonlocal for $v>\vc(p)$; a
\emph{smaller} critical visibility means a \emph{higher} resistance to noise. Relabellings map $\cL_d$ onto itself and
fix $p_0$, so they do not change $\vc$. For a strategy we write $\vc$ of its behaviour. In the paper behind the
problem, Gill measures noise resistance by mixing the behaviour with ``completely random, uniform
outcomes''~\cite[p.~140]{Gill} as long as local realism is still violated, that is, by $\vc$. For strategies on
maximally entangled states the noise clause can be read in three ways.
\begin{enumerate}
\item[(a)] \emph{CGLMP reading}: noise $p_0$, nonlocality detected by the CGLMP inequality $I_d\le2$.
\item[(b)] \emph{Gill's literal reading}: noise $p_0$, nonlocality against all Bell inequalities, that is, $\vc$.
\item[(c)] \emph{White-noise reading}: noise on the state, $v|\Phi_D\rangle\langle\Phi_D|+(1-v)\one/D^2$, and
nonlocality against all Bell inequalities. The noise then contributes the product behaviour
$\Tr P^x_a\,\Tr Q^y_b/D^2$, which equals $p_0$ if all outcomes of all measurements have trace $D/d$, for example for
rank-one measurements on $\Phi_d$.
\end{enumerate}
Reading~(a) holds for every $d$ (Corollary~\ref{cor:noise}), reading~(b) fails for every $d\ge4$
(Theorem~\ref{thm:literal}), and reading~(c) is open for $d\ge4$ (Section~\ref{subsec:white}).

\subsection{The CGLMP reading}\label{subsec:cglmpnoise}
\begin{corollary}\label{cor:noise}
Let $p$ be the behaviour of a projective strategy on $\Phi_D$. Then $I_d(p_v)=vI_d(p)$, so the CGLMP inequality detects
$p_v$ (that is, $I_d(p_v)>2$) if and only if $v>2/I_d(p)$. Consequently the CGLMP threshold $2/I_d(p)$ (defined when
$I_d(p)>2$) satisfies $2/I_d(p)\ge2/\IME(d)$, with equality if and only if the strategy is, up to a local unitary
$u\otimes\bar u$, the DKZ strategy tensored with the identity (in particular $d\mid D$). The same holds for white noise
added to the state, $v\,|\Phi_D\rangle\langle\Phi_D|+(1-v)\one/D^2$, whenever every outcome of every measurement has
$\Tr P^x_a=\Tr Q^y_b=D/d$ (for example rank-one measurements with $D=d$), because then the noise term produces $p_0$.
\end{corollary}

\begin{proof}
$I_d$ is linear in the behaviour, and $I_d(p_0)=0$ because each of the eight probabilities in each bracket
of~\eqref{eq:cglmp} equals $1/d$ under $p_0$. Apply Theorems~\ref{thm:A} and~\ref{thm:B}. Under the trace condition,
$\Tr\bigl((P^x_a\otimes Q^y_b)\one/D^2\bigr)=(D/d)^2/D^2=d^{-2}$.
\end{proof}

Corollary~\ref{cor:noise} is formally verified in Lean with the scope of Theorems~\ref{thm:A} and~\ref{thm:B}: without
assumptions for $2\le d\le20$, and assuming the cone condition (Theorem~\ref{thm:E}) for $d\ge21$
(Section~\ref{subsec:formal}). For measurements with outcomes of unequal traces, white noise on the state produces a
local behaviour whose CGLMP value need not vanish, and Corollary~\ref{cor:noise} says nothing about it. The corollary
concerns the CGLMP inequality only; for all Bell inequalities the corresponding statement is false for every $d\ge4$
(Theorem~\ref{thm:literal}).

\subsection{Gill's literal reading}\label{subsec:literal}
Let $\sigma_z,\sigma_x$ be Pauli matrices, let
\[
\mathsf A_1=\sigma_z,\qquad\mathsf A_2=\sigma_x,\qquad\mathsf B_1=\tfrac1{\sqrt2}(\sigma_z+\sigma_x),\qquad
\mathsf B_2=\tfrac1{\sqrt2}(\sigma_z-\sigma_x)
\]
be the optimal CHSH observables, and let $\Pi^x_\pm=\frac12(\one\pm\mathsf A_x)$ and $\Pi'^y_\pm=\frac12(\one\pm\mathsf B_y)$,
real symmetric projections of rank one on $\C^2$. For $d\ge3$ define a projective strategy $C_d$ on $\Phi_d$ as
follows.
\begin{itemize}
\item \emph{$d$ even.} Identify $\C^d=\C^2\otimes\C^{d/2}$ by $|i\rangle|j\rangle\mapsto|i\tfrac d2+j\rangle$, so that
$\Phi_d=\Phi_2\otimes\Phi_{d/2}$, and let $P^x_0=\Pi^x_+\otimes\one$, $P^x_1=\Pi^x_-\otimes\one$,
$Q^y_0=\Pi'^y_+\otimes\one$ and $Q^y_1=\Pi'^y_-\otimes\one$.
\item \emph{$d$ odd.} Decompose
$\C^d=\bigoplus_{k=1}^{(d-1)/2}\operatorname{span}\{|2k-2\rangle,|2k-1\rangle\}\oplus\C|d-1\rangle$ and let
$P^x_0=\bigl(\bigoplus_k\Pi^x_+\bigr)\oplus1$, $P^x_1=\bigl(\bigoplus_k\Pi^x_-\bigr)\oplus0$, and $Q^y_0$, $Q^y_1$
likewise with $\Pi'^y_\pm$.
\end{itemize}
In both cases $P^x_a=Q^y_b=0$ for $a,b\ge2$: the outcomes $2,\dots,d-1$ never occur. Put
\[
\veven(d)=\frac{4(d-1)}{(\sqrt2-1)d^2+4(d-1)},\qquad\vodd(d)=\frac4{(\sqrt2-1)d+4}.
\]

\begin{theorem}[Gill's noise clause; computer-assisted for $4\le d\le9$]\label{thm:literal}
\begin{enumerate}
\item For every $d\ge3$, $\vc(C_d)=\veven(d)$ if $d$ is even and $\vc(C_d)=\vodd(d)$ if $d$ is odd.
\item For every $d\ge2$, $\frac12\le\vc(\DKZ_d)\le2/\IME(d)$. For $4\le d\le9$, $\vc(\DKZ_d)\ge r(d)$, where
$r(4),\dots,r(9)=0.69,\ 0.687,\ 0.684,\ 0.683,\ 0.682,\ 0.681$.
\item For every $d\ge4$, $\vc(C_d)<\vc(\DKZ_d)$. Hence, on $\Phi_d$ and on every $\Phi_{kd}$, the DKZ measurements do
not realise the highest resistance of the violation of local realism to uniform outcome noise.
\end{enumerate}
\end{theorem}

The only computer-assisted step is Proposition~\ref{prop:dkzcert} below, which is used for $4\le d\le9$. We first
compute the behaviour of $C_d$ and a Bell inequality that detects it.

\begin{lemma}\label{lem:Cbeh}
Let $s_{xy}=-1$ if $x=y=2$ and $s_{xy}=1$ otherwise, let $\CH(a,b|x,y)=\frac14\bigl(1+(-1)^{a+b}s_{xy}/\sqrt2\bigr)$ for
$a,b\in\{0,1\}$ and $\CH(a,b|x,y)=0$ otherwise, and let $\delta_{00}(a,b|x,y)=[a=b=0]$. The behaviour of $C_d$ is $\CH$
if $d$ is even and $\frac{d-1}d\CH+\frac1d\delta_{00}$ if $d$ is odd.
\end{lemma}

\begin{proof}
All projections are real symmetric, so the behaviour is $\Tr(P^x_aQ^y_b)/d$. For $e,f\in\{\pm1\}$,
$\frac12\Tr\bigl(\frac{\one+e\mathsf A_x}2\cdot\frac{\one+f\mathsf B_y}2\bigr)=\frac14\bigl(1+\frac{ef}2\Tr(\mathsf A_x\mathsf B_y)\bigr)$
and $\frac12\Tr(\mathsf A_x\mathsf B_y)=s_{xy}/\sqrt2$. For even $d$,
$\Tr\bigl((\Pi\otimes\one)(\Pi'\otimes\one)\bigr)/d=\frac12\Tr(\Pi\Pi')$. For odd $d$ the trace splits into the
$(d-1)/2$ blocks, each contributing $\Tr(\Pi\Pi')/d$, and the last dimension, which contributes $[a=b=0]/d$.
\end{proof}

Let $\varsigma(0)=1$ and $\varsigma(a)=-1$ for $a\ne0$, let $E_{xy}(p)=\sum_{a,b}\varsigma(a)\varsigma(b)p(a,b|x,y)$, and
let $\mathcal E(p)=E_{11}(p)+E_{12}(p)+E_{21}(p)-E_{22}(p)$: the CHSH expression after coarse-graining every measurement
into the outcome $0$ and the rest.

\begin{lemma}[a coarse-grained CHSH witness]\label{lem:witness}
$\mathcal E\le2$ on $\cL_d$; $\mathcal E(p_0)=2(d-2)^2/d^2$, $\mathcal E(\CH)=2\sqrt2$ and $\mathcal E(\delta_{00})=2$.
Consequently the behaviour $p_v$ of $C_d$ violates $\mathcal E\le2$ if and only if $v>\veven(d)$ ($d$ even),
respectively $v>\vodd(d)$ ($d$ odd).
\end{lemma}

\begin{proof}
On $\delta_\lambda$, $\mathcal E=\varsigma(a_1)\bigl(\varsigma(b_1)+\varsigma(b_2)\bigr)+\varsigma(a_2)\bigl(\varsigma(b_1)-\varsigma(b_2)\bigr)=\pm2$,
and $\mathcal E$ is linear. Under $p_0$ the two coarse-grained outcomes are independent with mean $(2-d)/d$, so
$E_{xy}(p_0)=(d-2)^2/d^2$; under $\CH$, $E_{xy}=s_{xy}/\sqrt2$, and under $\delta_{00}$, $E_{xy}=1$. As
$\mathcal E(p_v)=v\mathcal E(p)+(1-v)\mathcal E(p_0)$ and $\mathcal E(p)>2>\mathcal E(p_0)$, the inequality is violated if
and only if $v>(2-\mathcal E(p_0))/(\mathcal E(p)-\mathcal E(p_0))$. Here $2-\mathcal E(p_0)=8(d-1)/d^2$, and
$\mathcal E(p)-\mathcal E(p_0)$ equals $(2\sqrt2d^2-2(d-2)^2)/d^2$ for $p=\CH$ and $2(d-1)\bigl((\sqrt2-1)d+4\bigr)/d^2$
for $p=\frac{d-1}d\CH+\frac1d\delta_{00}$; the quotients are $\veven(d)$ and $\vodd(d)$.
\end{proof}

The witness is optimal: at the threshold the behaviour is local. We give an explicit local model.

\begin{proposition}[a local model at the threshold]\label{prop:flagmodel}
Let $d\ge3$, $s=1/d$ and $T=4s(1-s)/\bigl(\sqrt2-(1-2s)^2\bigr)$. Then $T=\veven(d)$, and $T\CH+(1-T)p_0$ is local.
\end{proposition}

\begin{proof}
Multiplying numerator and denominator of $T$ by $d^2$ gives $\veven(d)$. Put $c_1=(1-T)(1-2s)$, $c_2=(1-T)(1-2s)^2$ and
$c_0=1-4c_1+3c_2$. Since $1-T=(\sqrt2-1)/\bigl(\sqrt2-(1-2s)^2\bigr)$, one has $c_1-c_2=2s(1-T)(1-2s)\ge0$ and
\[
c_0=\frac{4s\bigl[(1-s)-(\sqrt2-1)(1-3s)\bigr]}{\sqrt2-(1-2s)^2}>0,
\]
because $(\sqrt2-1)(1-3s)<1-3s\le1-s$. Let $\Lambda=(\Lambda_{A_1},\Lambda_{A_2},\Lambda_{B_1},\Lambda_{B_2})$ be a random
variable with values in $\{0,1,\ast\}^4$ and the following law:
\begin{enumerate}
\item with probability $c_2$, $\Lambda=(\ast,\ast,\ast,\ast)$;
\item for each $r\in\{0,1\}$, with probability $\frac12(c_1-c_2)$ each, $\Lambda=(\ast,r,r,1-r)$, $(r,\ast,r,r)$,
$(r,1-r,\ast,r)$ or $(r,r,r,\ast)$;
\item with probability $\frac18c_0$ each, $\Lambda$ is one of the eight $\lambda\in\{0,1\}^4$ with
$\sum_{x,y}s_{xy}(-1)^{\lambda_{A_x}+\lambda_{B_y}}=2$.
\end{enumerate}
The probabilities sum to $c_2+4(c_1-c_2)+c_0=1$. Alice outputs $\Lambda_{A_x}$ if it lies in $\{0,1\}$, and otherwise an
independent uniformly random element of $\{2,\dots,d-1\}$; Bob does the same with $\Lambda_{B_y}$. This is a local
model; we show that its behaviour $q$ is $T\CH+(1-T)p_0$. Fix $(x,y)$, and recall $1-2s=(d-2)/d$. If $a,b\ge2$, only
(1) contributes, and $q(a,b|x,y)=c_2/(d-2)^2=(1-T)/d^2$. If $a\ge2>b$, only the pattern of (2) with $\ast$ at $A_x$ and
the value $b$ at $B_y$ contributes, and $q(a,b|x,y)=\frac12(c_1-c_2)/(d-2)=(1-T)/d^2$; the case $b\ge2>a$ is
symmetric. Let $a,b\in\{0,1\}$. In each pattern of (2) whose $\ast$ is not at $A_x$ or $B_y$, one has
$(-1)^{\Lambda_{A_x}+\Lambda_{B_y}}=s_{xy}$, and the two values of $r$ give the two pairs $(a,b)$ with
$(-1)^{a+b}=s_{xy}$. Each strategy in (3) satisfies $(-1)^{\lambda_{A_x}+\lambda_{B_y}}=s_{xy}$ for exactly three of
the four pairs $(x,y)$, and the set of these eight strategies is invariant under $\lambda\mapsto1-\lambda$; so for each
$(x,y)$, six of them give a pair $(a,b)$ with $(-1)^{a+b}=s_{xy}$ (three for each such pair) and two give one of the
other two pairs (one each). Hence, for $a,b\in\{0,1\}$,
\begin{align*}
q(a,b|x,y)&=\tfrac14\bigl(1-2c_1+c_2\bigr)+\tfrac18(-1)^{a+b}s_{xy}(1-c_2),\\
T\CH(a,b|x,y)+(1-T)d^{-2}&=\tfrac14\bigl(T+4(1-T)s^2\bigr)+\tfrac{T}{4\sqrt2}(-1)^{a+b}s_{xy}.
\end{align*}
The constant terms agree, because $1-2c_1+c_2=1-(1-T)(1-2s)(1+2s)=T+4(1-T)s^2$, and so do the other terms, because
$1-c_2=\sqrt2\,T$ is equivalent to the definition of $T$.
\end{proof}

\begin{proof}[Proof of Theorem~\ref{thm:literal}(1)]
By Lemmas~\ref{lem:Cbeh} and~\ref{lem:witness}, $\vc(C_d)\le\veven(d)$ for even $d$ and $\vc(C_d)\le\vodd(d)$ for odd
$d$. For even $d$, Proposition~\ref{prop:flagmodel} gives equality. For odd $d$ let $s=1/d$; then
\[
vp_{C_d}+(1-v)p_0=vs\,\delta_{00}+(1-vs)\bigl(v'\CH+(1-v')p_0\bigr),\qquad v'=\frac{v(1-s)}{1-vs},
\]
and for $v=\vodd(d)=4s/(\sqrt2-1+4s)$ one finds $1-vs=\bigl(\sqrt2-(1-2s)^2\bigr)/(\sqrt2-1+4s)$ and hence $v'=T$. So
the behaviour at $v=\vodd(d)$ is a mixture of the deterministic behaviour $\delta_{00}$ and the local behaviour of
Proposition~\ref{prop:flagmodel}.
\end{proof}

Thus the coarse-grained CHSH inequality $\mathcal E\le2$ is an optimal witness for $C_d$. For the DKZ strategy we need
lower bounds on $\vc$, that is, local models. The first one is universal.

\begin{lemma}[the $\frac12$-lemma]\label{lem:half}
Let $p$ be a behaviour all of whose marginals are uniform: $\sum_bp(a,b|x,y)=\sum_bp(b,a|x,y)=1/d$ for all $a,x,y$.
Then $p_v$ is local for $0\le v\le\frac12$. The constant $\frac12$ is optimal: the behaviour
$\mathrm{PR}_d(a,b|x,y)=\frac1d[\,b-a\equiv[x=2][y=1]\pmod d\,]$, a $d$-outcome version of the box of Popescu and
Rohrlich~\cite{PR}, has uniform marginals and $\vc(\mathrm{PR}_d)=\frac12$.
\end{lemma}

\begin{proof}
The probability distributions on $\Z_d^4$ given by
\[
P_1(a_1,a_2,b_1,b_2)=p(a_1,b_1|1,1)\,p(a_2,b_2|2,2),\qquad P_2(a_1,a_2,b_1,b_2)=p(a_1,b_2|1,2)\,p(a_2,b_1|2,1)
\]
are local models. Under $P_1$ the pairs of outcomes
for the settings $(1,1)$ and $(2,2)$ have the laws $p(\cdot,\cdot|1,1)$ and $p(\cdot,\cdot|2,2)$, while for $(1,2)$
and $(2,1)$ the two outcomes are independent and uniformly distributed, with law $p_0(\cdot,\cdot|x,y)$; under $P_2$ the
roles are exchanged. Hence $\frac12(P_1+P_2)$ is a local model of $\frac12(p+p_0)$, and for $v\le\frac12$,
$p_v=2v\cdot\frac12(p+p_0)+(1-2v)p_0$ is local. For $\mathrm{PR}_d$ each of the four terms with $k=0$ in the first
bracket of~\eqref{eq:cglmp} equals $1$ and all other probabilities in~\eqref{eq:cglmp} vanish, so
$I_d(\mathrm{PR}_d)=4$; as $I_d(p_0)=0$, $I_d$ detects $(\mathrm{PR}_d)_v$ for every $v>\frac12$.
\end{proof}

The DKZ strategy has uniform marginals, so $\vc(\DKZ_d)\ge\frac12$. With the bases of Section~\ref{subsec:dkz},
$\langle\Phi_d|(|a\rangle_{A,x}\otimes|b\rangle_{B,y})=d^{-3/2}\sum_je^{2\pi ij(b-a-\alpha_x-\beta_y)/d}$, and
$|\sum_{j=0}^{d-1}e^{2\pi ij\theta/d}|^2=\sin^2(\pi\theta)/\sin^2(\pi\theta/d)$ with $\sin^2(\pi(\alpha_x+\beta_y))=\frac12$
give the closed form
\begin{equation}\label{eq:dkzclosed}
p_{\DKZ}(a,b|x,y)=\frac1dh_{xy}(b-a),\qquad h_{xy}(m)=\frac1{2d^2\sin^2\bigl(\pi(m-\delta_{xy})/d\bigr)},
\end{equation}
where $\delta_{xy}=\alpha_x+\beta_y$, that is, $\delta_{11}=\delta_{22}=\frac14$, $\delta_{12}=-\frac14$ and
$\delta_{21}=\frac34$; each $h_{xy}$ is a probability distribution on $\Z_d$.

\begin{lemma}[covariant local models]\label{lem:covariant}
Let $\mu$ be a probability distribution on $M_d=\{m\in\Z_d^{\{1,2\}^2}:m_{11}-m_{12}-m_{21}+m_{22}\equiv0\}$, with
one-dimensional marginals $\mu_{xy}$. Then $p(a,b|x,y)=\frac1d\mu_{xy}(b-a)$ is local.
\end{lemma}

\begin{proof}
Draw $m\sim\mu$ and, independently, $c$ uniformly from $\Z_d$, and output $A_1=c$, $A_2=c+m_{11}-m_{21}$,
$B_1=c+m_{11}$ and $B_2=c+m_{12}$. Then $B_y-A_x=m_{xy}$ for all $x,y$ (for $(2,2)$ by the constraint), and each $A_x$
is uniformly distributed and independent of $m$; so $\Pr(A_x=a,B_y=b)=\frac1d\mu_{xy}(b-a)$.
\end{proof}

\begin{lemma}[absorption]\label{lem:absorb}
Let $p$ have uniform marginals, let $Q$ be a local behaviour with uniform marginals, and let $0<v_0<v_1\le1$. If
$|Q(a,b|x,y)-p_{v_1}(a,b|x,y)|\le(v_1-v_0)/(2d^2v_0)$ for all $a,b,x,y$, then $p_{v_0}$ is local.
\end{lemma}

\begin{proof}
Let $E=p_{v_1}-Q$, whose marginals vanish, and $\gamma=v_0/(v_1-v_0)$. Then
\[
p_{v_0}=\frac{v_0}{v_1}p_{v_1}+\Bigl(1-\frac{v_0}{v_1}\Bigr)p_0=\frac{v_0}{v_1}Q+\Bigl(1-\frac{v_0}{v_1}\Bigr)\cdot
\tfrac12(p_0+S),\qquad S=p_0+2\gamma E.
\]
The hypothesis gives $S\ge0$; so $S$ is a behaviour with uniform marginals, and $\frac12(p_0+S)$ is local by
Lemma~\ref{lem:half}.
\end{proof}

For $p=p_{\DKZ}$ and $Q=\frac1d\mu_{xy}(b-a)$ as in Lemma~\ref{lem:covariant}, the condition of
Lemma~\ref{lem:absorb} reads $|\mu_{xy}(m)-v_1h_{xy}(m)-(1-v_1)/d|\le(v_1-v_0)/(2dv_0)$ for all $x,y,m$.

\begin{proposition}[computer-assisted]\label{prop:dkzcert}
For every $3\le d\le20$ there are rationals $v_0(d)<v_1(d)$ and a probability distribution $\mu$ on $M_d$ with rational
weights (common denominator $10^{16}$) whose marginals satisfy this condition. Hence $\vc(\DKZ_d)\ge v_0(d)$, where
$2/\IME(d)-1.02\cdot10^{-7}<v_0(d)<2/\IME(d)$; in particular $\vc(\DKZ_d)\ge r(d)$ for $4\le d\le9$. For example,
$v_0(4)=0.690549638999$ and $v_0(9)=0.681080610999$.
\end{proposition}

\begin{proof}
The weights are checked to be nonnegative with sum $1$ in integer arithmetic, and the marginals $\mu_{xy}(m)$ are
computed exactly. The numbers $h_{xy}(m)$ are enclosed in rational intervals: $\pi$ is enclosed by consecutive partial
sums of Machin's formula $\pi=16\arctan\frac15-4\arctan\frac1{239}$ (alternating series, enclosure width below
$10^{-50}$); $\sin^2(\pi s)$ for rational $s$ is reduced to $\sin(\pi r)$ with $r\in[0,\frac12]$, which lies between
$\sin(r\pi_-)$ and $\sin(r\pi_+)$ for an enclosure $[\pi_-,\pi_+]$ of $\pi$; and $\sin z$ for rational $0\le z\le2$ is
enclosed by consecutive partial sums of its Taylor series, which alternates with decreasing terms because $z^2<6$. The
largest deviation $|\mu_{xy}(m)-v_1h_{xy}(m)-(1-v_1)/d|$ over the enclosures is then compared exactly with
$(v_1-v_0)/(2dv_0)$; the margins are about $10^{-13}$ and the enclosure widths about $10^{-15}$. The distributions
$\mu$ were found by a linear program over the covariant local models ($d^3$ variables), solved in floating point
with HiGHS~\cite{HiGHS} slightly below its optimum and then rounded; this computation plays no role in the verification.
The certificates (\texttt{certificates/dkz\_d3.json}, \dots, \texttt{dkz\_d20.json}) and the checking program
\texttt{dkz\_cert.py} are in \texttt{complete-resolution/noise-literal/} of~\cite{repo-oqp27}.
\end{proof}

\begin{proof}[Proof of Theorem~\ref{thm:literal}(2) and (3)]
(2) The lower bound $\frac12$ is Lemma~\ref{lem:half}. Since $I_d(p_{\DKZ})=\IME(d)$ (Section~\ref{subsec:dkz}) and
$I_d(p_0)=0$, $I_d$ detects $(p_{\DKZ})_v$ for $v>2/\IME(d)$, so $\vc(\DKZ_d)\le2/\IME(d)$. The bounds $r(d)$ follow from
Proposition~\ref{prop:dkzcert}, since $v_0(d)\ge r(d)$ for $4\le d\le9$.

(3) For $d\ge10$ we show $\vc(C_d)<\frac12$. For even $d$, $\veven(d)<\frac12$ if and only if $(\sqrt2-1)d^2-4d+4>0$;
this holds at $d=9$ ($81\sqrt2>113$) and hence for all $d\ge9$, the derivative $2(\sqrt2-1)d-4$ being positive for
$d\ge5$. For odd $d$, $\vodd(d)<\frac12$ if and only if $d>4(\sqrt2+1)=9.657\ldots$. For $4\le d\le9$ we have
$\vc(C_d)<r(d)$, by exact comparison in $\Q(\sqrt2)$; numerically $\vc(C_d)=0.6442$, $0.6589$, $0.5729$, $0.5798$,
$0.5137$, $0.5176$ for $d=4,\dots,9$, against $r(d)=0.69,\dots,0.681$. (For $d=8,9$ the value exceeds $\frac12$, which is
why certificates are needed up to $d=9$.) Finally, $C_d\otimes\one_k$ and $\DKZ_d\otimes\one_k$ on
$\Phi_{kd}=\Phi_d\otimes\Phi_k$ have the behaviours of $C_d$ and $\DKZ_d$.
\end{proof}

\begin{remark}[exact values for the DKZ strategy]\label{rem:exactdkz}
For $3\le d\le20$ the lower bounds of Proposition~\ref{prop:dkzcert} are within $1.02\cdot10^{-7}$ of the upper bound
$2/\IME(d)$: up to this accuracy the CGLMP inequality is an optimal witness for $\DKZ_d$ against all Bell inequalities.
For $d=3,4,5,6$ the equality $\vc(\DKZ_d)=2/\IME(d)$ holds exactly. The link laws of the covariant local behaviours of
Lemma~\ref{lem:covariant} form the polytope spanned by the vectors $(e_{m_{11}},e_{m_{12}},e_{m_{21}},e_{m_{22}})$,
$m\in M_d$; its complete facet lists ($66$, $216$, $1020$ and $2462$ facets for $d=3,4,5,6$) were computed by an exact
double description, and in exact arithmetic in $\Q(\cos\frac\pi{4d})$ the link laws of $(p_{\DKZ})_v$ at
$v=2/\IME(d)$ satisfy all of them (computer-assisted; not needed for Theorem~\ref{thm:literal}).
\end{remark}

\begin{proposition}[$d=2$ and $d=3$]\label{prop:small}
\begin{enumerate}
\item For every quantum behaviour $p$ with $d=2$ (any state, any measurements), $\vc(p)\ge1/\sqrt2=\vc(\DKZ_2)$.
\item (computer-assisted) For every projective strategy on $\Phi_3$, $\vc\ge2/\IME(3)=3\sqrt3-\frac92=0.696152\ldots$,
with equality if and only if the strategy is $\DKZ_3$ up to relabellings and a local unitary $u\otimes\bar u$.
\item The strategy on $\Phi_2$ with the three-outcome measurements $\{\Pi^x_+,\Pi^x_-,0\}$ and $\{\Pi'^y_+,\Pi'^y_-,0\}$
has $\vc=8/(9\sqrt2-1)=0.682132\ldots<\vc(\DKZ_3)$. The same behaviour arises on every $\Phi_D$ with $D$ even, in
particular on $\Phi_6$, which also carries $\DKZ_3\otimes\one_2$.
\end{enumerate}
So the clause holds for $d=2$; for $d=3$ it holds on $\Phi_3$ and fails if every maximally entangled state is allowed.
\end{proposition}

\begin{proof}
(1) The local polytope $\cL_2$ is described by positivity and the eight CHSH inequalities~\cite{Fine}. Every CHSH
expression vanishes on $p_0$ and is at most $2\sqrt2$ on quantum behaviours~\cite{Tsirelson}, so $p_v$ satisfies all
of them for $v\le1/\sqrt2$; and $I_2(\DKZ_2)=2\sqrt2$ gives $\vc(\DKZ_2)\le1/\sqrt2$.

(2) The polytope $\cL_3$ has exactly $1116$ facets: $36$ positivity facets, $648$ liftings of CHSH inequalities (every
measurement coarse-grained into one outcome against the other two) and the $432$ images of the CGLMP inequality under
relabellings. This is the classification of Collins and Gisin~\cite{CollinsGisin}; we recomputed the list by an exact
double description, verifying every facet. Let $p$ be the behaviour and $v\le2/\IME(3)$. Positivity holds. A facet of
CGLMP type is a positive multiple of $2-I_3(g\cdot p)$ for a relabelling $g$; $g$ maps projective strategies on $\Phi_3$
to projective strategies on $\Phi_3$ (for the exchange of the parties use $\Tr(X^{\mathsf T}Y)=\Tr(Y^{\mathsf T}X)$) and
fixes $p_0$, so Theorem~\ref{thm:A} gives $I_3(g\cdot p_v)=vI_3(g\cdot p)\le v\IME(3)\le2$. For a lifted CHSH facet let
$\tilde{\mathsf A}_x,\tilde{\mathsf B}_y$ be the Hermitian involutions of the coarse-grained measurements. Up to
relabelling, the CHSH value of $p$ is
$\frac13[\Tr(\tilde{\mathsf A}_1^{\mathsf T}(\tilde{\mathsf B}_1+\tilde{\mathsf B}_2))+\Tr(\tilde{\mathsf A}_2^{\mathsf T}(\tilde{\mathsf B}_1-\tilde{\mathsf B}_2))]$,
which by H\"older's inequality is at most
$\frac13(\|\tilde{\mathsf B}_1+\tilde{\mathsf B}_2\|_1+\|\tilde{\mathsf B}_1-\tilde{\mathsf B}_2\|_1)$. By Jordan's
lemma $\C^3$ splits into subspaces of dimension at most two invariant under $\tilde{\mathsf B}_1$ and
$\tilde{\mathsf B}_2$; a two-dimensional one contributes at most $4\sqrt2$ and a one-dimensional one at most $2$, so the
CHSH value of $p$ is at most $(4\sqrt2+2)/3$, while that of $p_0$ is $\pm\frac29$. Hence no lifted CHSH facet is
violated for $v\le4/(3\sqrt2+1)=0.762974\ldots$, which exceeds $2/\IME(3)$. So $\vc\ge2/\IME(3)$, with equality for
$\DKZ_3$ by Theorem~\ref{thm:literal}(2). If $\vc=2/\IME(3)$, then for every $v>2/\IME(3)$ some facet of CGLMP type is
violated by $p_v$; as there are finitely many, one relabelling $g$ works for a sequence $v\downarrow2/\IME(3)$, so
$I_3(g\cdot p)\ge\IME(3)$, and Theorem~\ref{thm:B} applies to the relabelled strategy.

(3) As in Lemma~\ref{lem:Cbeh}, the behaviour is $\CH$, now with $d=3$. Lemma~\ref{lem:witness} and
Proposition~\ref{prop:flagmodel} for $d=3$ give $\vc=\veven(3)=8/(9\sqrt2-1)$, which is smaller than $3\sqrt3-\frac92$.
On $\Phi_D$ with $D$ even, the measurements $\Pi^x_\pm\otimes\one_{D/2}$, $\Pi'^y_\pm\otimes\one_{D/2}$ give the same
behaviour.
\end{proof}

\begin{remark}[unused outcomes]\label{rem:scope}
The strategy $C_d$ never produces the outcomes $2,\dots,d-1$, while the noise $p_0$ gives each of them probability
$1/d$; this is what lets a two-outcome CHSH structure beat DKZ, and $\vc(C_d)\sim4/((\sqrt2-1)d)\to0$. On $\Phi_d$,
projective measurements cannot avoid this: a PVM on $\C^d$ all of whose $d$ outcomes occur consists of $d$ nonzero
orthogonal projections, so it is rank-one, and for rank-one measurements on $\Phi_d$ the noise $p_0$ is white noise on
the state (reading~(c)). A projective counterexample on $\Phi_d$ that uses every outcome would therefore settle
reading~(c). With POVMs every outcome can be used. For $0<\varepsilon\le\frac1{100}$ the effects
$M^x_a=(1-\varepsilon)P^x_a+\frac\varepsilon d\one$ and $N^y_b=(1-\varepsilon)Q^y_b+\frac\varepsilon d\one$, with
$(P,Q)=C_d$, are invertible, every pair of outcomes has probability at least $\varepsilon^2/d^2$, and still
$\vc(M,N)<\vc(\DKZ_d)$ for every $d\ge4$. Indeed, the behaviour of $(M,N)$ is obtained from that of $C_d$ by the local
channel that keeps each outcome with probability $1-\varepsilon$ and otherwise replaces it by a uniformly random one;
this channel fixes $p_0$ and preserves locality, so $\vc(M,N)\ge\vc(C_d)$, and the witness $\mathcal E$ gives
$\vc(M,N)\le w_\varepsilon(d)=(2-2e^2)/(\mathcal E_\varepsilon-2e^2)$, where $e=(2-d)/d$,
$\mathcal E_\varepsilon=2\sqrt2(1-\varepsilon)^2+2\varepsilon^2e^2$ for even $d$, and
$\mathcal E_\varepsilon=(1-\varepsilon)^2\bigl(2\sqrt2\frac{d-1}d+\frac2d\bigr)+4\varepsilon(1-\varepsilon)\frac ed+2\varepsilon^2e^2$
for odd $d$. Exact computations show $w_\varepsilon(d)\le w_{1/100}(d)<v_0(d)$ for $4\le d\le10$ (for example
$w_{1/100}(4)=0.6602$ and $w_{1/100}(9)=0.5380$) and $w_{1/100}(d)<\frac12$ for $d\ge11$. Projective measurements on
$\Phi_D$ with $D>d$ can also use every outcome and still beat DKZ: one adds $d-2$ extra basis vectors carrying the
outcomes $2,\dots,d-1$; for $(d,D)=(4,12)$, $(5,11)$, $(6,12)$, $(7,13)$, $(8,12)$, $(9,13)$ the resulting critical
visibility is below $v_0(d)$, and for $(4,12)$ the strategy $\DKZ_4\otimes\one_3$ lives on the same state.
\end{remark}

\begin{remark}[prior work]\label{rem:literalcredit}
The mechanism behind $C_d$ is known. Ac\'{\i}n, Durt, Gisin and Latorre~\cite[eq.~(14)]{ADGL} found that the CHSH
strategy on a maximally entangled state of Schmidt rank two embedded in $\C^d\otimes\C^d$, with two-outcome
measurements and white noise, has critical visibility $\veven(d)$, which tends to $0$, and concluded that resistance to
noise is not a good measure of nonlocality; for even $d$ the strategy $C_d$ produces the same coarse-grained statistics
on $\Phi_d$ itself. In the two-outcome setting the exactness of this threshold is elementary; in the $d$-outcome
scenario it is Proposition~\ref{prop:flagmodel}. Gruca, Laskowski and \.Zukowski~\cite{GLZ} found numerically, for
$d\le5$, that states of Schmidt rank two are the most noise-robust. Baek, Ryu and Lee~\cite{BRL} observed by linear
programming that merging outcomes increases the white-noise robustness $R_w=(1-\vc)/\vc$ of the DKZ behaviour on
$\Phi_4$ and $\Phi_{16}$, from $0.448$ to $0.552$ for $d=4$, and along the chain $0.4755$, $0.6216$, $1.0243$, $1.7673$
for $d=16$; these values agree with $(\IME(4)-2)/2=0.44812$, $(1-\veven(4))/\veven(4)=0.55229$ and
$(1-\veven(16))/\veven(16)=1.76731$, so merged-outcome strategies beating DKZ on $\Phi_4$ and $\Phi_{16}$ under uniform
noise were already visible there; they did not relate this to Problem~27. The thresholds $2/\IME(d)$ of the DKZ
measurements are the numerical values of~\cite{KGZMZ,DKZ,CGLMP,GLZ}, and Chen et al.~\cite{Chen01} gave an analytic
result for $d=3$; we did not check whether it contains a local model, so the novelty of our lower bound for $d=3$ is
uncertain. As far as we found, the following are new: an exact theorem for every $d\ge4$ on $\Phi_d$ itself, for both
parities, the exact critical visibility of $C_d$ in the $d$-outcome scenario, rigorous lower bounds for
$\vc(\DKZ_d)$, and the consequence for Problem~27.
\end{remark}

\subsection{The white-noise reading}\label{subsec:white}
In reading~(c), the noise behaviour of a projective strategy on $\Phi_D$ is $n(a,b|x,y)=\Tr P^x_a\,\Tr Q^y_b/D^2$, the
product of the marginals of its behaviour $p$, and the question is whether $vp+(1-v)n$ is local for all
$v\le2/\IME(d)$, for every projective strategy on every $\Phi_D$, with equality of the critical visibility only for
DKZ (up to $u\otimes\bar u$ and an inert identity). For strategies with $n=p_0$ this is reading~(b), and on $\Phi_d$
these are exactly the strategies that use every outcome (Remark~\ref{rem:scope}), which $C_d$ does not; for them the
answer is positive for $d=2$ and, on $\Phi_3$, for $d=3$ (Proposition~\ref{prop:small}). For $d\ge4$ the question is
open. Numerical searches have found no strategy more robust than DKZ for $d\le8$ (for $d\le5$ see
also~\cite{GLZ}; our searches are summarised in \texttt{complete-resolution/LEDGER.md} of~\cite{repo-oqp27}). A proof
has to control all facets of $\cL_d$, not only the CGLMP inequality: already for $d=4$ the polytope of link laws of
covariant local behaviours (Lemma~\ref{lem:covariant}), Masanes' polytope of correlation functions~\cite{Masanes}, has
facets that are not images of the CGLMP inequality.

\subsection{The Kullback--Leibler clause}\label{subsec:kl}
Van Dam, Gill and Gr\"unwald~\cite{vDGG} measure the strength of a nonlocality proof with behaviour $q$ by its
Kullback--Leibler divergence from the local behaviours. For a probability distribution $\sigma$ on the setting pairs
$\{1,2\}^2$ let
\[
U_q(\sigma)=\inf_{p\in\cL_d}\sum_{x,y}\sigma_{xy}\,D\bigl(q(\cdot,\cdot|x,y)\,\big\|\,p(\cdot,\cdot|x,y)\bigr),\qquad
D(r\|s)=\sum_ir_i\log_2\frac{r_i}{s_i},
\]
and let $S^{\mathrm{UNI}}(q)=U_q(\frac14)$ (uniform settings), $S^{\mathrm{UC}}(q)$ the supremum of $U_q(\sigma)$ over
product distributions $\sigma=\sigma_A\otimes\sigma_B$, and $S^{\mathrm{COR}}(q)=\sup_\sigma U_q(\sigma)$; so
$S^{\mathrm{UNI}}\le S^{\mathrm{UC}}\le S^{\mathrm{COR}}$. All values are in bits.

The competitor is built from two copies of the CHSH strategy. Write $d=4k+\rho$ with $k\ge1$ and
$\rho\in\{0,1,2,3\}$, and split the computational basis of $\C^d$, and the set of outcomes, into consecutive groups of
sizes $4,\dots,4$ ($k$ times) and $\rho$. On a group of size four, $\C^4=\C^2\otimes\C^2$, Alice measures in the
product basis $e^x_{a_1}\otimes e^x_{a_2}$, where $e^x_0,e^x_1$ are the eigenvectors of $\mathsf A_x$ for the eigenvalues
$1,-1$, with outcome $2a_1+a_2$ within the group, and Bob likewise with $\mathsf B_y$: two copies of the CHSH strategy
with shared settings. On the last group both parties measure the single basis vector ($\rho=1$), the eigenbases of
$\mathsf A_x$ and $\mathsf B_y$ ($\rho=2$), or the DKZ bases for three outcomes ($\rho=3$). The resulting strategy $B_d$ on
$\Phi_d$ consists of four orthonormal bases of $\C^d$, so every outcome occurs.

\begin{theorem}[the Kullback--Leibler clause; computer-assisted]\label{thm:kl}
For every $d\ge4$ and $X\in\{\mathrm{UNI},\mathrm{UC},\mathrm{COR}\}$,
\[
S^X(\DKZ_d)\le0.0687802606<0.0703203697\le S^{\mathrm{UNI}}(B_d)\le S^X(B_d).
\]
If $4\mid d$, then $S^{\mathrm{UNI}}(B_d)\ge0.0879004622$. The same holds on $\Phi_{kd}$ for $\DKZ_d\otimes\one_k$ and
$B_d\otimes\one_k$. Hence for every $d\ge4$ the DKZ measurements do not give the best discrimination in any of the
three senses; the margin is at least $0.00154$ bits, and at least $0.0191$ bits if $4\mid d$.
\end{theorem}

\begin{proof}[Sketch of proof]
\emph{Test factors.} Let $r\ge0$ be a function of $(a,b,x,y)$, positive where $q$ is, with
$\sum_{x,y}\sigma_{xy}r(a_x,b_y|x,y)\le1$ for every $\lambda$. For $p\in\cL_d$, Jensen's inequality gives
$\sum\sigma_{xy}q\log(q/p)\ge\sum\sigma_{xy}q\log r-\log\sum\sigma_{xy}rp$, and the last sum is at most $1$ because it is
linear in $p$ and at most $1$ on every $\delta_\lambda$. Hence $U_q(\sigma)\ge\sum_{x,y}\sigma_{xy}\sum_{a,b}q\log r$.

\emph{The competitor.} The measurements of $B_d$ are block diagonal with disjoint outcome groups, so its behaviour is
the mixture $\sum_j\frac{m_j}dq_j$ of the behaviours $q_j$ of its blocks (of sizes $m_j$), each supported on the outcomes
of its group. For uniform $\sigma$, test factors of the blocks add: if each $r_j$ is a test factor of $q_j$ in its own
scenario with all entries at most $2$, then $r=r_j$ on the cells of block $j$, and $r=0$ elsewhere, is a test factor of
the mixture. Indeed, for a deterministic $\lambda$ only the setting pairs whose two outcomes lie in the same group
contribute; if these groups coincide, the contribution is at most that of a deterministic strategy of one block, and
otherwise the contributing pairs are $\{(1,1),(2,2)\}$ or $\{(1,2),(2,1)\}$ and contribute at most $\frac14(2+2)=1$.
Exact rational test factors for the four kinds of blocks, with all entries at most $1.24801$ and feasibility checked by
exhaustive exact enumeration, certify the values $0$ (one basis vector), $0.046273845411$ (CHSH), $0.057783024051$
($\DKZ_3$) and $L_4=0.087900462237$ (two copies of CHSH). Hence $S^{\mathrm{UNI}}(B_d)\ge(4kL_4+\rho L_\rho)/d$, where
$L_\rho$ is the value of the last block; this bound is nondecreasing in $k$, and its minimum over $d\ge4$ is
$\frac45L_4\ge0.0703203697$, attained at $d=5$.

\emph{The DKZ strategy.} The partial fraction expansion $\sum_{n\in\Z}(z+n)^{-2}=\pi^2/\sin^2(\pi z)$ shows that $h_{xy}$
in~\eqref{eq:dkzclosed} is the image under $\Z\to\Z_d$ of the probability distribution
$h^\infty_{xy}(m)=1/\bigl(2\pi^2(m-\delta_{xy})^2\bigr)$ on $\Z$, which does not depend on $d$. Let $\mu$ be a probability
distribution of a triple $(t,n_1,n_2)\in\Z^3$, and let $\mu_{11},\mu_{12},\mu_{21},\mu_{22}$ be the distributions of
$n_1,n_2,n_1-t,n_2-t$. With a uniformly random $c\in\Z_d$, the outputs $A_1=c$, $A_2=c+t$, $B_1=c+n_1$, $B_2=c+n_2$
modulo $d$ form a local model with behaviour $\frac1d\bar\mu_{xy}(b-a)$, where $\bar\mu_{xy}$ is the image of
$\mu_{xy}$ in $\Z_d$. Its divergence from $p_{\DKZ}$ on the pair $(x,y)$ is $D(h_{xy}\|\bar\mu_{xy})$, which is at most
$D(h^\infty_{xy}\|\mu_{xy})$ by the data-processing inequality. Hence
\[
S^{\mathrm{UNI}}(\DKZ_d)\le T(\mu):=\frac14\sum_{x,y}D\bigl(h^\infty_{xy}\big\|\mu_{xy}\bigr)\qquad\text{for every }d\ge2.
\]
An explicit distribution $\mu^*$, with rational weights in a finite window and a tail proportional to $1/(t^2-\frac14)$
(of total mass exactly $1$ by telescoping), gives $T(\mu^*)\le0.068780260568$: the sum over $|m|\le3000$ is evaluated in
interval arithmetic, and for $|m|>3000$ every term is nonpositive because there $\mu^*_{xy}(m)\ge h^\infty_{xy}(m)$.
Finally, the relabellings $q\mapsto q'$ with $q'(a,b|x,y)=q(-a,[y=1]-b|3-x,y)$ and
$q'(a,b|x,y)=q(-a-[x=2],-b|x,3-y)$ fix $p_{\DKZ}$ and generate a group acting transitively on the setting pairs; as
$\sigma\mapsto U_q(\sigma)$ is concave and invariant under this group, uniform settings are optimal, and
$S^{\mathrm{UNI}}=S^{\mathrm{UC}}=S^{\mathrm{COR}}$ for $\DKZ_d$.
\end{proof}

The certificates and programs are in the directory \texttt{complete-resolution/kl-divergence/}
of~\cite{repo-oqp27}. There, the program \texttt{blocks\_cert.py} checks the block test factors, which are stored in
\texttt{blocks\_testfactors.json}; the program \texttt{zcert.py} checks the distribution $\mu^*$, which is stored in
\texttt{zmodel\_cert\_A10.json}; and the program \texttt{verify\_main.py} re-checks the theorem independently from the
stored exact data.

\begin{corollary}\label{cor:klsup}
For $X\in\{\mathrm{UNI},\mathrm{UC},\mathrm{COR}\}$, $d\ge2$ and $k\ge1$, $S^X(\DKZ_d)\le S^X(\DKZ_{kd})$, and
\[
0.0687464\le\sup_dS^X(\DKZ_d)\le0.0687711 .
\]
\end{corollary}

\begin{proof}[Sketch of proof]
The infimum defining $U_q(\frac14)$ may be restricted to local models invariant under common shifts of all outcomes
(averaging over the shifts does not increase the divergence, by convexity), and these are the models of the proof of
Theorem~\ref{thm:kl} with $\Z_{kd}$ in place of $\Z$; reducing modulo $d$ gives the monotonicity. The lower bound is the
value at $d=128$ of an exact rational test factor invariant under common shifts, and the upper bound follows from a
two-sided rigorous enclosure of the tail of $T(\mu^*)$.
\end{proof}

\begin{proposition}[$d=3$; computer-assisted]\label{prop:kl3}
There is a neighbourhood of $\DKZ_3$ in the set of projective strategies on $\Phi_3$ on which each of the three
strengths is at most its value at $\DKZ_3$, $S(\DKZ_3)=0.05778302549332865$ bits, with equality only for strategies obtained
from $\DKZ_3$ by changing the phases of the basis vectors and by a local unitary $u\otimes\bar u$.
\end{proposition}

\begin{proof}[Sketch of proof]
For $d=3$, \eqref{eq:cglmp} reads $I_3=4-\sum_{x,y}\E[g_{xy}]$, where the level $g_{xy}\in\{0,1,2\}$ of a pair of
outcomes $(a,b)$ is the residue modulo $3$ of $a-b$, $b-a$, $b-a-1$, $a-b$ for $(x,y)=(1,1),(1,2),(2,1),(2,2)$; so the
CGLMP inequality is $\sum_{x,y}\E[g_{xy}]\ge2$. With
$\beta=\bigl[(16+24\sqrt3)-\sqrt{3280-960\sqrt3}\bigr]/54$, the test factor $r=1+\beta(\frac12-g_{xy})$ is feasible, and
$p_{\DKZ}/r$ is a mixture with positive weights of the $30$ deterministic strategies on the CGLMP facet; so $r$ is an
optimal test factor and gives the value. Every probability of a CGLMP level is stationary at $\DKZ_3$ (an exact
identity in $\Q(\sqrt3)$), so $\DKZ_3$ is a critical point. Moving the local model inside the CGLMP facet so that the
first- and second-order terms of the four divergences are equal, an interval Cholesky factorisation shows that the
averaged Hessian is at most $-\frac1{200}$ times the identity on a $16$-dimensional slice transversal to the orbit of
the phases and local unitaries. The rank of the measurements is locally constant, which gives the statement among all
projective strategies.
\end{proof}

Whether $\DKZ_3$ is a global maximiser is open; $330$ numerical searches found no strategy with a larger strength. As
$B_d$ consists of rank-one projective measurements, Theorem~\ref{thm:kl} also refutes the clause if POVMs or other ranks
are allowed. The Kullback--Leibler clause was first shown to fail for $d=4$ by Zhang~\cite{Zhang}, with an exact
certificate, and in~\cite{repo-cglmp} we gave exact certificates for $4\le d\le9$; numerically, $\DKZ_3$ appears to be
optimal~\cite{vDGG,AGG}. As far as we found, the additivity of test factors over blocks with disjoint outcomes and the
description of every $\DKZ_d$ as the reduction modulo $d$ of one distribution on $\Z$, which gives a bound uniform in
$d$, are new.

%======================================================================
\section{Open questions}\label{sec:open}
%======================================================================

(1) \emph{POVMs.} Theorems~\ref{thm:A} and~\ref{thm:B} concern projective measurements; the reduction of
Section~\ref{sec:reduction} uses projections in an essential way. For $3\le d\le8$ POVMs do not increase the maximum,
on maximally entangled states of every local dimension, and every optimal POVM strategy is projective, hence DKZ by
Theorem~\ref{thm:B}: we verified this with exact sum-of-squares certificates over $\mathbb Q(\zeta_{4d})$ in the tracial
moment relaxation with POVM constraints only (level $1{+}AB$ with doubly-localizing constraints), checked
independently (repository folder \texttt{complete-resolution/povm/}~\cite{repo-oqp27}). For $d\ge9$ the question is
open; the numerical evidence of~\cite{LVN} for $d=3$ pointed the same way. (2) \emph{All states.} The maximal value of $I_d$ over all states is
known exactly only for $d\le8$ (\cite{IR,repo-cglmp}). (3) \emph{Robustness of rigidity.} The deficit identity of
Remark~\ref{rem:deficit} suggests a quantitative version of Theorem~\ref{thm:B}; it would require a lower bound for the
defect in Theorem~\ref{thm:C} in terms of $\|[B,g]\|_F^2$. (4) \emph{An analytic proof of Theorem~\ref{thm:E}.} The
mechanism (Remark~\ref{rem:poisson}) and the continuum limit of the Gaussian modulation are explicit, and the large-$d$
argument of Section~\ref{subsec:regimeIII} is an expansion with certified remainders; a proof without interval
arithmetic for all $d$ seems within reach but is not attempted here. Such a proof, once formalised, would also complete
the Lean proof of Theorems~\ref{thm:A} and~\ref{thm:B} for $d\ge21$. (5) \emph{Noise on the state.} Is the critical
visibility of every projective strategy on a maximally entangled state, for white noise on the state and against all
Bell inequalities, at least $2/\IME(d)$, with equality only for DKZ (reading~(c) of Section~\ref{subsec:noisesetting})?
This is open for every $d\ge4$, already for rank-one measurements on $\Phi_d$. (6) \emph{Kullback--Leibler for small
$d$.} Whether $\DKZ_3$ maximises the statistical strength globally is open (Proposition~\ref{prop:kl3} gives strict
local optimality); the case $d=2$ is not treated here, and the exact value of $\sup_dS(\DKZ_d)$
(Corollary~\ref{cor:klsup}) is not known.

%======================================================================
\section{Verification and reproducibility}\label{sec:verification}
%======================================================================

\subsection{Computer-assisted components}
In the proof of Theorems~\ref{thm:A} and~\ref{thm:B}, only Theorem~\ref{thm:E} (and the convexity statement of
Lemma~\ref{lem:P}(1)) depend on computer calculations; the computations for Section~\ref{sec:clauses} are described in
Section~\ref{subsec:clausecomp}. All verifications for Theorem~\ref{thm:E} use interval arithmetic with outward
rounding (\texttt{mpmath}~\cite{mpmath}); floating-point computations
are used only to find candidates (linear-programming columns, approximate roots) and play no role in the verification.
Table~\ref{tab:programs} lists the programs; they are in the accompanying repository~\cite{repo-oqp27} (directory
\texttt{cone-certificates/}), together with their outputs and with written accounts of every bound they use. The
program \texttt{verify\_cert.py} of Regime~II has the SHA-256 hash
\begin{center}
\small\texttt{ff56740d15d40a07c99deb6433fb1b0f40b5a2d5b356b3ec9da0270ff1fab6a7},
\end{center}
which is stored in each of its $1800$ certificates and checked by \texttt{audit\_cert\_g.py}. The single command
\texttt{python audit\_alld.py} re-checks all three regimes and prints ALL-D CERTIFIED (every $d\ge2$).

\begin{table}[ht]
\small
\begin{tabular}{@{}>{\raggedright\arraybackslash}p{0.2\textwidth}>{\raggedright\arraybackslash}p{0.3\textwidth}>{\raggedright\arraybackslash}p{0.42\textwidth}@{}}
\toprule
Statement & Programs & Output and margins\\
\midrule
Prop.~\ref{prop:regimeI} ($d\le200$) & \texttt{verify\_cone.py}, \texttt{run\_verify.py}, \texttt{lp\_reverify.py};
check \texttt{crosscheck\_certs.py} & 199 certificates \texttt{cert\_d\{d\}.json}, all re-verified;
$d\min_kc_k\ge0.0816$\\
Prop.~\ref{prop:regimeII} ($201\le d\le2000$) & \texttt{verify\_cert.py} (tests \texttt{test\_verify\_cert.py}); audit
\texttt{audit\_cert\_g.py} & 1800 certificates \texttt{cert\_g\_d\{d\}.json}, all pass; $d^2\min\hat a_k\ge0.06319$ (even
$d$), $\ge0.12638$ (odd $d$); $d\min\Delta^2(P_v-H)_S\ge0.660$; derivative ratio $\ge4.0009$\\
Prop.~\ref{prop:regimeIII} ($d\ge2001$) & \texttt{a06\_zoneO.py}, \texttt{a07\_zoneI.py}, \texttt{a08\_A2.py},
\texttt{a09\_smallm.py}, \texttt{a10\_pointwise.py}, \texttt{a11\_slack.py}, \texttt{a12\_fp\_alld.py},
\texttt{a13\_constants.py}, \texttt{a14\_identities.py} (libraries \texttt{a02}--\texttt{a05}) & $Q_O\ge7.419$,
$\Lambda(R_2)\ge6.022$, $D_B(1)\ge-2.765192\varepsilon^2$, $|\Upsilon_*-\Upsilon_e|\le6.6468\cdot10^{-4}\varepsilon^2$,
$\hat a_k\ge0.1186944\,w_k\varepsilon^2$, $\Delta^2W\ge0.620154\,\varepsilon$, $|\mathsf G-\mathsf G_4|\le138.103\,\varepsilon^3$,
$|\partial_t\mathsf G_4|\le11.5703\,d$\\
Lemma~\ref{lem:P}(1) & \texttt{t35\_polya.py} & $\inf\mathcal A\ge0.1054878997577$\\
Theorem~\ref{thm:E}, all $d\ge2$ & \texttt{audit\_alld.py} & ALL-D CERTIFIED\\
\bottomrule
\end{tabular}
\caption{Programs for the computer-assisted part of the proof of Theorems~\ref{thm:A} and~\ref{thm:B}.}\label{tab:programs}
\end{table}

\subsection{The computations of Section~\ref{sec:clauses}}\label{subsec:clausecomp}
The results of Section~\ref{sec:clauses} use computer calculations at the following points, all in exact rational or
algebraic arithmetic or with rigorous rational enclosures; floating-point computations were used only to find
certificates. (i) Proposition~\ref{prop:dkzcert}: \texttt{dkz\_cert.py} checks the certificates
\texttt{certificates/dkz\_d3.json}, \dots, \texttt{dkz\_d20.json}. (ii) Proposition~\ref{prop:small}(2): the facet list
of $\cL_3$ is computed by \texttt{facets223.py} (exact double description, two insertion orders, every facet
re-verified) and stored in \texttt{facets223.txt}; \texttt{phi3.py} checks the remaining exact identities. (iii)
Remark~\ref{rem:exactdkz}: \texttt{dkz\_exact.py}. (iv) Remark~\ref{rem:scope}: \texttt{povm.py}. These programs are in
\texttt{complete-resolution/noise-literal/} of~\cite{repo-oqp27}, and the single command \texttt{python
verify\_theorem.py} there re-checks (i)--(iv) and the exact identities of Section~\ref{subsec:literal} (17 items, about
three minutes). (v) Theorem~\ref{thm:kl}: \texttt{zcert.py} (the bound on $T(\mu^*)$) and \texttt{blocks\_cert.py} (the
block test factors), with the independent re-check \texttt{verify\_main.py}; Corollary~\ref{cor:klsup}: \texttt{dkz\_lower.py}
and the programs of the independent verification; Proposition~\ref{prop:kl3}: \texttt{d3\_cert.py} (exact arithmetic
and an interval Cholesky factorisation). These programs are in \texttt{complete-resolution/kl-divergence/}.

\subsection{Independent verification}\label{subsec:independent}
Each analytic proof was refereed by an independent checker, who re-derived every step and wrote the numerical checks
from scratch, without importing any of our programs; the reports and scripts are in the repository~\cite{repo-oqp27}
(directories \texttt{proofs/} and \texttt{complete-resolution/}). There were four such verifications for the proof of
Theorems~\ref{thm:A}--\ref{thm:E}, and three for Section~\ref{sec:clauses}.
\begin{enumerate}
\item Theorems~\ref{thm:D} and~\ref{thm:C}, by the proofs of Section~\ref{sec:strip}: verdict correct. The independent
re-implementation (adaptive quadrature between exactly located branch points of $\rho$, and $\cD$ by a 30-digit matrix
exponential) reproduces~\eqref{eq:2bmv} at $82$ complex points $(a,t)$ on $13$ instances ($M=3,\dots,6$, all ranks for
$M\le5$), with worst relative error $6.8\cdot10^{-13}$, and the defect formula~\eqref{eq:star} at $51$ pairs (instance,
$\lambda$) on $9$ instances, including spectra up to $\pm150$, clustered spectra and $M=6$, to $1.3\cdot10^{-12}$
(absolute). The marginal identity $\int\rho(s,\tau)\dd s=\|BgP\|_F^2$ holds to $10^{-13}$, and the back-projection
integral in Step~2 equals $\rho$ to $10^{-20}$ (40-digit arithmetic, determinants). Fixed-order quadrature rules between
branch points are not sufficient: $\rho$ has complex singularities close to the real $s$-axis.
\item The distribution-free proof of Appendix~\ref{app:RI}, including the general case of
Remark~\ref{rem:RIgeneralB}: verdict correct. The identities involved, among them the Radon identity of
Theorem~\ref{thm:RI} on $22$ instances with $M\le8$, were confirmed to 34--46 digits.
\item The reduction chain (Section~\ref{sec:reduction}), the continuum theorem (Section~\ref{sec:continuum}), the cell
inequalities (Section~\ref{sec:cells}), the logic of the cone condition (Section~\ref{sec:cone}) and the rigidity argument
(Section~\ref{sec:proofsAB}): verdict correct.
\item The rank-one argument of Remark~\ref{rem:rankone}: verdict correct (no mathematical error; some expository gaps,
none of which changes a statement). The check includes the zero count on $1800$
extreme instances and the inequality at $M=3,\dots,8$ with entries up to $300$, re-evaluated at 50 digits.
\item Theorem~\ref{thm:literal}, first verification: verdict correct, no error. The checker re-derived every step,
recomputed the facets of $\cL_3$ by an independent double description in different coordinates, certified the
threshold of $C_d$ independently, re-verified all eighteen certificates of Proposition~\ref{prop:dkzcert} in 200-bit
interval arithmetic, confirmed the exact values of Remark~\ref{rem:exactdkz} for $d=3,4,5$ with explicit local models,
and ran eleven tamper tests on our verification program, all of which were detected.
\item Theorem~\ref{thm:literal}, second verification: verdict correct, no gap. The checker rebuilt the result from
scratch, with independent certificates for the DKZ strategy, and found the explicit local model of
Proposition~\ref{prop:flagmodel}.
\item Theorem~\ref{thm:kl}, Corollary~\ref{cor:klsup} and Proposition~\ref{prop:kl3}: verdict correct, no error. The
independent code reproduces $T(\mu^*)\le0.068780260568$ digit for digit, re-certifies the block values with independent
test factors, and re-proves the local optimality at $d=3$ with a different chart, slice and right inverse (averaged
Hessian at most $-\frac1{1000}$ times the identity on the slice); the sharper bounds of Corollary~\ref{cor:klsup} come
from this verification.
\end{enumerate}

\subsection{Further numerical checks}
The following checks are numerical; they are not part of any proof, and were made to guard against errors in formulas,
signs and constants.
\begin{itemize}
\item Sections~\ref{sec:reduction} and~\ref{sec:cells}: on random non-commuting families ($d=3,\dots,12$, $M=2,3,4$) the
identities of Lemmas~\ref{lem:projform} and~\ref{lem:linearform} and Proposition~\ref{prop:average} hold to $2\cdot10^{-12}$;
with the certificates of Proposition~\ref{prop:regimeI} at $d=25$ and $d=40$, the identity
$\ip v\delta=\sum_kc_k\ip{u^{\ell^{(k)}}}\delta$ holds to $1.6\cdot10^{-11}$ on random non-commuting configurations.
\item Section~\ref{sec:continuum}: the closed form of Lemma~\ref{lem:herglotz} agrees with the Herglotz integral to
$2\cdot10^{-16}$; the deficit identity~\eqref{eq:deficit} holds term by term on cell-embedded fields ($d=4$, $M=2,3$) to
$10^{-8}$; the values in Lemma~\ref{lem:classical} at $\alpha=\beta=\frac14$ agree to 40 digits.
\item Consistency: Theorems~\ref{thm:A} and~\ref{thm:B} were proved for $3\le d\le20$ by exact sum-of-squares certificates
in~\cite{repo-cglmp}, by a method sharing no step with the present one; and gradient ascents of $F$ over non-commuting
families ($d=3,\dots,6$, $M=2,3$) that reach $\FD$ end at commuting families, as Theorem~\ref{thm:B} predicts.
\end{itemize}

\subsection{Formalisation}\label{subsec:formal}
The argument of this paper has been formalised in the Lean~4 theorem prover~\cite{Lean} with the mathematical library
Mathlib~\cite{Mathlib}, in the versions
\begin{center}
\small Lean 4.33.1,\qquad Mathlib revision \texttt{0df444a360eaa60ab8c11dca51a86af692955474},
\end{center}
in $98$ files of about $33{,}400$ lines (directory \texttt{lean/} of~\cite{repo-oqp27}). No file uses \texttt{sorry},
\texttt{admit}, an \texttt{axiom} declaration or \texttt{native\_decide}, and \texttt{\#print axioms} reports, for every
main theorem, only the three standard axioms \texttt{propext}, \texttt{Classical.choice} and \texttt{Quot.sound}.
Table~\ref{tab:lean} lists the main theorems (file \texttt{Main.lean} and the files it imports, and the file
\texttt{CglmpNoise.lean}, which imports \texttt{Main.lean} and formalises Corollary~\ref{cor:noise}).

\begin{table}[ht]
\small
\begin{tabular}{@{}>{\raggedright\arraybackslash}p{0.35\textwidth}>{\raggedright\arraybackslash}p{0.41\textwidth}>{\raggedright\arraybackslash}p{0.18\textwidth}@{}}
\toprule
Lean theorem & Statement & Assumptions\\
\midrule
\nolinkurl{OQP27.maxEntClause_le_twenty} & for every $2\le d\le20$: Theorems~\ref{thm:A} and~\ref{thm:B}, and the DKZ
strategy attains $\IME(d)$ & none\\
\nolinkurl{OQP27.maxEntClause_all} & Theorems~\ref{thm:A} and~\ref{thm:B} for every $d\ge2$ & CONE$_d$ for $d\ge21$ only
(\nolinkurl{Hyp_ConeCertPos_large})\\
\nolinkurl{OQP27.stripInequality}, \nolinkurl{OQP27.stripEquality} & Theorem~\ref{thm:C} for every matrix size $M$, and
its equality case & none\\
\nolinkurl{OQP27.StripL3b.theorem1_bmv2} & Theorem~\ref{thm:D}, for all $(a,t)\in\C^2$ & none\\
\nolinkurl{OQP27.StripL3b.hyp_RI} & the Radon identity, Theorem~\ref{thm:RI} & none\\
\nolinkurl{OQP27.Red.hyp_reduction} & the exact reduction: strategy $\to$ clock model $\to$ Q-configuration $\to$ linear
form (Theorem~\ref{thm:reduction}, Lemma~\ref{lem:linearform}) & none\\
\nolinkurl{OQP27.Cell.continuumCell} & the continuum theorem for step fields: Theorem~\ref{thm:C} implies every cell
inequality (Proposition~\ref{prop:average}) & none\\
\nolinkurl{OQP27.ClassB.classicalTheoremB_all} & the classical (commuting) theorem for every $d$
(Theorem~\ref{thm:commutative}), including the discrete Hilbert-transform rearrangement ($N=4d$, $|A|=|B|=d$) & none\\
\nolinkurl{OQP27.coneCertPos_le_twenty} & Theorem~\ref{thm:E} for $2\le d\le20$: the certificates of
Proposition~\ref{prop:regimeI} re-evaluated by the Lean kernel in exact rational interval arithmetic & none\\
\nolinkurl{OQP27.ConeCertificate.fullCheck_sound} & soundness of the certificate checker & none\\
\nolinkurl{OQP27.dkz_cglmp} & the DKZ strategy attains $\IME(d)$, for every $d$ & none\\
\nolinkurl{OQP27.cglmpOf_noisy}, \nolinkurl{OQP27.noisy_violation_iff} & $I_d(vp+(1-v)p_0)=vI_d(p)$ for every behaviour
$p$; the CGLMP threshold of a strategy is $2/I_d$ & none\\
\nolinkurl{OQP27.noisy_violation_imp_le_twenty}, \nolinkurl{OQP27.noisy_threshold_rigid_le_twenty} &
Corollary~\ref{cor:noise} for $2\le d\le20$ & none\\
\nolinkurl{OQP27.noisy_violation_imp}, \nolinkurl{OQP27.noisy_threshold_rigid} & Corollary~\ref{cor:noise} for every
$d\ge2$ & CONE$_d$ for $d\ge21$ only\\
\bottomrule
\end{tabular}
\caption{The main theorems of the Lean formalisation. The last three rows are in the file \texttt{CglmpNoise.lean}.}\label{tab:lean}
\end{table}

Thus the maximally entangled clause of Problem~27B, Theorems~\ref{thm:A} and~\ref{thm:B}, is fully machine-checked for
$2\le d\le20$. For $d\ge21$ the formal proof is complete except for the single input CONE$_d$
(\nolinkurl{Hyp_ConeCertPos_large}), which is established by the interval-arithmetic computations of
Section~\ref{sec:cone}, outside Lean. Theorems~\ref{thm:C} and~\ref{thm:D} are formalised in full, for every matrix size
and with no assumption: the identity~\eqref{eq:2bmv} on $\C^2$, with $\rho\ge0$ measurable, supported in $K$ and bounded
by $C(\tau(1-\tau))^{-1/2}$, and the identity~\eqref{eq:star}, that is, the strip inequality with its exact defect, together
with its equality case. The formal proof of Theorem~\ref{thm:D} follows the distribution-free route of
Appendix~\ref{app:RI}: the Radon identity, then Tonelli's theorem and the identity theorem. In the uniqueness step~(b) of
the proof of Theorem~\ref{thm:C}, the formal proof uses the Fourier inversion theorem of Mathlib in place of
Lemma~\ref{lem:U}. The continuum theorem is formalised for step fields, with the explicit Herglotz functions of
Lemma~\ref{lem:herglotz}. The file \texttt{CglmpNoise.lean} defines the CGLMP expression of an arbitrary behaviour and
the uniform behaviour $p_0$; besides the theorems of Table~\ref{tab:lean}, \nolinkurl{OQP27.dkz_noisy_violation_iff}
states that the noisy DKZ behaviour violates the CGLMP inequality exactly when $v\IME(d)>2$. These theorems, too, depend
only on the three standard axioms.

A few auxiliary statements of this paper are proved on paper only, because the formal proof does not need them: the
continuity of $\rho$ (the formal proof uses its measurability, and a bound in Frobenius norm instead of the operator norm),
the marginal identity $\int\rho(s,\tau)\dd s=\|BgP\|_F^2$ of Theorem~\ref{thm:D}, and the closed formula~\eqref{eq:hlambda}
for $h_\lambda$ (in Lean, $h_\lambda$ is defined as the Poisson integral of Lemma~\ref{lem:h}(1)). Likewise, the formal
rigidity statement is the main assertion of Theorem~\ref{thm:B}: $d$ divides $D$, and a unitary $u$ maps the measurements to
the DKZ measurements tensored with the identity. The uniqueness of $u$ and the strict inequality for $d\nmid D$, which
follow from it by the short arguments at the end of Section~\ref{sec:proofsAB}, are not part of the formal statement.

\begin{remark}[the finite form of the cone condition]\label{rem:finiteform}
In Lean the cone condition is stated in finite form: $v=\sum_j\alpha_ju^{\ell_j}$ for finitely many $\ell_j\in\Delta_d$
and $\alpha_j\ge0$, with $\alpha_j>0$ for at least one $\ell_j\in\Delta_d^+$. For $21\le d\le200$ this is
Proposition~\ref{prop:regimeI}. For $d\ge201$ it follows from Theorem~\ref{thm:E}. Indeed, the measure $\mu$ of Propositions~\ref{prop:regimeII}
and~\ref{prop:regimeIII} is a probability measure
carried by $\Delta_d^+$ and $\ell\mapsto u^\ell$ is continuous, so the barycentre $\int u^\ell\mu(\dd\ell)$ lies in the convex
hull of $\{u^\ell:\ell\in\Delta_d^+\}$ (in finite dimension the barycentre of a probability measure carried by a set lies
in its convex hull), and by Carath\'eodory's theorem it is a convex combination of at most $d$ such vectors. Finally,
each $t_s$ is a positive multiple of $u^\ell$ for the cell vector $\ell$ with $\ell_0=a$, $\ell_s=d-a$ ($0<a<d$) and all
other cells empty. Indeed, by Lemma~\ref{lem:kernel}, $K^\ell(x,y)=0$ unless $x$ and $y$ are congruent to $0$ or $s$
modulo $d$, so $u^\ell$ is supported on $\{s,d-s\}$; its two entries are equal, by the reflections $\xi\mapsto d+a-\xi$
and $\xi\mapsto2d+a-\xi$ of $\T_N$, and positive, because they are sums of integrals of $\kappa(\xi-\eta)$ over cells
with $\xi-\eta\in(0,2d)$, where $\kappa>0$.
\end{remark}

\subsection*{Data and code availability}
Certificates, verification programs, logs, the Lean formalisation, the reports of the independent verifications and the
notes referred to above are available at~\cite{repo-oqp27},
\url{https://github.com/anshM123/anshM123-IQOQI-Vienna-Open-Problem-27}, under the MIT license. The material for
Section~\ref{sec:clauses} is in its directory \texttt{complete-resolution/}: the file \texttt{LEDGER.md} (the wording of
Problem~27 and the verdict on each clause), \texttt{noise-cglmp/} (Corollary~\ref{cor:noise}; the Lean file is
\texttt{lean/OQP27/CglmpNoise.lean}), \texttt{noise-literal/} (Theorem~\ref{thm:literal}, Propositions~\ref{prop:dkzcert}
and~\ref{prop:small}, Remarks~\ref{rem:exactdkz} and~\ref{rem:scope}) and \texttt{kl-divergence/} (Theorem~\ref{thm:kl},
Corollary~\ref{cor:klsup} and Proposition~\ref{prop:kl3}), each with the reports of its independent verifications. The
exact certificates for $d\le20$ of~\cite{repo-cglmp} are available at \url{https://github.com/anshM123/CGLMP}.

\appendix
%======================================================================
\section{The model for large \texorpdfstring{$d$}{d}}\label{app:model}
%======================================================================

This appendix defines the model $\Upsilon_e$ of Proposition~\ref{prop:regimeIII} and records the decomposition used in its
verification. Throughout $d\ge2001$, $\varepsilon=1/d$, and for real $b\in(0,d)$ we write $x=b/d$ and $w=1/b$.

\emph{Ingredients.} Let $\kappa_n(b)=\frac{\dd^n}{\dd b^n}[b\psi(b+1)]$ and $\mathrm B_{2j}(b)=b^{2j-1}\kappa_{2j}(b)/(2j-2)!$ (so that
$\mathrm B_{2j}(b)\to1$ as $b\to\infty$), $k_j=2(2j-2)!/(2^jj!)$. Let $g_r(x)=\sum_{i\text{ odd}}g_ix^i$ (Lemma~\ref{lem:singular}),
$a_n=(-1)^ng_r^{(n)}(1)/n!$, $P_n(x,\varepsilon)=\prod_{k<n}\frac{x+k\varepsilon}{1+k\varepsilon}$ (so $P_n=\E X_b^n$),
$\Gamma(x,\varepsilon)=\sum_{i\text{ odd}}g_iP_i(x,\varepsilon)=\Xb_r(b)/d^2$, and $S(x,\varepsilon)=\Gamma(x,\varepsilon)-\sum_na_nP_n(x,\varepsilon)$
(the second sum is $\Xb_r(d-b)/d^2$, since $X_{d-b}$ and $1-X_b$ have the same law). With
$c_j(b)=[\Xb^{(2j)}(b)-\Xb^{(2j)}(d-b)]/(2^jj!)$, the decoupled equation~\eqref{eq:decoupled} reads
$\sum_{j\ge1}c_j(m)t^j+(\text{remainder})=\tau_m$, and in the scaled unknown $u=t/(\varepsilon m^2)$, after multiplication by
$\frac\pi{2m}$, its coefficients are
\[
\mathrm{Ch}_j(b)=\tfrac\pi2\varepsilon^jb^{2j-1}c_j(b)=\varepsilon^{j-1}\Bigl\{k_j\Bigl[\mathrm B_{2j}(b)-\Bigl(\frac x{1-x}\Bigr)^{2j-1}\mathrm B_{2j}(d-b)\Bigr]
+\frac\pi2\,\frac{x^{2j-1}\partial_x^{2j}S(x,\varepsilon)}{2^jj!}\Bigr\}.
\]
Let $\mathrm p=g_r-\frac{\varepsilon^2}{12}g_r''+\frac{\varepsilon^4}{240}g_r''''$ (an Euler--Maclaurin interpolant), $\tilde P(b)=d^2[\mathrm p(x)-x(\mathrm p(1)-g_r(1))]$,
$\tau_e(b)=[\tilde P-\Xb_r](b)-[\tilde P-\Xb_r](d-b)$ and $\hat\tau_e=\frac\pi2\tau_e/b$.

\emph{The model.} $u_e(b)$ is the root in $(0,1)$ of $\sum_{j=1}^4\mathrm{Ch}_j(b)u^j=\hat\tau_e(b)$, and $\Upsilon_e(b)=\varepsilon b^2u_e(b)$;
$\Upsilon_e(0)=0$ and $\Upsilon_e(d-b)=\Upsilon_e(b)$. The scan of \texttt{a10\_pointwise.py} certifies that this root exists and is unique,
so $u_e$ is real-analytic; with the exact discrete targets and Gaussian order $7$ the model reproduces the certified
solutions at $d=1000$ to relative accuracy $10^{-18}$ (numerical check).

\emph{Decomposition.} Expanding $u_e=U_0(w)+xU_1(w)+x^2R_2(x,w)$ in $x$ at fixed $w$, one has $b^2U_0=A_0b^2/\mathrm B(b)$ and
$b^3U_1=\varphi_2(b)$ for an explicit function $\varphi_2$, so that
\[
\Upsilon_e(b)=\varepsilon A_0b^2+A_0\varepsilon f(b)+\varepsilon^2\varphi_2(b)+\varepsilon^3b^4R_2(x,w),
\]
and, since $\Delta^4$ annihilates quadratic polynomials,
$D_B[\Upsilon_e](m)=\varepsilon^2A_2(m)+\varepsilon^3\Delta^4[b^4R_2](m)-A_0\varepsilon\Delta^4[f(d-\cdot)](m)$ with $A_2=\Delta^4\varphi_2$. For $m\ge3$,
$\Delta^4h(m)=\int_{-2}^2h''''(m+s)M_4(s)\dd s$ with the cardinal B-spline $M_4\ge0$ of mass $1$, and
$\frac{\dd^4}{\dd b^4}[b^4R_2]=\Lambda(R_2)$ with $\Lambda=(\vartheta+1)(\vartheta+2)(\vartheta+3)(\vartheta+4)$, $\vartheta=x\partial_x-w\partial_w$. In the
\emph{inner zone} ($x\le0.013$) the certified bounds are $A_2(m)\ge0$ for $m\ge2$ and $\Lambda(R_2)\ge6.022$; in the \emph{outer zone}
($x\in[0.011,0.501]$) one uses instead $\varepsilon^{-3}\Delta^4[\Upsilon_e-\varepsilon A_0b^2/\mathrm B](m)=\int Q_O(x_s,\varepsilon)M_4(s)\dd s$ with
$Q_O=\partial_x^4[x^2(u_e-A_0/\mathrm B(1/w))]$ at fixed $\varepsilon$, certified $\ge7.419$. Since the stencil of $\Delta^4$ has width
$4/d\le0.002$, every $m$ is covered by one of the two zones; $m=1,2$ are treated directly. All enclosures are Taylor models in
$(x,\varepsilon)$ or $(w,\varepsilon)$ with $\varepsilon\in[0,1/2001]$: thin centre jets of the model root (Krawczyk iteration, then
implicit differentiation degree by degree), Lagrange remainders from jets over the full box, special functions (polygamma
and Binet functions, Pochhammer sums with rigorous tails) by interval series. For boxes with $x>0.42$ the exactly
antisymmetric coefficients are divided by $1-2x$ before enclosing (they vanish at $x=\frac12$).

\emph{Pointwise lemma and slack.} With the exact targets $\tau_m$ and the Gaussian expansion to order $7$ plus remainder,
\texttt{a10\_pointwise.py} bounds the difference between the decoupled solution and the model by $6.6468\cdot10^{-4}\varepsilon^2$,
using analytic bounds for the Euler--Maclaurin defect of $\tilde P$ (at most $235.6\,\varepsilon^6$ for its second differences), for
the Gaussian remainder, and for the higher coefficients $\mathrm{Ch}_5,\dots,\mathrm{Ch}_7$. For the slack, let
$\mathrm{Dl}_e=\tilde P-\Xb_r$ (the model of $\mathrm{Dl}$, by Lemma~\ref{lem:singular}) and
$W_e(b)=\mathrm{Dl}_e(b)-\sum_{j=1}^4\Upsilon_e(b)^j\,\Xb^{(2j)}(b)/(2^jj!)$, which is symmetric under $b\mapsto d-b$ by the model
equation. Its leading term is exactly linear, $-(1-\frac2\pi)b$, so $\Delta^2W_e=O(\varepsilon)$; the density
$\Delta^2W_e/\varepsilon$ and the perturbation from $W_e$ to $W$ are certified as stated in Section~\ref{subsec:regimeIII}.

%======================================================================
\section{A distribution-free proof of Theorem~\ref{thm:D}}\label{app:RI}
%======================================================================

The proof of Theorem~\ref{thm:D} in Section~\ref{subsec:proofD} uses the Paley--Wiener--Schwartz theorem, Fourier
inversion along lines and the structure of distributions supported on a line, and the proof of Theorem~\ref{thm:C} in
Section~\ref{subsec:proofC} uses the uniqueness of Fourier transforms. Here we give a second proof of both theorems,
which uses no distributions. Its tools are linear algebra, the Cauchy integral formula on a circle and Hurwitz's theorem,
differentiation under the integral sign, dominated convergence, Weierstrass' theorem on locally uniform limits of
holomorphic functions, the theorems of Fubini and Tonelli, the identity theorem in one variable, the Weierstrass
approximation theorem, and a Gaussian approximate identity in one variable (Lemma~\ref{lem:U}). Its key step is a Radon
identity for the pencil (Theorem~\ref{thm:RI}), which we prove through
a complex inviscid Burgers equation satisfied by the roots of the pencil. This is the route followed by the Lean
formalisation (Section~\ref{subsec:formal}).

\subsection{The Radon identity}
Throughout this appendix $H$ is a Hermitian $M\times M$ matrix, $B$ an orthogonal projection of rank $r$, $P=\one-B$, and
$H_d:=BHB+PHP$. Let $\lambda_1,\dots,\lambda_M$ be the eigenvalues of $H$ and $\lambda^0_1,\dots,\lambda^0_M$ those of $H_d$,
which consist of the eigenvalues $\mu_1,\dots,\mu_{M-r}$ of $PHP|_{\operatorname{ran}P}$ and $\beta_1,\dots,\beta_r$ of
$BHB|_{\operatorname{ran}B}$; note that $\Tr H=\Tr H_d$. We write $\C_\pm=\{z:\pm\Im z>0\}$, and $\operatorname{Log}$ is the
principal logarithm on $\C\setminus(-\infty,0]$. For $0<\tau<1$ and $c\in\C$ put
\[
p(y;\tau,c)=\det\bigl(H+c-y(P-\tau)\bigr),\qquad Y(\tau,c)=(P-\tau)^{-1}(H+c)=\Bigl(\frac P{1-\tau}-\frac B\tau\Bigr)(H+c).
\]
Since $H+c-y(P-\tau)=(P-\tau)(Y(\tau,c)-y)$, the roots of $p(\cdot\,;\tau,c)$, with multiplicity, are the eigenvalues of
$Y(\tau,c)$, and the leading coefficient $\det(\tau-P)=(\tau-1)^{M-r}\tau^r$ does not depend on $c$. Let
$y_1(\tau),\dots,y_M(\tau)$ be the roots for $c=0$, and put
\[
\mathsf L(H)=\frac1{2\pi}\int_0^1\sum_{i=1}^M\bigl|\Im y_i(\tau)\bigr|\dd\tau,\qquad \mathsf R(H)=\Tr H_+-\Tr(H_d)_+ .
\]

\begin{theorem}[Radon identity]\label{thm:RI}
The function $\tau\mapsto\sum_i|\Im y_i(\tau)|$ is continuous on $(0,1)$ and bounded by $M\|H\|/\sqrt{\tau(1-\tau)}$,
and $\mathsf L(H)=\mathsf R(H)$.
\end{theorem}

If $r\in\{0,M\}$, then $P-\tau$ is a nonzero multiple of $\one$, all roots are real and $H_d=H$, so both sides vanish.
We assume $1\le r\le M-1$ from now on. Applied to $H=g-\xi P-w$, Theorem~\ref{thm:RI} says that the integrals of the
density $\rho$ of Theorem~\ref{thm:D} along the lines $s=w+\xi\tau$ are the real slices $U_\xi(w)$ of
Lemma~\ref{lem:slices} (Section~\ref{subsec:RItoD}); this explains the name.

\subsection{The complexified pencil}
In this subsection $\Im c>0$ and $0<\tau<1$, and we put $\Gamma_c=\|H\|+|c|$. For a root $y$ of $p(\cdot\,;\tau,c)$ we
fix a unit vector $v$ with $(H+c)v=y(P-\tau)v$, and we put $a=\ip v{(P-\tau)v}=\|Pv\|^2-\tau\in[-\tau,1-\tau]$.

\begin{lemma}[location and count]\label{lem:RIa}
\begin{enumerate}
\item $\Im c=a\Im y$. Hence $a\ne0$, $y\in\C_+$ if and only if $a>0$, $\Im y\ge\Im c/(1-\tau)$ if $y\in\C_+$, and
$\Im y\le-\Im c/\tau$ if $y\in\C_-$.
\item No root of $p(\cdot\,;\tau,c)$ is real, and exactly $M-r$ roots, counted with multiplicity, lie in $\C_+$.
\item For $H_d$ in place of $H$ the roots are the numbers $(\mu_k+c)/(1-\tau)\in\C_+$ and $-(\beta_j+c)/\tau\in\C_-$.
\end{enumerate}
\end{lemma}

\begin{proof}
(1) $\ip v{(H+c)v}=\ip v{Hv}+c$ and $\ip v{(H+c)v}=ya$, where $\ip v{Hv}$ and $a$ are real; taking imaginary parts gives
$\Im c=a\Im y$. As $\Im c>0$, $a\ne0$ and $\Im y=\Im c/a$, with $0<a\le1-\tau$ or $-\tau\le a<0$.
(2) No root is real by (1). For $\sigma\in[0,1]$, part (1) applies to the Hermitian matrix $\sigma H$, so no root of
$\det(\sigma H+c-y(P-\tau))$ is real. The leading coefficient does not depend on $\sigma$, so the roots, as an unordered
$M$-tuple, depend continuously on $\sigma$, and the number of roots in $\C_+$ is constant on $[0,1]$. At $\sigma=0$ the
polynomial is $(c-y(1-\tau))^{M-r}(c+y\tau)^r$, with the root $c/(1-\tau)\in\C_+$ of multiplicity $M-r$ and the root
$-c/\tau\in\C_-$ of multiplicity $r$.
(3) $H_d+c-y(P-\tau)$ is block diagonal, equal to $PHP+c-y(1-\tau)$ on $\operatorname{ran}P$ and to $BHB+c+y\tau$ on
$\operatorname{ran}B$.
\end{proof}

\begin{lemma}[bounds]\label{lem:RIb}
\begin{enumerate}
\item Every root satisfies $|y|\le\Gamma_c/\min(\tau,1-\tau)$.
\item If $y\in\C_+$ and $\tau\le\frac12$, or $y\in\C_-$ and $\tau\ge\frac12$, then $|y|\le\Gamma_c/\sqrt{\tau(1-\tau)}$.
\item Let $\Sigma_\pm(\tau,c)$ be the sums of the roots in $\C_\pm$, and $\Sigma^0_+(\tau,c)=(\Tr_P(PHP)+(M-r)c)/(1-\tau)$,
$\Sigma^0_-(\tau,c)=-(\Tr_B(BHB)+rc)/\tau$ the corresponding sums for $H_d$. Then $\Sigma_++\Sigma_-=\Sigma^0_++\Sigma^0_-$, and
\[
|\Sigma_+(\tau,c)-\Sigma^0_+(\tau,c)|\le M\Gamma_c\bigl(2+(\tau(1-\tau))^{-1/2}\bigr).
\]
\end{enumerate}
\end{lemma}

\begin{proof}
We have $|y|\,|(P-\tau)v|=|(H+c)v|\le\Gamma_c$. (1) $|(P-\tau)v|^2=(1-\tau)^2|Pv|^2+\tau^2|Bv|^2\ge\min(\tau,1-\tau)^2$.
(2) With $|Pv|^2=a+\tau$, $|(P-\tau)v|^2=(1-\tau)^2(a+\tau)+\tau^2(1-\tau-a)=\tau(1-\tau)+a(1-2\tau)$, and $a(1-2\tau)\ge0$
in both cases by Lemma~\ref{lem:RIa}(1). (3) $\Sigma_++\Sigma_-=\Tr Y(\tau,c)$, and
$\Tr[(\frac P{1-\tau}-\frac B\tau)(H-H_d)]=\Tr[(\frac P{1-\tau}-\frac B\tau)(PHB+BHP)]=0$, so $\Tr Y(\tau,c)$ is the same for $H$
and for $H_d$. For $\tau\le\frac12$, part (2) gives $|\Sigma_+|\le(M-r)\Gamma_c/\sqrt{\tau(1-\tau)}$, and
$|\Sigma^0_+|\le(M-r)\Gamma_c/(1-\tau)\le2(M-r)\Gamma_c$. For $\tau\ge\frac12$ write $\Sigma_+-\Sigma^0_+=\Sigma^0_--\Sigma_-$
and use part (2) for the roots in $\C_-$ and $|\Sigma^0_-|\le r\Gamma_c/\tau\le2r\Gamma_c$.
\end{proof}

Let $\mathcal L_+(\tau,c)=\sum\operatorname{Log}y$, the sum over the roots $y$ in $\C_+$ with multiplicity, and
$\mathcal L^0_+(\tau,c)=\sum_k\operatorname{Log}\frac{\mu_k+c}{1-\tau}$, the same sum for $H_d$.

\begin{lemma}[the Burgers identity]\label{lem:RIc}
On $(0,1)\times\C_+$ the functions $\Sigma_+$ and $\mathcal L_+$ are continuously differentiable in $\tau$ and holomorphic
in $c$, their partial derivatives are jointly continuous, and
\[
\partial_\tau\mathcal L_+=\partial_c\Sigma_+,\qquad \partial_\tau\mathcal L^0_+=\partial_c\Sigma^0_+=\frac{M-r}{1-\tau}.
\]
\end{lemma}

\begin{proof}
Fix $(\tau_0,c_0)\in(0,1)\times\C_+$, put $\eta=\frac12\Im c_0$ and $R_0=2(\|H\|+|c_0|+1)/\min(\tau_0,1-\tau_0)$, and let
$U$ be a neighbourhood of $(\tau_0,c_0)$ on which $\Im c\ge\eta$, $|c|\le|c_0|+1$ and
$\min(\tau,1-\tau)\ge\frac12\min(\tau_0,1-\tau_0)$. By Lemmas~\ref{lem:RIa}(1) and~\ref{lem:RIb}(1), for $(\tau,c)\in U$ the
roots in $\C_+$ lie in $Q=\{y:\Im y\ge\eta,\ |y|\le R_0\}$. Let $\Omega$ be the open disc with centre $iT$ and radius
$T-\frac\eta2$, where $T=R_0^2/\eta+1$. Its closure lies in $\{\Im y\ge\frac\eta2\}\subseteq\C_+$, and $Q\subseteq\Omega$: for
$y=u+iw\in Q$, $|y-iT|^2=|y|^2-2wT+T^2\le R_0^2-2\eta T+T^2<(T-\frac\eta2)^2$, because $\eta T>R_0^2$. So for
$(\tau,c)\in U$ the roots in $\C_+$ lie in $\Omega$, the roots in $\C_-$ lie outside $\bar\Omega$, and $p(\cdot\,;\tau,c)$ has
no zero on $\partial\Omega$. Writing $p=\det(\tau-P)\prod_i(y-y_i)$, so that $p_y/p=\sum_i(y-y_i)^{-1}$, the Cauchy integral
formula gives, for every $\varphi$ holomorphic near $\bar\Omega$,
\begin{equation}\label{eq:RIcauchy}
\sum_{y_i\in\C_+}\varphi(y_i)=\frac1{2\pi i}\oint_{\partial\Omega}\varphi(y)\,\frac{p_y(y;\tau,c)}{p(y;\tau,c)}\dd y .
\end{equation}
We take $\varphi(y)=y$ and $\varphi=\operatorname{Log}$. The integrand is a smooth function of $(y,\tau,\Re c,\Im c)$ on
$\partial\Omega\times U$ and is holomorphic in $c$, so differentiation under the integral sign gives the regularity
statements. Next, $p(y;\tau,c)=f(y,c+y\tau)$ with the polynomial $f(y,w')=\det(H+w'-yP)$, so $p_\tau=y\,p_c$. With
$q=p_c/p$ we get $\partial_\tau(p_y/p)=\partial_y(p_\tau/p)=\partial_y(yq)$ and $\partial_c(p_y/p)=\partial_yq$, hence
\[
\partial_\tau\mathcal L_+-\partial_c\Sigma_+=\frac1{2\pi i}\oint_{\partial\Omega}\bigl[\operatorname{Log}(y)\,\partial_y(yq)-y\,\partial_yq\bigr]\dd y
=\frac1{2\pi i}\oint_{\partial\Omega}\partial_y\bigl[(y\operatorname{Log}y-y)\,q\bigr]\dd y=0,
\]
because $\operatorname{Log}(y)(q+yq_y)-yq_y=\partial_y[(y\operatorname{Log}y-y)q]$ is the derivative of a function that is
holomorphic near $\partial\Omega$. For $H_d$ the identity follows from the explicit formulas
$\mathcal L^0_+=\sum_k\operatorname{Log}(\mu_k+c)-(M-r)\log(1-\tau)$ and $\Sigma^0_+=(\sum_k\mu_k+(M-r)c)/(1-\tau)$.
\end{proof}

For a simple root $y(\tau,c)$, formula~\eqref{eq:RIcauchy} on a small circle around $y$ (in place of $\partial\Omega$) reduces Lemma~\ref{lem:RIc} to
$\partial_\tau y=y\,\partial_cy$, the complex inviscid Burgers equation; the lemma is the form of this equation, summed
over the roots in $\C_+$, that remains valid at multiple roots.

\begin{lemma}[the two ends]\label{lem:RId}
Fix $c\in\C_+$.
\begin{enumerate}
\item As $\tau\to1^-$, $\mathcal L_+(\tau,c)-\mathcal L^0_+(\tau,c)\to0$.
\item Let $T_0=PHP+c-PHB(BHB+c)^{-1}BHP$ on $\operatorname{ran}P$ (the matrix $BHB+c$ is invertible on $\operatorname{ran}B$
since $\Im c>0$). Then $\spec T_0\subseteq\C_+$, and as $\tau\to0^+$,
\[
\mathcal L_+(\tau,c)-\mathcal L^0_+(\tau,c)\to\theta_0(c):=\sum_{t\in\spec T_0}\operatorname{Log}t-\sum_k\operatorname{Log}(\mu_k+c),
\qquad e^{\theta_0(c)}=\frac{\det(H+c)}{\det(H_d+c)} .
\]
\end{enumerate}
\end{lemma}

\begin{proof}
(1) $(1-\tau)Y(\tau,c)=(P-\frac{1-\tau}\tau B)(H+c)\to P(H+c)$ as $\tau\to1$. With respect to
$\operatorname{ran}P\oplus\operatorname{ran}B$, $P(H+c)$ is block upper triangular with diagonal blocks $PHP+c$ and $0$, so its
eigenvalues are the numbers $\mu_k+c\in\C_+$ and $0$ ($r$ times). By continuity of eigenvalues, for $\tau$ near $1$ exactly
$M-r$ eigenvalues $z_k$ of $(1-\tau)Y$ lie near the $\mu_k+c$, hence in $\C_+$, labelled so that $z_k\to\mu_k+c$, and $r$
lie near $0$. Multiplication by $1/(1-\tau)>0$ preserves $\C_+$, and $Y$ has exactly $M-r$ eigenvalues in $\C_+$
(Lemma~\ref{lem:RIa}(2)); so the roots in $\C_+$ are the $z_k/(1-\tau)$, and since
$\operatorname{Log}(z/(1-\tau))=\operatorname{Log}z-\log(1-\tau)$ for $z\in\C_+$,
$\mathcal L_+-\mathcal L^0_+=\sum_k[\operatorname{Log}z_k-\operatorname{Log}(\mu_k+c)]\to0$.

(2) By the Schur complement formula with respect to $\operatorname{ran}P\oplus\operatorname{ran}B$,
$p_0(y):=\det(H+c-yP)=\det_B(BHB+c)\det_P(T_0-y)$, a polynomial of exact degree $M-r$ whose roots form $\spec T_0$. For
real $y$ the matrix $H-yP$ is Hermitian and $\Im c>0$, so $p_0(y)\ne0$: $\spec T_0$ contains no real number. As
$\tau\to0$, $p(y;\tau,c)=\det(H+c-yP+y\tau)\to p_0(y)$ coefficientwise. By Hurwitz's theorem, applied on a circle that
encloses $\spec T_0$ and on small circles around its points, for small $\tau$ exactly $M-r$ roots of $p(\cdot\,;\tau,c)$ lie
in a fixed disc, and they converge to the points of $\spec T_0$ with multiplicity. The other $r$ roots leave every
bounded set: $\tau Y(\tau,c)=(\frac\tau{1-\tau}P-B)(H+c)\to-B(H+c)$, whose eigenvalues are $0$ ($M-r$ times) and the
numbers $-(\beta_j+c)$, which lie in $\C_-$ at distance at least $\Im c$ from $\R$; so $r$ eigenvalues of $Y$ are of the form
$-(\beta_j+c+o(1))/\tau$ and lie in $\C_-$ for small $\tau$. Since $\C_-$ contains exactly $r$ roots, the $M-r$ bounded roots
are those in $\C_+$. Their limits lie in the closure of $\C_+$ and are not real, so $\spec T_0\subseteq\C_+$ and, by
continuity of $\operatorname{Log}$ on $\C_+$, $\mathcal L_+\to\sum_{t\in\spec T_0}\operatorname{Log}t$. Also
$\mathcal L^0_+=\sum_k\operatorname{Log}(\mu_k+c)-(M-r)\log(1-\tau)\to\sum_k\operatorname{Log}(\mu_k+c)$. Finally
$e^{\theta_0}=\det_PT_0/\det_P(PHP+c)=\det(H+c)/(\det_B(BHB+c)\det_P(PHP+c))=\det(H+c)/\det(H_d+c)$.
\end{proof}

\begin{lemma}[boundary values]\label{lem:RIe}
Let $c_0\in\R$, $0<\tau<1$, and let $y_1,\dots,y_M$ be the roots of $p(\cdot\,;\tau,c_0)$. As $c\to c_0$ within $\C_+$,
$\Im\Sigma_+(\tau,c)\to\frac12\sum_i|\Im y_i|$ and $\Im\Sigma^0_+(\tau,c)=(M-r)\Im c/(1-\tau)\to0$. If $H+c_0>0$, all $y_i$
are real.
\end{lemma}

\begin{proof}
For real $y$, $p(y;\tau,c_0)$ is the determinant of a Hermitian matrix, hence real; so the non-real roots of
$p(\cdot\,;\tau,c_0)$ come in conjugate pairs. The roots depend continuously on $c$, since the leading coefficient does not.
If $z$ is a root of $p(\cdot\,;\tau,c_0)$ of multiplicity $m_z$, then for $c$ near $c_0$ exactly $m_z$ roots lie near $z$; they lie
in $\C_+$ if $z\in\C_+$, in $\C_-$ if $z\in\C_-$, and their imaginary parts tend to $0$ if $z$ is real. Summing over the roots
in $\C_+$ gives $\Im\Sigma_+(\tau,c)\to\sum_{\Im z>0}m_z\Im z=\frac12\sum_i|\Im y_i|$. If $A=H+c_0>0$ and $y$ is a non-real
root with eigenvector $v$, then $\ip v{Av}=y\ip v{(P-\tau)v}$ with both inner products real, which forces
$\ip v{(P-\tau)v}=0=\ip v{Av}$, a contradiction.
\end{proof}

\subsection{Proof of Theorem~\ref{thm:RI}}
\emph{Step 1.} For $c\in\C_+$ let $\Psi(c)=\int_0^1[\Sigma_+(\tau,c)-\Sigma^0_+(\tau,c)]\dd\tau$; by
Lemma~\ref{lem:RIb}(3) the integral converges absolutely, uniformly for $c$ in compact subsets of $\C_+$. For
$0<\delta<\frac12$ the function $\Psi_\delta(c)=\int_\delta^{1-\delta}(\Sigma_+-\Sigma^0_+)\dd\tau$ is holomorphic on $\C_+$
(Lemma~\ref{lem:RIc} and differentiation under the integral sign), and by Lemma~\ref{lem:RIc} and the fundamental theorem of
calculus,
\[
\Psi_\delta'(c)=\int_\delta^{1-\delta}\partial_c(\Sigma_+-\Sigma^0_+)\dd\tau=\int_\delta^{1-\delta}\partial_\tau(\mathcal L_+-\mathcal L^0_+)\dd\tau
=(\mathcal L_+-\mathcal L^0_+)(1-\delta,c)-(\mathcal L_+-\mathcal L^0_+)(\delta,c).
\]
As $\delta\to0$, $\Psi_\delta\to\Psi$ locally uniformly on $\C_+$, so $\Psi$ is holomorphic and $\Psi_\delta'\to\Psi'$
(Weierstrass); by Lemma~\ref{lem:RId}, $\Psi'(c)=-\theta_0(c)$.

\emph{Step 2.} The function $\theta_1(c)=\sum_j\operatorname{Log}(\lambda_j+c)-\sum_j\operatorname{Log}(\lambda^0_j+c)$ is
holomorphic on $\C_+$ and $e^{\theta_1}=\det(H+c)/\det(H_d+c)=e^{\theta_0}$. Hence $\theta_0-\theta_1$ is continuous with values
in $2\pi i\Z$, so it equals a constant $2\pi ik$ on the connected set $\C_+$. Let
\[
\mathcal E(c)=\Tr\bigl[(H_d+c)\operatorname{Log}(H_d+c)\bigr]-\Tr\bigl[(H+c)\operatorname{Log}(H+c)\bigr],
\]
that is, $\mathcal E(c)=\sum_j(\lambda^0_j+c)\operatorname{Log}(\lambda^0_j+c)-\sum_j(\lambda_j+c)\operatorname{Log}(\lambda_j+c)$.
Since $\frac{\dd}{\dd c}[(\lambda+c)\operatorname{Log}(\lambda+c)]=\operatorname{Log}(\lambda+c)+1$ and both spectra have $M$
elements, $\mathcal E'=-\theta_1$. Hence $(\Psi-\mathcal E)'=-\theta_0+\theta_1=-2\pi ik$, and
$\Psi(c)=\mathcal E(c)-2\pi ikc+\kappa_0$ on $\C_+$ for a constant $\kappa_0\in\C$.

\emph{Step 3.} Fix $c_0\in\R$ and let $\epsilon\to0^+$. By Lemma~\ref{lem:RIe}, for every $\tau$,
$\Im(\Sigma_+-\Sigma^0_+)(\tau,c_0+i\epsilon)\to\frac12\sum_i|\Im y_i(\tau;c_0)|$, where the $y_i(\tau;c_0)$ are the roots of
$p(\cdot\,;\tau,c_0)$, that is, the roots of Theorem~\ref{thm:RI} for $H+c_0$ in place of $H$; and by Lemma~\ref{lem:RIb}(3)
the integrand is dominated by $M(\|H\|+|c_0|+1)(2+(\tau(1-\tau))^{-1/2})$ for $\epsilon\le1$. By dominated convergence,
$\Im\Psi(c_0+i\epsilon)\to\pi\mathsf L(H+c_0)$. On the other side, as $z\to t\in\R$ within $\C_+$,
$z\operatorname{Log}z\to t\log|t|+i\pi\min(t,0)$ (with limit $0$ at $t=0$). Since $\min(t,0)=t-t_+$ and
$\sum_j\lambda_j=\sum_j\lambda^0_j$,
\begin{align*}
\Im\mathcal E(c_0+i\epsilon)&\to\pi\Bigl[\sum_j\min(\lambda^0_j+c_0,0)-\sum_j\min(\lambda_j+c_0,0)\Bigr]\\
&=\pi\Bigl[\sum_j(\lambda_j+c_0)_+-\sum_j(\lambda^0_j+c_0)_+\Bigr]=\pi\mathsf R(H+c_0).
\end{align*}
Taking imaginary parts in $\Psi=\mathcal E-2\pi ikc+\kappa_0$ gives
$\mathsf L(H+c_0)=\mathsf R(H+c_0)-2kc_0+\Im\kappa_0/\pi$ for every real $c_0$.

\emph{Step 4.} For $c_0>\|H\|$ the matrices $H+c_0$ and $H_d+c_0$ are positive definite ($\|H_d\|\le\|H\|$). By
Lemma~\ref{lem:RIe} all roots are then real, so $\mathsf L(H+c_0)=0$, and $\mathsf R(H+c_0)=\Tr(H+c_0)-\Tr(H_d+c_0)=0$.
Hence $-2kc_0+\Im\kappa_0/\pi=0$ for all $c_0>\|H\|$, so $k=0$ and $\Im\kappa_0=0$, and $\mathsf L(H+c_0)=\mathsf R(H+c_0)$ for
every real $c_0$; the case $c_0=0$ is the identity. Finally, $\tau\mapsto\sum_i|\Im y_i(\tau)|$ is continuous on $(0,1)$ by
continuity of the roots, and every non-real root satisfies $|y|\le\|H\|/\sqrt{\tau(1-\tau)}$, by the argument of
Lemma~\ref{lem:pencil}(3). \qed

One can show in addition that $\kappa_0=0$, that is, $\Psi=\mathcal E$ on $\C_+$; we do not need this.

\subsection{Theorem~\ref{thm:D} from the Radon identity}\label{subsec:RItoD}
Let $g$ be Hermitian, let $\rho$ and the roots $\xi_i(s,\tau)$ be as in Theorem~\ref{thm:D}, and let $\lambda_j(\xi)$,
$\lambda^0_j(\xi)$ and $U_\xi$ be as in Lemma~\ref{lem:slices}. The properties of $\rho$ stated in Theorem~\ref{thm:D} are those
of Lemma~\ref{lem:pencil}, whose proof is elementary. Fix $\xi,w\in\R$ and put $H=g-\xi P-w$. For $s=w+\xi\tau$,
\[
g-s-x(P-\tau)=H-(x-\xi)(P-\tau),
\]
so the roots of $p_{s,\tau}$ are the numbers $\xi+y_i(\tau)$, with $y_i(\tau)$ the roots of Theorem~\ref{thm:RI} for this $H$,
and $\rho(w+\xi\tau,\tau)=\frac1{2\pi}\sum_i|\Im y_i(\tau)|$. Moreover $H_d=g_d-\xi P-w$. Theorem~\ref{thm:RI} therefore gives
\begin{equation}\label{eq:radonslice}
\int_0^1\rho(w+\xi\tau,\tau)\dd\tau=\sum_j\bigl(\lambda_j(\xi)-w\bigr)_+-\sum_j\bigl(\lambda^0_j(\xi)-w\bigr)_+=U_\xi(w),
\end{equation}
the last equality because $(l-w)_+-(w-l)_+=l-w$ and $\sum_j\lambda_j(\xi)=\sum_j\lambda^0_j(\xi)$. Thus the integrals of $\rho$
along all non-horizontal lines are the real slices $U_\xi$.

For real $a>0$, $\int_\R e^{aw}(l-w)_+\dd w=e^{al}/a^2$, and for real $a<0$, $\int_\R e^{aw}(w-l)_+\dd w=e^{al}/a^2$. Using the
middle expression in~\eqref{eq:radonslice} for $a>0$ and the definition of $U_\xi$ for $a<0$, we get, for real $a\ne0$,
\[
\int_\R e^{aw}U_\xi(w)\dd w=a^{-2}\sum_j\bigl(e^{a\lambda_j(\xi)}-e^{a\lambda^0_j(\xi)}\bigr)=a^{-2}\cD(a,a\xi),
\]
as in Lemma~\ref{lem:slices}(2). Since $\rho\ge0$, Tonelli's theorem and the substitution $s=w+\xi\tau$ at fixed $\tau$ give
\[
\int_\R e^{aw}U_\xi(w)\dd w=\int_0^1\!\!\int_\R e^{aw}\rho(w+\xi\tau,\tau)\dd w\dd\tau=\int_0^1\!\!\int_\R e^{as-a\xi\tau}\rho(s,\tau)\dd s\dd\tau .
\]
Every real $t$ is of the form $a\xi$, so $\cG(a,t)=\cD(a,t)/a^2$ equals $I(a,t):=\iint e^{as-t\tau}\rho(s,\tau)\dd s\dd\tau$ for
all real $a\ne0$ and all real $t$. Both functions are entire on $\C^2$: $\cG$ by Step~0 of Section~\ref{subsec:proofD}, and $I$
because $\rho$ is integrable with compact support. For fixed real $t$, $a\mapsto\cG(a,t)-I(a,t)$ is entire and vanishes on
$\R\setminus\{0\}$, hence identically; then for fixed $a\in\C$, $t\mapsto\cG(a,t)-I(a,t)$ is entire and vanishes on $\R$,
hence identically. This is~\eqref{eq:2bmv}. For the last statement of Theorem~\ref{thm:D}, put
$m(\tau)=\int_\R\rho(s,\tau)\dd s$, which is continuous on $(0,1)$ and integrable. Since $\cG=I$ on $\C^2$, Step~4 of
Section~\ref{subsec:proofD} gives $\int_0^1e^{-t\tau}m(\tau)\dd\tau=I(0,t)=\cG(0,t)=\|BgP\|_F^2\int_0^1e^{-t\tau}\dd\tau$, that is,
$\int_0^1e^{-t\tau}(m(\tau)-\|BgP\|_F^2)\dd\tau=0$ for every real $t$. Taking $t=0,1,2,\dots$ and
using that the polynomials in $e^{-\tau}$ are dense in $C[0,1]$ (Weierstrass), $\int_0^1\varphi\,(m-\|BgP\|_F^2)=0$ for every
continuous $\varphi$; so $m=\|BgP\|_F^2$ almost everywhere, hence everywhere on $(0,1)$ by continuity. \qed

\subsection{Theorem~\ref{thm:C} without distributions}
The proof of Theorem~\ref{thm:C} in Section~\ref{subsec:proofC} uses Theorem~\ref{thm:D} at $t=\mp ia$, Lemmas~\ref{lem:poisson}
and~\ref{lem:h} (residues, geometric series, and the fact that two harmonic functions that agree on an open subset of the
strip agree everywhere), and Tonelli's theorem. Only step~(b) uses distributions and the uniqueness of Fourier
transforms; we replace it by the following lemma.

\begin{lemma}\label{lem:U}
Let $Z\colon\R\to\C$ be continuous with $|Z(w)|\le Ce^{-c|w|}$ for some $c>0$. If $\int_\R e^{aw}Z(w)\dd w=0$ for all $a$ in a
nonempty open interval $J\subseteq(-c,c)$, then $Z=0$.
\end{lemma}

\begin{proof}
The function $A(a)=\int_\R e^{aw}Z(w)\dd w$ is holomorphic on $\{|\Re a|<c\}$ (differentiation under the integral sign) and
vanishes on $J$, hence on the whole strip. So $\hat Z(k)=\int e^{-ikw}Z(w)\dd w=A(-ik)=0$ for all real $k$. For $\epsilon>0$ and
$w_0\in\R$, Fubini's theorem and $\int_\R e^{iku-\epsilon^2k^2/2}\dd k=(\sqrt{2\pi}/\epsilon)e^{-u^2/(2\epsilon^2)}$ give
\[
0=\frac1{2\pi}\int_\R\hat Z(k)e^{ikw_0-\epsilon^2k^2/2}\dd k=\int_\R Z(w)\,\frac{e^{-(w-w_0)^2/(2\epsilon^2)}}{\epsilon\sqrt{2\pi}}\dd w ,
\]
and the right side tends to $Z(w_0)$ as $\epsilon\to0$, since $Z$ is bounded and continuous.
\end{proof}

\emph{Step (b) without distributions.} Let $\omega$, $\sigma$, $V$ and $W$ be as in the proof of Theorem~\ref{thm:C}. By
\eqref{eq:laplaceV}, $\hat\omega(a)-\hat\sigma(a)=O(a^2)$ as $a\to0$, so $\omega-\sigma$ has zero mass and zero first moment,
and $W(w)=\int(w-u)_+\dd(\omega-\sigma)(u)=\int(u-w)_+\dd(\omega-\sigma)(u)$. For $0<a<\pi$, the second expression and Fubini's
theorem, which applies because $\iint e^{aw}(u-w)_+\dd w\dd(\omega+\sigma)(u)=a^{-2}\int e^{au}\dd(\omega+\sigma)(u)<\infty$, give
$\int e^{aw}W(w)\dd w=a^{-2}(\hat\omega(a)-\hat\sigma(a))=\int e^{aw}V(w)\dd w$. The functions $W$ and $V$ are continuous and
$O(e^{-\pi|w|})$ (step (b) of Section~\ref{subsec:proofC}), so Lemma~\ref{lem:U} with $Z=W-V$, $c=\pi$ and $J=(0,\pi)$ gives
$W=V$. Steps~(c) and~(d) are unchanged; step~(d) uses only Theorem~\ref{thm:D} and the formula
$\cG(0,t)=\|BgP\|_F^2(1-e^{-t})/t$ of Step~4 of Section~\ref{subsec:proofD}.

\subsection{Remarks}

\begin{remark}[arbitrary Hermitian $B$; Stahl's theorem]\label{rem:RIgeneralB}
Let $B$ be a Hermitian matrix with distinct eigenvalues $b_1<\dots<b_n$ ($n\ge2$) and spectral projections $\Pi_k$, let
$H_d=\sum_k\Pi_kH\Pi_k$, and let $G_B=(b_1,b_n)\setminus\{b_2,\dots,b_{n-1}\}$. For $\tau\in G_B$ let $y_i(\tau)$ be the
eigenvalues of $(B-\tau)^{-1}H$. The proof of Theorem~\ref{thm:RI} extends to this setting: in a gap $(b_l,b_{l+1})$
exactly $\sum_{k>l}\operatorname{rank}\Pi_k$ roots lie in $\C_+$, the bounds of Lemma~\ref{lem:RIb} hold with constants that
depend on the gap and are integrable over it, and at an interior eigenvalue $b_k$ the boundary terms coming from the two
adjacent gaps cancel. The result is
\[
\frac1{2\pi}\int_{G_B}\sum_i|\Im y_i(\tau)|\dd\tau=\Tr H_+-\Tr(H_d)_+ .
\]
The argument of Section~\ref{subsec:RItoD} then shows that, for Hermitian $A$ and $B$,
\[
\Tr e^{aA-tB}-\sum_ke^{-tb_k}\Tr_{\Pi_k}e^{a\Pi_kA\Pi_k}=a^2\iint e^{as-t\tau}\rho_{A,B}(s,\tau)\dd s\dd\tau,
\]
where $\rho_{A,B}(s,\tau)=\frac1{2\pi}\sum_i|\Im x_i(s,\tau)|$ for $\tau\in G_B$, with $x_i(s,\tau)$ the roots of
$\det(A-s-x(B-\tau))$, and $\rho_{A,B}=0$ otherwise. At $a=1$ this says that $t\mapsto\Tr e^{A-tB}$ is the Laplace transform of
the positive measure $\sum_k\Tr_{\Pi_k}(e^{\Pi_kA\Pi_k})\delta_{b_k}+(\int_\R e^s\rho_{A,B}(s,\tau)\dd s)\dd\tau$ on $[b_1,b_n]$,
which is Stahl's theorem~\cite{Stahl}, the former conjecture of Bessis, Moussa and Villani~\cite{BMV}, here with a proof that
uses neither distribution theory nor Fourier analysis. We make no priority claim about this by-product; for the history of
the theorem and of the Fourier-slice mechanism see~\cite{Stahl,Eremenko,LS,He}. The details are in the accompanying
notes~\cite{repo-oqp27}, and they were checked in the independent verification of Section~\ref{subsec:independent}; they are
not used in this paper.
\end{remark}

\begin{remark}[formalisation]\label{rem:RIlean}
The Lean theorem \nolinkurl{OQP27.StripL3b.hyp_RI} is Theorem~\ref{thm:RI}, and \nolinkurl{OQP27.StripL3b.theorem1_bmv2}
is Theorem~\ref{thm:D}, proved from it as in Section~\ref{subsec:RItoD}. In the uniqueness step of the proof of
Theorem~\ref{thm:C}, the formal proof uses the Fourier inversion theorem of Mathlib in place of Lemma~\ref{lem:U}.
\end{remark}

%======================================================================
\raggedbottom
\begin{thebibliography}{99}

\bibitem{ADGL}
A.~Ac\'{\i}n, T.~Durt, N.~Gisin and J.~I. Latorre, \emph{Quantum nonlocality in two three-level systems}, Phys. Rev. A
\textbf{65} (2002), 052325.

\bibitem{AGG}
A.~Ac\'{\i}n, R.~Gill and N.~Gisin, \emph{Optimal Bell tests do not require maximally entangled states}, Phys. Rev.
Lett. \textbf{95} (2005), 210402.

\bibitem{BRL}
K.~Baek, J.~Ryu and J.~Lee, \emph{Robustness measures for quantifying nonlocality}, New J. Phys. \textbf{27} (2025),
053001; arXiv:2311.07077.

\bibitem{BGP}
J.-D. Bancal, N.~Gisin and S.~Pironio, \emph{Looking for symmetric Bell inequalities}, J. Phys. A \textbf{43} (2010),
385303.

\bibitem{BMV}
D.~Bessis, P.~Moussa and M.~Villani, \emph{Monotonic converging variational approximations to the functional integrals
in quantum statistical mechanics}, J. Math. Phys. \textbf{16} (1975), 2318--2325.

\bibitem{Boole}
G.~Boole, \emph{On the comparison of transcendents, with certain applications to the theory of definite integrals},
Philos. Trans. R. Soc. Lond. \textbf{147} (1857), 745--803.

\bibitem{Chen01}
J.-L. Chen, D.~Kaszlikowski, L.~C. Kwek, C.~H. Oh and M.~\.Zukowski, \emph{Entangled three-state systems violate local
realism more strongly than qubits: an analytical proof}, Phys. Rev. A \textbf{64} (2001), 052109.

\bibitem{Tsirelson}
B.~S. Cirel'son, \emph{Quantum generalizations of Bell's inequality}, Lett. Math. Phys. \textbf{4} (1980), 93--100.

\bibitem{CollinsGisin}
D.~Collins and N.~Gisin, \emph{A relevant two qubit Bell inequality inequivalent to the CHSH inequality}, J. Phys. A
\textbf{37} (2004), 1775.

\bibitem{CGLMP}
D.~Collins, N.~Gisin, N.~Linden, S.~Massar and S.~Popescu, \emph{Bell inequalities for arbitrarily high-dimensional
systems}, Phys. Rev. Lett. \textbf{88} (2002), 040404.

\bibitem{CLM}
L.~Colzani, E.~Laeng and L.~Monz\'on, \emph{Variations on a theme of Boole and Stein--Weiss}, J. Math. Anal. Appl.
\textbf{363} (2010), 225--229.

\bibitem{DKZ}
T.~Durt, D.~Kaszlikowski and M.~\.Zukowski, \emph{Violations of local realism with quantum systems described by
$N$-dimensional Hilbert spaces up to $N=16$}, Phys. Rev. A \textbf{64} (2001), 024101.

\bibitem{Eremenko}
A.~Eremenko, \emph{Herbert Stahl's proof of the BMV conjecture}, Sb. Math. \textbf{206} (2015), 87--92; arXiv:1312.6003.

\bibitem{Fine}
A.~Fine, \emph{Hidden variables, joint probability, and the Bell inequalities}, Phys. Rev. Lett. \textbf{48} (1982),
291--295.

\bibitem{Gill}
R.~D. Gill, \emph{Better Bell inequalities (passion at a distance)}, in: Asymptotics: Particles, Processes and Inverse
Problems, IMS Lecture Notes--Monograph Series \textbf{55} (2007), 135--148; arXiv:math/0610115.

\bibitem{GLZ}
J.~Gruca, W.~Laskowski and M.~\.Zukowski, \emph{Nonclassicality of pure two-qutrit entangled states}, Phys. Rev. A
\textbf{85} (2012), 022118.

\bibitem{He}
O.~Hein\"avaara, \emph{Tracial joint spectral measures}, arXiv:2310.03227 (2023).

\bibitem{Hormander}
L.~H\"ormander, \emph{The Analysis of Linear Partial Differential Operators I}, 2nd ed., Springer, Berlin, 1990.

\bibitem{HiGHS}
Q.~Huangfu and J.~A.~J. Hall, \emph{Parallelizing the dual revised simplex method}, Math. Program. Comput. \textbf{10}
(2018), 119--142.

\bibitem{IR}
M.~Ioannou and D.~Rosset, \emph{Noncommutative polynomial optimization under symmetry}, arXiv:2112.10803 (2021).

\bibitem{mpmath}
F.~Johansson et al., \emph{mpmath: a Python library for arbitrary-precision floating-point arithmetic},
\url{https://mpmath.org}.

\bibitem{KGZMZ}
D.~Kaszlikowski, P.~Gnaci\'nski, M.~\.Zukowski, W.~Miklaszewski and A.~Zeilinger, \emph{Violations of local realism by two
entangled $N$-dimensional systems are stronger than for two qubits}, Phys. Rev. Lett. \textbf{85} (2000), 4418--4421.

\bibitem{KKCZO}
D.~Kaszlikowski, L.~C. Kwek, J.-L. Chen, M.~\.Zukowski and C.~H. Oh, \emph{Clauser--Horne inequality for three-state
systems}, Phys. Rev. A \textbf{65} (2002), 032118.

\bibitem{Kulpa}
W.~Kulpa, \emph{The Poincar\'e--Miranda theorem}, Amer. Math. Monthly \textbf{104} (1997), 545--550.

\bibitem{Laeng}
E.~Laeng, \emph{On the $L^p$ norms of the Hilbert transform of a characteristic function}, J. Funct. Anal. \textbf{262}
(2012), 4534--4539.

\bibitem{LVN}
B.~Lang, T.~V\'ertesi and M.~Navascu\'es, \emph{Closed sets of correlations: answers from the zoo}, J. Phys. A \textbf{47}
(2014), 424029.

\bibitem{Lewin}
L.~Lewin, \emph{Polylogarithms and Associated Functions}, North-Holland, New York, 1981.

\bibitem{LS}
E.~H. Lieb and R.~Seiringer, \emph{Equivalent forms of the Bessis--Moussa--Villani conjecture}, J. Stat. Phys.
\textbf{115} (2004), 185--190.

\bibitem{Masanes}
Ll.~Masanes, \emph{Tight Bell inequality for $d$-outcome measurements correlations}, Quantum Inf. Comput. \textbf{3}
(2003), 345--358.

\bibitem{MS-classical}
A.~Mishra and A.~Senthilkumar, \emph{A rearrangement inequality for the discrete Hilbert transform, with an application to
optimal CGLMP measurements}, preprint (2026), \url{https://github.com/anshM123/CGLMP}.

\bibitem{repo-cglmp}
A.~Mishra and A.~Senthilkumar, \emph{CGLMP on maximally entangled states: exact optimality, rigidity and Tsirelson
bounds}, preprint and data repository (2026), \url{https://github.com/anshM123/CGLMP}.

\bibitem{repo-oqp27}
A.~Mishra and A.~Senthilkumar, accompanying repository of this paper: certificates, verification programs, logs, Lean~4
formalisation, independent verification reports and notes (2026),
\url{https://github.com/anshM123/anshM123-IQOQI-Vienna-Open-Problem-27}.

\bibitem{Lean}
L.~de~Moura and S.~Ullrich, \emph{The Lean~4 theorem prover and programming language}, in: CADE~28, Lecture Notes in
Comput. Sci. \textbf{12699}, Springer, Cham, 2021, 625--635.

\bibitem{Mathlib}
The mathlib Community, \emph{The Lean mathematical library}, in: Proceedings of the 9th ACM SIGPLAN International
Conference on Certified Programs and Proofs (CPP 2020), ACM, New York, 2020, 367--381.

\bibitem{OQP}
IQOQI Vienna, \emph{Open Quantum Problems}, Problem~27: \emph{The power of CGLMP inequalities}; archived copy of 29 October
2023:
\url{https://web.archive.org/web/20231029013544/https://oqp.iqoqi.oeaw.ac.at/the-power-of-cglmp-inequalities}.

\bibitem{PR}
S.~Popescu and D.~Rohrlich, \emph{Quantum nonlocality as an axiom}, Found. Phys. \textbf{24} (1994), 379--385.

\bibitem{SATWAP}
A.~Salavrakos, R.~Augusiak, J.~Tura, P.~Wittek, A.~Ac\'{\i}n and S.~Pironio, \emph{Bell inequalities tailored to maximally
entangled states}, Phys. Rev. Lett. \textbf{119} (2017), 040402.

\bibitem{SSKA}
S.~Sarkar, D.~Saha, J.~Kaniewski and R.~Augusiak, \emph{Self-testing quantum systems of arbitrary local dimension with
minimal number of measurements}, npj Quantum Inf. \textbf{7} (2021), 151.

\bibitem{Stahl}
H.~R. Stahl, \emph{Proof of the BMV conjecture}, Acta Math. \textbf{211} (2013), 255--290.

\bibitem{SW}
E.~M. Stein and G.~Weiss, \emph{An extension of a theorem of Marcinkiewicz and some of its applications}, J. Math. Mech.
\textbf{8} (1959), 263--284.

\bibitem{vDGG}
W.~van Dam, R.~D. Gill and P.~D. Gr\"unwald, \emph{The statistical strength of nonlocality proofs}, IEEE Trans. Inform.
Theory \textbf{51} (2005), 2812--2835; arXiv:quant-ph/0307125.

\bibitem{Zagier}
D.~Zagier, \emph{The dilogarithm function}, in: Frontiers in Number Theory, Physics, and Geometry~II, Springer, Berlin,
2007, 3--65.

\bibitem{Zhang}
Y.~Zhang, \emph{CGLMP measurements need not maximize statistical strength}, Zenodo (2026),
\url{https://doi.org/10.5281/zenodo.23022433}.

\bibitem{Zygmund}
A.~Zygmund, \emph{Trigonometric Series}, 3rd ed., Cambridge University Press, Cambridge, 2002.

\end{thebibliography}

\end{document}
